% Homeostasis — Biology Upper Sixth — Cours signé OnBuch+
% Official MINESEC syllabus. Compiler avec : tectonic BIO04-homeostasis.tex
\newcommand{\DOCMATIERE}{Biology}
\newcommand{\DOCNIVEAU}{Upper Sixth}
\newcommand{\DOCMODULE}{Upper Sixth — Biology}
\newcommand{\DOCLECON}{4}
\newcommand{\DOCTITRE}{Homeostasis}
% OnBuch+ — préambule « cours signés » ENRICHI (compilé avec tectonic / XeLaTeX)
% Système visuel « Pop » d'OnBuch+ : fond crème (jamais blanc), bordures encre
% épaisses, ombres « dures » décalées, titres Archivo Black en capitales.
% Version MATHÉMATIQUES : graphes (pgfplots/tikz), boîtes pédagogiques
% (définition, propriété, méthode, attention, le savais-tu, exercice résolu,
% exercice, TP, bilan), page de garde et signature OnBuch+ ancrée en bas.
%
% Macros attendues dans main.tex AVANT \input{preamble} :
%   \DOCMATIERE  \DOCNIVEAU  \DOCMODULE  \DOCLECON (numéro)  \DOCTITRE
\documentclass[11pt]{article}

\usepackage{fontspec}
\usepackage[english]{babel}
\usepackage[a4paper,margin=2cm,top=3.4cm,bottom=2.8cm,headheight=1.4cm,footskip=1.3cm]{geometry}
\usepackage{amsmath,amssymb,amsfonts}
\usepackage{mathtools} % \xmapsto, \coloneqq... souvent utilisés par les modèles
\usepackage{cancel} % simplifications barrées (bilans de forces)
% \abs{z} (module, valeur absolue) et \norm{u} (norme), utilisés par les modèles
\providecommand{\abs}{}\let\abs\relax\DeclarePairedDelimiter{\abs}{\lvert}{\rvert}
\providecommand{\norm}{}\let\norm\relax\DeclarePairedDelimiter{\norm}{\lVert}{\rVert}
\usepackage[table]{xcolor} % [table] : fournit aussi \rowcolors (alternance de lignes)
\usepackage{pagecolor}
\usepackage{tikz}
\usetikzlibrary{arrows.meta,positioning,calc,shapes.geometric,shapes.misc,shapes.arrows,decorations.pathmorphing,decorations.pathreplacing,decorations.markings,patterns,shadows,mindmap,trees,fit,backgrounds,matrix,intersections,angles,quotes}
\usepackage{pgfplots}
\usepackage[version=4]{mhchem} % équations nucléaires \ce{^{235}_{92}U}
\usepackage[european,siunitx]{circuitikz} % schémas électriques
\pgfplotsset{compat=1.18}
\usepgfplotslibrary{fillbetween,groupplots} % aires entre courbes (calcul intégral)
\usepackage{enumitem}
\usepackage{fancyhdr}
% [most] charge toutes les bibliothèques tcolorbox (listings, minted, raster,
% documentation...) et gonfle inutilement la mémoire TeX sur un document long
% (~300 boîtes) : on ne charge que ce qui est réellement utilisé.
\usepackage[skins,breakable]{tcolorbox}
\usepackage{graphicx}
\usepackage{colortbl}
\usepackage{booktabs}
\usepackage{tabularx}
\usepackage{diagbox} % case d'en-tête diagonale des tableaux à double entrée
% Colonnes extensibles centrée / gauche / droite (C, L, R), souvent utilisées par les modèles
\newcolumntype{C}{>{\centering\arraybackslash}X}
\newcolumntype{L}{>{\raggedright\arraybackslash}X}
\newcolumntype{R}{>{\raggedleft\arraybackslash}X}
\usepackage{array}
\usepackage{multicol}
\usepackage{siunitx}
% parse-numbers=false : accepte aussi \SI{2,5 \times 10^{-3}}{...}
\sisetup{locale=UK,output-decimal-marker={.},group-minimum-digits=4,parse-numbers=false}
\DeclareSIUnit\ohm{\ensuremath{\symup{\Omega}}} % U+2126 absent de la police
\DeclareSIUnit\division{div} % réglages d'oscilloscope (ms/div, V/div)
\DeclareSIUnit\tr{tr} % tours (vitesse de rotation, ex. \SI{1500}{\tr\per\minute})
\usepackage{qrcode}
% \degree : commande fréquemment utilisée par les modèles pour un angle ou une
% température en dehors de \SI{}{} (non fournie par les paquets chargés).
\providecommand{\degree}{^{\circ}}
% \qed (fin de démonstration), utilisé par les modèles sans amsthm
\providecommand{\qed}{\hfill$\square$}
% \barre : double barre des valeurs interdites dans un tableau de variations (tkz-tab)
\providecommand{\barre}{\ensuremath{\|}}
% environnement proof (amsthm non chargé) utilisé par les modèles
\ifdefined\proof\else\newenvironment{proof}[1][Proof]{\par\noindent\textit{#1.}\ }{\qed\par}\fi
\usepackage{lastpage}
\usepackage{hyperref}
\hypersetup{hidelinks}

% ============================================================
% Palette « Pop » OnBuch+ — mêmes hex que la classe P (plus_ui.dart)
% ============================================================
\definecolor{poporange}{HTML}{F59321}
\definecolor{popdark}{HTML}{E07A0C}
\definecolor{popcream}{HTML}{FFF6E8}
\definecolor{popink}{HTML}{1C1714}
\definecolor{poppurple}{HTML}{7A5AE0}
\definecolor{popgreen}{HTML}{2E9E6A}
\definecolor{poppink}{HTML}{E0457B}
\definecolor{popblue}{HTML}{2F7DE1}
\definecolor{popgold}{HTML}{C98A12}
\definecolor{popmuted}{HTML}{8A7F76}
\definecolor{popline}{HTML}{EADFD2}
\definecolor{popgray}{HTML}{8A7F76}
\colorlet{popgrey}{popgray}
% Alias tolérants (le modèle nomme parfois les couleurs différemment)
\colorlet{popwhite}{white}
\colorlet{popblack}{popink}
\colorlet{popred}{poppink}
\colorlet{popyellow}{popgold}
% Teintes claires pour fonds de tableaux / zones
\colorlet{popblueL}{popblue!10!white}
\colorlet{poporangeL}{poporange!14!white}
\colorlet{popgreenL}{popgreen!12!white}
\colorlet{poppurpleL}{poppurple!12!white}
\colorlet{poppinkL}{poppink!12!white}
\colorlet{popgoldL}{popgold!14!white}

\pagecolor{popcream}

% ============================================================
% Typographie
% ============================================================
\defaultfontfeatures{Ligatures=TeX}
\setmainfont{PlusJakartaSans-Medium.ttf}[Path=../../../fonts/,BoldFont=PlusJakartaSans-Bold.ttf]
% Maths en sans-serif (Fira Math) avec chiffres et lettres droites de Plus
% Jakarta, pour que les formules se fondent dans le texte.
\usepackage[math-style=ISO,bold-style=ISO]{unicode-math}
\setmathfont{FiraMath-Regular.otf}[Path=../../../fonts/]
\setmathfont{PlusJakartaSans-Medium.ttf}[Path=../../../fonts/,range={up/num,up/{Latin,latin},\mathup}]
% (icomma retiré : décimales avec point en anglais)
% Caractères mathématiques Unicode absents des polices de texte (Plus Jakarta,
% Archivo Black) : rendus par leur équivalent mathématique, en texte comme en
% formule — sinon ils s'affichent en carré vide (ex. « Arithmétique dans ℤ »).
\usepackage{newunicodechar}
\newunicodechar{ℝ}{\ensuremath{\mathbb{R}}}
\newunicodechar{ℤ}{\ensuremath{\mathbb{Z}}}
\newunicodechar{ℂ}{\ensuremath{\mathbb{C}}}
\newunicodechar{ℕ}{\ensuremath{\mathbb{N}}}
\newunicodechar{ℚ}{\ensuremath{\mathbb{Q}}}
\newunicodechar{𝔻}{\ensuremath{\mathbb{D}}}
\newunicodechar{α}{\ensuremath{\alpha}}
\newunicodechar{β}{\ensuremath{\beta}}
\newunicodechar{γ}{\ensuremath{\gamma}}
\newunicodechar{δ}{\ensuremath{\delta}}
\newunicodechar{ε}{\ensuremath{\varepsilon}}
\newunicodechar{θ}{\ensuremath{\theta}}
\newunicodechar{λ}{\ensuremath{\lambda}}
\newunicodechar{μ}{\ensuremath{\mu}}
\newunicodechar{σ}{\ensuremath{\sigma}}
\newunicodechar{τ}{\ensuremath{\tau}}
\newunicodechar{φ}{\ensuremath{\varphi}}
\newunicodechar{ψ}{\ensuremath{\psi}}
\newunicodechar{ω}{\ensuremath{\omega}}
\newunicodechar{Σ}{\ensuremath{\Sigma}}
% majuscules produites par \MakeUppercase dans les titres (α-aminés -> Α)
\newunicodechar{Α}{\ensuremath{\alpha}}
\newunicodechar{Β}{\ensuremath{\beta}}
\newunicodechar{Γ}{\ensuremath{\Gamma}}
\newunicodechar{Δ}{\ensuremath{\Delta}}
\newunicodechar{Ε}{\ensuremath{\varepsilon}}
\newunicodechar{Θ}{\ensuremath{\Theta}}
\newunicodechar{Λ}{\ensuremath{\Lambda}}
\newunicodechar{Μ}{\ensuremath{\mu}}
\newunicodechar{Τ}{\ensuremath{\tau}}
\newunicodechar{Φ}{\ensuremath{\Phi}}
\newunicodechar{Ψ}{\ensuremath{\Psi}}
\newunicodechar{Ω}{\ensuremath{\Omega}}
\newunicodechar{↦}{\ensuremath{\mapsto}}
\newunicodechar{∈}{\ensuremath{\in}}
\newunicodechar{∉}{\ensuremath{\notin}}
\newunicodechar{∧}{\ensuremath{\wedge}}
\newunicodechar{⇔}{\ensuremath{\Leftrightarrow}}
\newunicodechar{⟺}{\ensuremath{\Longleftrightarrow}}
\newunicodechar{⇒}{\ensuremath{\Rightarrow}}
\newunicodechar{⟹}{\ensuremath{\Longrightarrow}}
\newunicodechar{∘}{\ensuremath{\circ}}
\newunicodechar{∪}{\ensuremath{\cup}}
\newunicodechar{∩}{\ensuremath{\cap}}
\newunicodechar{≡}{\ensuremath{\equiv}}
\newunicodechar{⊂}{\ensuremath{\subset}}
\newunicodechar{⊥}{\ensuremath{\perp}}
\newunicodechar{∃}{\ensuremath{\exists}}
\newunicodechar{∀}{\ensuremath{\forall}}
\newunicodechar{∛}{\ensuremath{\sqrt[3]{\,}}}
\newunicodechar{∜}{\ensuremath{\sqrt[4]{\,}}}
\newunicodechar{⃗}{\ensuremath{{}^{\to}}}
\newunicodechar{ⁿ}{\ifmmode^{n}\else\textsuperscript{\ensuremath{n}}\fi}
\newunicodechar{ˣ}{\ifmmode^{x}\else\textsuperscript{\ensuremath{x}}\fi}
\newunicodechar{ᵇ}{\ifmmode^{b}\else\textsuperscript{\ensuremath{b}}\fi}
\newunicodechar{ʳ}{\ifmmode^{r}\else\textsuperscript{\ensuremath{r}}\fi}
\newunicodechar{ᵘ}{\ifmmode^{u}\else\textsuperscript{\ensuremath{u}}\fi}
\newunicodechar{ʸ}{\ifmmode^{y}\else\textsuperscript{\ensuremath{y}}\fi}
\newunicodechar{ᵃ}{\ifmmode^{a}\else\textsuperscript{\ensuremath{a}}\fi}
\newunicodechar{ᵉ}{\ifmmode^{e}\else\textsuperscript{\ensuremath{e}}\fi}
\newunicodechar{ᵏ}{\ifmmode^{k}\else\textsuperscript{\ensuremath{k}}\fi}
\newunicodechar{ᵖ}{\ifmmode^{p}\else\textsuperscript{\ensuremath{p}}\fi}
\newunicodechar{ᴺ}{\ifmmode^{N}\else\textsuperscript{\ensuremath{N}}\fi}
\newunicodechar{ᶜ}{\ifmmode^{c}\else\textsuperscript{\ensuremath{c}}\fi}
\newunicodechar{ʰ}{\ifmmode^{h}\else\textsuperscript{\ensuremath{h}}\fi}
\newunicodechar{ᵠ}{\ifmmode^{q}\else\textsuperscript{\ensuremath{q}}\fi}
\newunicodechar{ₙ}{\ifmmode_{n}\else\textsubscript{\ensuremath{n}}\fi}
\newunicodechar{ₐ}{\ifmmode_{a}\else\textsubscript{\ensuremath{a}}\fi}
\newunicodechar{ᵢ}{\ifmmode_{i}\else\textsubscript{\ensuremath{i}}\fi}
\newunicodechar{ᵣ}{\ifmmode_{r}\else\textsubscript{\ensuremath{r}}\fi}
\newunicodechar{ₖ}{\ifmmode_{k}\else\textsubscript{\ensuremath{k}}\fi}
\newunicodechar{ᵦ}{\ifmmode_{\beta}\else\textsubscript{\ensuremath{\beta}}\fi}
\newunicodechar{ₓ}{\ifmmode_{x}\else\textsubscript{\ensuremath{x}}\fi}
\newfontfamily\popdisplay{ArchivoBlack-Regular.ttf}[Path=../../../fonts/]
\newfontfamily\poplogo{SpaceGrotesk-Bold.ttf}[Path=../../../fonts/]
\newfontfamily\popsemi{PlusJakartaSans-SemiBold.ttf}[Path=../../../fonts/]
\newfontfamily\popxbold{PlusJakartaSans-ExtraBold.ttf}[Path=../../../fonts/]
\newfontfamily\popxboldls{PlusJakartaSans-ExtraBold.ttf}[Path=../../../fonts/,LetterSpace=3.0]

\setlist[enumerate]{leftmargin=1.7em,itemsep=5pt,topsep=4pt}
\setlist[itemize]{leftmargin=1.7em,itemsep=4pt,topsep=4pt,label={\color{poporange}\textbullet}}
\setlength{\parindent}{0pt}
\setlength{\parskip}{5pt}
\renewcommand{\arraystretch}{1.25}

% pgfplots : style Pop par défaut
\pgfplotsset{
  every axis/.append style={axis line style={popink,line width=1pt},tick style={popink},
    grid style={popline,line width=0.5pt},label style={font=\small},tick label style={font=\footnotesize}},
  popcurve/.style={poporange,line width=1.8pt,no markers,smooth},
}
% Même style disponible dans un \draw TikZ simple (hors pgfplots)
\tikzset{popcurve/.style={poporange,line width=1.8pt}}
\tikzset{popgrid/.style={popline,very thin}}
\colorlet{popcurve}{poporange} % le nom du style est parfois utilisé comme couleur

% ============================================================
% Pill / squiggle
% ============================================================
\newcommand{\poppill}[3]{%
  \tikz[baseline=-0.55ex]\node[draw=popink,line width=1.2pt,rounded corners=2.6pt,%
    fill=#2,inner xsep=6.5pt,inner ysep=3pt]{%
    \popxboldls\fontsize{7}{7}\selectfont\color{#3}\MakeUppercase{#1}};%
}
\newcommand{\popsquiggle}[1]{%
  \tikz[baseline=-0.4ex]{
    \draw[#1,line width=2.6pt,line cap=round]
      plot [smooth] coordinates {(0,0.09) (0.34,0.01) (0.68,0.21) (1.02,0.01) (1.36,0.10)};
  }%
}
% Mot-clé surligné (marqueur orange sous le texte)
\newcommand{\cle}[1]{{\popxbold\color{popdark}#1}}

% ============================================================
% En-tête / pied
% ============================================================
\pagestyle{fancy}
\fancyhf{}
\renewcommand{\headrulewidth}{0pt}
\renewcommand{\footrulewidth}{0pt}
\lhead{\raisebox{-0.15ex}{\poplogo\fontsize{13}{13}\selectfont\color{popink}OnBuch\color{poporange}+}}
\rhead{\poppill{\DOCMATIERE\ \textbullet\ \DOCNIVEAU\ \textbullet\ Chapter \DOCLECON}{white}{popink}}
\lfoot{\popxbold\fontsize{7.6}{7.6}\selectfont\color{popmuted}\MakeUppercase{Signed course OnBuch+}\ \textbullet\ {\popsemi\DOCTITRE}}
\rfoot{\tikz[baseline=-0.6ex]\node[rounded corners=8.5pt,fill=popink,minimum height=17pt,minimum width=17pt,inner xsep=5pt,inner sep=0]{\popxbold\fontsize{8}{8}\selectfont\color{popcream}\thepage\,/\,\pageref*{LastPage}};}

% ============================================================
% Titres
% ============================================================
\newcommand{\chaptitre}[2]{%
  \begin{tcolorbox}[enhanced,colback=poporange,colframe=popink,boxrule=2.6pt,arc=11pt,
      drop shadow=popink,left=15pt,right=15pt,top=12pt,bottom=11pt,before skip=3pt,after skip=18pt]
    \poppill{#2}{popink}{popcream}\par\vspace{9pt}
    {\raggedright\hyphenpenalty=10000\exhyphenpenalty=10000\popdisplay\fontsize{22}{24}\selectfont\color{popcream}\MakeUppercase{#1}\par}
  \end{tcolorbox}
}
% Section numérotée : pastille orange avec le numéro + titre Archivo
\newcounter{popsec}
\newcommand{\coursec}[1]{%
  \stepcounter{popsec}\setcounter{popsub}{0}%
  \par\vspace{16pt}\noindent
  \begin{minipage}{\linewidth}
  \tikz[baseline=-0.7ex]\node[draw=popink,line width=1.6pt,fill=poporange,rounded corners=4pt,minimum size=22pt,inner sep=0]{\popdisplay\fontsize{12}{12}\selectfont\color{popcream}\thepopsec};%
  \hspace{8pt}{\popdisplay\fontsize{15}{17}\selectfont\color{popink}\MakeUppercase{#1}}\par
  \vspace{4pt}\noindent{\color{poporange}\rule{36pt}{3.6pt}}
  \end{minipage}\par\nopagebreak\vspace{8pt}
}
\newcounter{popsub}
\newcommand{\courssub}[1]{%
  \stepcounter{popsub}%
  \par\vspace{9pt}\noindent{\popxbold\fontsize{11.5}{13}\selectfont\color{popink}{\color{poporange}\thepopsec.\thepopsub}\hspace{6pt}#1}\par\nopagebreak\vspace{4pt}
}

% ============================================================
% Boîtes pédagogiques — même recette : fond blanc, bordure épaisse colorée,
% ombre dure encre, badge plein. Titre optionnel [#1].
% ============================================================
\tcbset{popbase/.style={breakable,enhanced,colback=white,boxrule=2.3pt,arc=9pt,
  shadow={3pt}{-3pt}{0pt}{popink},left=14pt,right=14pt,top=11pt,bottom=11pt,before skip=11pt,after skip=15pt,
  fontupper=\popsemi\fontsize{10.6}{14.5}\selectfont\color{popink},
  fontlower=\popsemi\fontsize{10.6}{14.5}\selectfont\color{popink}}}
% Coche dessinée (le glyphe ✓ n'existe pas dans les polices du document)
\newcommand{\popcheck}[1]{\tikz[baseline=-0.5ex]\draw[#1,line width=1.6pt,line cap=round,line join=round](-0.13,0.0)--(-0.04,-0.09)--(0.13,0.1);}
\providecommand{\coche}{\popcheck{popgreen}}
\providecommand{\resultat}[1]{\par\smallskip\noindent\fcolorbox{popgreen}{popgreenL}{\parbox{\dimexpr\linewidth-2\fboxsep-2\fboxrule\relax}{\textbf{Result.} #1}}\par\smallskip}
\newcommand{\popboxhead}[3]{\poppill{#1}{#2}{white}\ifx\relax#3\relax\else\hspace{6pt}{\popxbold\fontsize{10.6}{12}\selectfont\color{popink}#3}\fi\par\vspace{7pt}}

\newtcolorbox{aretenir}[1][]{popbase,colframe=popgreen,before upper={\popboxhead{Key points}{popgreen}{#1}}}
\newtcolorbox{exemplebox}[1][]{popbase,colframe=poporange,before upper={\popboxhead{Exemple}{poporange}{#1}}}
\newtcolorbox{definition}[1][]{popbase,colframe=popblue,before upper={\popboxhead{Definition}{popblue}{#1}}}
\newtcolorbox{propriete}[1][]{popbase,colframe=poppurple,before upper={\popboxhead{Property}{poppurple}{#1}}}
\newtcolorbox{methode}[1][]{popbase,colframe=popgold,before upper={\popboxhead{Method}{popgold}{#1}}}
\newtcolorbox{attention}[1][]{popbase,colframe=poppink,before upper={\popboxhead{Warning}{poppink}{#1}}}
\newtcolorbox{experience}[1][]{popbase,colframe=popblue,colback=popblueL,before upper={\popboxhead{Experiment}{popblue}{#1}}}
% « Le savais-tu ? » : carte violette pleine (contexte, culture, Cameroun)
\newtcolorbox{savaistu}[1][]{popbase,colback=poppurple,colframe=popink,
  fontupper=\popsemi\fontsize{10.4}{14.2}\selectfont\color{white},
  before upper={\poppill{Did you know?}{white}{poppurple}\ifx\relax#1\relax\else\hspace{6pt}{\popxbold\color{white}#1}\fi\par\vspace{7pt}}}
% Exercice résolu : énoncé en haut, \tcblower puis la solution sur fond crème
\newcounter{popexo}\newcounter{popexr}
\newtcolorbox{exoresolu}[1][]{popbase,colframe=poporange,colbacklower=popcream,
  segmentation style={popink,line width=1.2pt,dashed},
  before upper={\stepcounter{popexr}\popboxhead{Worked example \thepopexr}{poporange}{#1}},
  before lower={\poppill{Solution}{popink}{popcream}\par\vspace{6pt}}}
% Exercice (non résolu en place) — la correction est dans la partie Corrigés
\newtcolorbox{exercice}[1][]{popbase,colframe=popink,
  before upper={\stepcounter{popexo}\popboxhead{Exercise \thepopexo}{popink}{#1}}}
% Corrigé
\newtcolorbox{corrige}[1][]{popbase,colframe=popgreen,colback=popgreenL,
  before upper={\popboxhead{Solution}{popgreen}{#1}}}
% Objectifs / prérequis / bilan
\newtcolorbox{objectifs}{popbase,colframe=popink,colback=poporangeL,
  before upper={\popboxhead{Chapter objectives}{popink}{}}}
\newtcolorbox{prerequis}{popbase,colframe=popink,colback=popblueL,
  before upper={\popboxhead{Prerequisites}{popblue}{}}}
\newtcolorbox{bilan}{popbase,colframe=popink,colback=white,boxrule=2.8pt,
  before upper={\popboxhead{Fiche bilan}{poporange}{L'essentiel à connaître pour l'examen}}}
% Figure encadrée avec légende
\newtcolorbox{popfigure}[1][]{enhanced,breakable=false,colback=white,colframe=popline,boxrule=1.4pt,arc=8pt,
  left=10pt,right=10pt,top=10pt,bottom=8pt,before skip=10pt,after skip=12pt,halign=center,
  lower separated=false,halign lower=center,
  fontlower=\popsemi\fontsize{9}{11.5}\selectfont\color{popmuted},
  #1}
\newcommand{\legende}[1]{\par\vspace{6pt}{\popsemi\fontsize{9}{11.5}\selectfont\color{popmuted}#1}}

% ============================================================
% Signature OnBuch+
% ============================================================
\newcommand{\poplogoinline}[1]{{\poplogo\fontsize{#1}{#1}\selectfont\color{popink}OnBuch\color{poporange}+}}
\newcommand{\signatureonbuch}{%
  \begin{tcolorbox}[enhanced,colback=popink,colframe=popink,boxrule=2.2pt,arc=10pt,
      drop shadow=poporange,left=14pt,right=14pt,top=10pt,bottom=10pt,before skip=0pt,after skip=0pt]
    \tikz[baseline=-0.7ex]\node[circle,fill=popgreen,draw=popcream,line width=1.3pt,minimum size=22pt,inner sep=0]{\popcheck{white}};%
    \hspace{9pt}\begin{minipage}[c]{0.62\linewidth}
      {\popxbold\fontsize{12}{13}\selectfont\color{popcream}Signed\ \textbullet\ {\poplogo OnBuch\color{poporange}+}}\par\vspace{2pt}
      {\raggedright\popsemi\fontsize{8}{10.5}\selectfont\color{popline}\MakeUppercase{OnBuch+ teaching team}\\\MakeUppercase{Follows the official MINESEC syllabus}\par}
    \end{minipage}\hfill
    {\poplogo\fontsize{22}{22}\selectfont\color{popcream}OnBuch\color{poporange}+}
  \end{tcolorbox}%
}

% ============================================================
% Page de garde
% #1 titre · #2 matière · #3 niveau · #4 module · #5 n° leçon · #6 durée ·
% #7 description / objectifs (texte ou itemize)
% ============================================================
\newcommand{\pagegarde}[7]{%
  \thispagestyle{empty}
  \newgeometry{a4paper,margin=1.7cm,top=1.6cm,bottom=1.4cm}%
  \begin{tikzpicture}[remember picture,overlay]
    \fill[poporange!18] ([xshift=-3.2cm,yshift=-3.6cm]current page.north east) circle (4.6cm);
    \fill[poppurple!14] ([xshift=2.4cm,yshift=7.6cm]current page.south west) circle (3.8cm);
  \end{tikzpicture}%
  {\poplogo\fontsize{30}{30}\selectfont\color{popink}OnBuch\color{poporange}+}\hfill
  \raisebox{0.6ex}{\poppill{Signed course \textbullet\ Premium}{popink}{popcream}}\par
  \vspace{1pt}\popsquiggle{poppurple}\par
  \vspace{8pt}
  \begin{tcolorbox}[enhanced,colback=poporange,colframe=popink,boxrule=2.8pt,arc=13pt,
      drop shadow=popink,left=18pt,right=18pt,top=12pt,bottom=13pt,after skip=10pt]
    \poppill{#2 \textbullet\ #3}{popink}{popcream}\hspace{5pt}\poppill{#4}{white}{popink}\par\vspace{8pt}
    {\popdisplay\fontsize{14}{14}\selectfont\color{popink}CHAPTER #5}\par\vspace{4pt}
    {\raggedright\hyphenpenalty=10000\exhyphenpenalty=10000\popdisplay\fontsize{25}{27}\selectfont\color{popcream}\MakeUppercase{#1}\par}
    \vspace{8pt}
    {\popxbold\fontsize{9.5}{9.5}\selectfont\color{popink}\MakeUppercase{Suggested time: #6}}
  \end{tcolorbox}
  \begin{tcolorbox}[enhanced,colback=white,colframe=popink,boxrule=2.2pt,arc=10pt,
      drop shadow=popink,left=16pt,right=16pt,top=10pt,bottom=10pt,after skip=8pt]
    \poppill{In this chapter}{poppurple}{white}\par\vspace{5pt}
    \popsemi\fontsize{9.3}{12.6}\selectfont\color{popink}#7
  \end{tcolorbox}
  \noindent\begin{minipage}[t]{0.487\linewidth}
    \begin{tcolorbox}[enhanced,colback=popgreen,colframe=popink,boxrule=2.2pt,arc=10pt,
        drop shadow=popink,left=11pt,right=11pt,top=10pt,bottom=10pt,height=3.1cm,valign=center]
      \tcbox[colback=white,colframe=popink,boxrule=1.3pt,arc=5pt,boxsep=0pt,nobeforeafter,
        left=3pt,right=3pt,top=3pt,bottom=3pt,valign=center]{\qrcode[height=1.9cm]{https://play.google.com/store/apps/details?id=com.luvvix.onbuch&hl=en}}%
      \hspace{8pt}\begin{minipage}[c]{0.5\linewidth}
        {\popdisplay\fontsize{11}{12.5}\selectfont\color{white}\MakeUppercase{Download OnBuch}}\par\vspace{4pt}
        {\raggedright\popsemi\fontsize{8.4}{11}\selectfont\color{white}Free \textbullet\ Google Play\\AI tutor, past papers, results\par}
      \end{minipage}
    \end{tcolorbox}
  \end{minipage}\hfill
  \begin{minipage}[t]{0.487\linewidth}
    \begin{tcolorbox}[enhanced,colback=poppurple,colframe=popink,boxrule=2.2pt,arc=10pt,
        drop shadow=popink,left=11pt,right=11pt,top=10pt,bottom=10pt,height=3.1cm,valign=center]
      \tikz[baseline=-0.7ex]\node[circle,fill=white,draw=popink,line width=1.3pt,minimum size=1.6cm,inner sep=0]{\popdisplay\fontsize{24}{24}\selectfont\color{poppurple}+};%
      \hspace{8pt}\begin{minipage}[c]{0.6\linewidth}
        {\popdisplay\fontsize{11}{12.5}\selectfont\color{white}\MakeUppercase{Go OnBuch+}}\par\vspace{4pt}
        {\raggedright\popsemi\fontsize{8.4}{11}\selectfont\color{white}All signed courses, unlimited AI tutor, PDF sheets — OnBuch+ tab in the app\par}
      \end{minipage}
    \end{tcolorbox}
  \end{minipage}
  \vspace{8pt}
  \popcontenu
  \vfill
  \signatureonbuch
  \restoregeometry
  \newpage
}

% Rangée « Ce cours contient » (page de garde)
\newcommand{\poptile}[3]{% #1 couleur #2 gros texte #3 légende
  \begin{tcolorbox}[enhanced,colback=#1,colframe=popink,boxrule=2pt,arc=9pt,shadow={2.5pt}{-2.5pt}{0pt}{popink},
      left=6pt,right=6pt,top=9pt,bottom=9pt,halign=center,valign=center,height=1.75cm,width=\linewidth,nobeforeafter]
    {\popdisplay\fontsize{12.5}{13}\selectfont\color{white}\MakeUppercase{#2}}\par\vspace{3pt}
    {\centering\popsemi\fontsize{7.8}{9.5}\selectfont\color{white}#3\par}
  \end{tcolorbox}}
\newcommand{\popcontenu}{%
  \par\noindent{\popxboldls\fontsize{7.5}{7.5}\selectfont\color{popmuted}THIS SIGNED COURSE CONTAINS}\par\vspace{7pt}
  \noindent\begin{minipage}[t]{0.235\linewidth}\poptile{popblue}{Course}{complete, illustrated, step by step}\end{minipage}\hfill
  \begin{minipage}[t]{0.235\linewidth}\poptile{poporange}{Methods}{and worked examples}\end{minipage}\hfill
  \begin{minipage}[t]{0.235\linewidth}\poptile{poppink}{Exercises}{exam style + solutions}\end{minipage}\hfill
  \begin{minipage}[t]{0.235\linewidth}\poptile{popgreen}{Summary sheet}{the essentials to remember}\end{minipage}\par}

% Sommaire
\newtcolorbox{sommaire}{enhanced,colback=white,colframe=popink,boxrule=2.2pt,arc=10pt,
  drop shadow=popink,left=18pt,right=18pt,top=13pt,bottom=13pt,before skip=6pt,after skip=20pt,
  before upper={\poppill{Contents}{poppurple}{white}\par\vspace{9pt}\popsemi\fontsize{10.6}{14.5}\selectfont\color{popink}}}

% Signature de fin de cours — poussée tout en bas de la dernière page
\newcommand{\findecours}{%
  \par\vfill\nopagebreak
  \begin{center}\popsquiggle{poporange}\end{center}\vspace{4pt}
  \signatureonbuch
}
\AtBeginDocument{\def\llbracket{\mathopen{[\mkern-3.5mu[}}\def\rrbracket{\mathclose{]\mkern-3.5mu]}}} % intervalles d'entiers (glyphe ⟦ absent de la police)
% Glyphes absents de Fira Math : redéfinis à partir de glyphes disponibles
\AtBeginDocument{%
  \def\setminus{\mathbin{\backslash}}\let\smallsetminus\setminus
  \def\vdots{\mathord{\vbox{\baselineskip3.2pt\lineskiplimit0pt\kern4pt\hbox{.}\hbox{.}\hbox{.}}}}%
  \def\ddots{\mathinner{\mkern1mu\raise6.5pt\hbox{.}\mkern2mu\raise3.6pt\hbox{.}\mkern2mu\raise0.7pt\hbox{.}\mkern1mu}}%
  \def\triangle{\mathord{\scalebox{0.9}{$\symup{\Delta}$}}}%
  \def\nabla{\mathord{\raisebox{\depth}{\scalebox{1}[-1]{$\symup{\Delta}$}}}}%
  \def\checkmark{\popcheck{popdark}}%
  \def\star{\mathbin{\ast}}\def\bigstar{\mathord{\ast}}%
  \def\diamond{\mathbin{\scalebox{0.6}{$\Diamond$}}}%
  \def\bigcup{\mathop{\mathchoice{\raisebox{-0.3em}{\scalebox{1.7}{$\cup$}}}{\scalebox{1.25}{$\cup$}}{\cup}{\cup}}\displaylimits}%
  \def\bigcap{\mathop{\mathchoice{\raisebox{-0.3em}{\scalebox{1.7}{$\cap$}}}{\scalebox{1.25}{$\cap$}}{\cap}{\cap}}\displaylimits}%
  \def\lesseqqgtr{\mathrel{\lessgtr}}\def\bigoplus{\mathop{\oplus}\displaylimits}%
}
\providecommand{\ssi}{if and only if} % abréviation fréquente
\AtBeginDocument{\providecommand{\wideparen}{\overparen}} % arc de cercle

%% FIGFIX-BEGIN v3 — correctifs globaux des figures (docs/onbuch-generation/tools/figfix)
\usepackage{adjustbox}\usepackage{etoolbox}
% toute figure TikZ/pgfplots plus large que la ligne est réduite à la largeur disponible
\BeforeBeginEnvironment{tikzpicture}{\begin{adjustbox}{max width=\linewidth}}
\AfterEndEnvironment{tikzpicture}{\end{adjustbox}}
% légendes pgfplots lisibles par-dessus les courbes
\pgfplotsset{every axis/.append style={legend style={fill=white,fill opacity=0.92,text opacity=1,draw=black!15,inner sep=3pt}}}
% étiquettes d'arêtes (syntaxe edge["..."]) avec fond blanc
\tikzset{every edge quotes/.append style={fill=white,inner sep=1.2pt}}
% bulles de cartes mentales : texte à la ligne dans la bulle, bulle un peu plus grande
\tikzset{every concept/.append style={text width=2.0cm,align=center,minimum size=2.3cm,inner sep=2pt}}
%% FIGFIX-END
\begin{document}

\pagegarde{Homeostasis}{Biology}{Upper Sixth}{Upper Sixth — Biology}{4}{13 periods}{This chapter studies homeostasis: the maintenance of a stable internal environment in living organisms. It examines the roles of the liver, the skin and the kidneys in temperature, blood glucose, water and excretion control, the feedback mechanisms involved, and excretory products in both animals and plants.
\vspace{4pt}
\begin{multicols}{2}\small
\begin{itemize}
\item By the end of this chapter I can define homeostasis and explain its importance for cell survival.
\item By the end of this chapter I can describe the structure of the liver and explain its homeostatic roles (blood glucose regulation, deamination, detoxification, bile production, heat production).
\item By the end of this chapter I can explain negative and positive feedback mechanisms using biological examples.
\item By the end of this chapter I can describe the structure of mammalian skin and relate each structure to its functions, especially thermoregulation.
\item By the end of this chapter I can identify the main nitrogenous excretory products of animals (ammonia, urea, uric acid) and relate them to habitat and water availability.
\item By the end of this chapter I can describe the structure of the kidney and nephron and explain the stages of urine formation (ultrafiltration, selective reabsorption, secretion).
\end{itemize}\end{multicols}}

\begin{sommaire}
\begin{multicols}{2}
\begin{enumerate}
\item Section 1: The Homeostatic Principle and Feedback Mechanisms
\item Section 2: The Liver and Its Role in Homeostasis
\item Section 3: The Skin — Structure and Functions
\item Section 4: Excretory Products and Excretory Structures in Animals
\item Section 5: The Renal System — Structure, Urine Formation and Elimination
\item Section 6: Osmoregulation — Kidneys and Animal Adaptations
\item Section 7: Excretory Products of Plants, Integration and Evaluation
\item Integration activity
\item Methods and worked examples
\item Exercises
\item Solutions to the exercises
\item Summary sheet
\end{enumerate}
\end{multicols}
\end{sommaire}

\begin{objectifs}
\begin{itemize}
\item By the end of this chapter I can define homeostasis and explain its importance for cell survival.
\item By the end of this chapter I can describe the structure of the liver and explain its homeostatic roles (blood glucose regulation, deamination, detoxification, bile production, heat production).
\item By the end of this chapter I can explain negative and positive feedback mechanisms using biological examples.
\item By the end of this chapter I can describe the structure of mammalian skin and relate each structure to its functions, especially thermoregulation.
\item By the end of this chapter I can identify the main nitrogenous excretory products of animals (ammonia, urea, uric acid) and relate them to habitat and water availability.
\item By the end of this chapter I can describe the structure of the kidney and nephron and explain the stages of urine formation (ultrafiltration, selective reabsorption, secretion).
\item By the end of this chapter I can explain the role of ADH and the kidneys in osmoregulation, and compare osmoregulation in freshwater, marine and terrestrial animals.
\item By the end of this chapter I can state the excretory products of plants and how plants eliminate or store them.
\item By the end of this chapter I can interpret experimental data (graphs, tables) on blood glucose, body temperature or urine composition in GCE-style questions.
\end{itemize}
\end{objectifs}

\begin{prerequis}
\begin{itemize}
\item Cell structure and the concept of the internal environment (tissue fluid, blood plasma) from Lower Sixth.
\item Enzymes and the effect of temperature and pH on enzyme activity.
\item Transport in animals: composition and functions of blood.
\item Carbohydrates, proteins and lipids: structure and digestion (biomolecules).
\item Basic knowledge of hormones and the endocrine system.
\end{itemize}
\end{prerequis}

\newpage

\coursec{Section 1: The Homeostatic Principle and Feedback Mechanisms}

During the dry season in Maroua, a farmer works under a $40\ ^{\circ}\mathrm{C}$ sun all morning, yet his body temperature stays close to $37\ ^{\circ}\mathrm{C}$, while a tilapia in a drying pond near Bafoussam must survive water that becomes saltier by the day. After a heavy meal of fufu and eru, a student's blood sugar rises sharply, but within two hours it returns to normal without any conscious effort. How does the body keep its internal environment stable while the external environment changes so brutally? This chapter explores \cle{homeostasis} — the set of mechanisms, centred on the liver, skin and kidneys, that keep our internal conditions constant and allow survival. In this first section, we build the foundation: what the internal environment is, why it must be kept stable, and the logic of the \cle{feedback mechanisms} that achieve this stability.

\courssub{The Internal Environment and the Need for Stability}

\textbf{Observation.} A red blood cell placed in pure water swells and bursts; placed in very salty water it shrivels. Yet inside your body, the same cells bathe continuously in a fluid whose composition barely changes from morning to night, from the cold of Mount Cameroon to the heat of Waza. The cells never "see" the outside world directly — they live in a carefully regulated internal fluid.

Every cell in the body is surrounded by fluid. This fluid is not blood itself, but a set of related liquids that together form what we call the \cle{internal environment} (the French physiologist Claude Bernard called it the \emph{milieu int\'erieur}):

\begin{itemize}
\item \textbf{Blood plasma}: the liquid part of blood, circulating in vessels, carrying nutrients, gases, hormones and wastes.
\item \textbf{Tissue fluid (interstitial fluid)}: the fluid that leaks out of blood capillaries and directly bathes the cells. It is the \emph{immediate} environment of every cell — exchanges of \ce{O2}, \ce{CO2}, glucose and wastes occur across it.
\item \textbf{Lymph}: tissue fluid that has drained into lymph vessels; it is returned to the blood, so the three fluids are in constant communication.
\end{itemize}

Because cells exchange materials only with this internal fluid, its composition determines whether cells live or die. Three reasons explain why stability is vital:

\begin{enumerate}
\item \textbf{Enzyme activity.} Enzymes are proteins with precise three-dimensional shapes. A rise of a few degrees in temperature or a small shift in pH alters the bonds holding the enzyme's shape, changing the active site and slowing — or stopping — metabolism.
\item \textbf{Membrane integrity.} Cell membranes are held together by interactions that are sensitive to temperature and pH. Extreme conditions disrupt membranes, destroying the cell's control over what enters and leaves.
\item \textbf{Osmotic balance.} If the water potential of the tissue fluid falls (fluid becomes too concentrated), cells lose water by osmosis and shrivel; if it rises too much, cells take up water and may burst. Animal cells have no cell wall to protect them.
\end{enumerate}

\begin{definition}[Homeostasis]
\cle{Homeostasis} is the maintenance of a relatively constant internal environment despite changes in the external environment (and despite changes produced by the body's own activity). "Relatively constant" means that each controlled variable is kept within \textbf{narrow limits} around an average value (the \textbf{set point}), not at one fixed value.
\end{definition}

\begin{attention}[Common mistake: "exactly constant"]
Many candidates write that homeostasis keeps conditions "exactly constant" or "perfectly stable". This is \textbf{wrong} and loses marks. No variable is fixed: blood glucose rises after a meal and falls during exercise. The correct phrasing is that variables are kept \textbf{within narrow limits around a set point}. Examiners specifically reward the phrase "within narrow limits".
\end{attention}

\textbf{Variables controlled in mammals.} The principal variables under homeostatic control are:

\begin{center}
\begin{tabularx}{\linewidth}{|l|X|}
\hline
\rowcolor{popblueL} \textbf{Controlled variable} & \textbf{Why it must be controlled} \\
\hline
Body temperature & Enzyme reaction rates; membrane fluidity \\
\hline
Blood glucose concentration & Respiratory substrate for all cells; too high damages tissues, too low starves the brain \\
\hline
Blood water potential (osmotic concentration) & Prevents cells shrinking or bursting by osmosis \\
\hline
Blood pH & Enzyme shape and activity; blood pH is held very close to 7.4 \\
\hline
Blood pressure & Ensures adequate perfusion of tissues without damaging vessels \\
\hline
Blood \ce{CO2} and \ce{O2} levels & \ce{O2} for aerobic respiration; \ce{CO2} affects blood pH \\
\hline
\end{tabularx}
\end{center}

\courssub{The Components of a Homeostatic Control System}

\textbf{Intuition.} Think of the air-conditioning unit in a Douala office. A thermometer inside the unit \emph{detects} the room temperature; an electronic circuit \emph{compares} it with the desired value and \emph{switches} the compressor on or off; the compressor \emph{acts} to cool the room. Every biological control system has exactly the same architecture.

\begin{definition}[Components of a control system]
Any homeostatic control system contains:
\begin{itemize}
\item a \cle{receptor} (detector): a cell or structure that detects a change (the \textbf{stimulus}) in a controlled variable and sends information (nerve impulses or hormones) to the control centre;
\item a \cle{coordination / control centre}: usually a region of the brain (e.g.\ the hypothalamus) or an endocrine gland (e.g.\ the pancreas), which receives the information, compares the value with the \textbf{set point}, and coordinates a response;
\item an \cle{effector}: a muscle or gland that carries out the corrective response;
\item a \cle{response}: the action that returns the variable towards its set point.
\end{itemize}
\end{definition}

\begin{methode}[Analysing any homeostatic control loop in 4 steps]
Whenever the GCE asks you to "describe how X is regulated", structure your answer as follows:
\begin{enumerate}
\item \textbf{Identify the variable} being controlled and its set point (e.g.\ blood glucose, about $90\ \mathrm{mg}$ per $100\ \mathrm{cm^3}$ of blood).
\item \textbf{Name the specific receptor} that detects the change (e.g.\ the $\beta$-cells of the islets of Langerhans in the pancreas detect rising glucose).
\item \textbf{Name the control centre and the messenger} (nervous pathway or hormone) linking receptor to effector (e.g.\ insulin secreted into the blood).
\item \textbf{Name the specific effector and its response}, and state how this brings the variable back to the set point (e.g.\ liver and muscle cells convert glucose to glycogen, so blood glucose falls).
\end{enumerate}
\end{methode}

\begin{attention}[GCE tip: be specific]
Never write "the brain detects the change" or "the body responds". Examiners require the \textbf{specific} receptor (e.g.\ thermoreceptors in the hypothalamus and skin; chemoreceptors in the carotid bodies; osmoreceptors in the hypothalamus) and the \textbf{specific} effector (e.g.\ arterioles of the skin, sweat glands, the liver, the adrenal medulla). Vague answers score zero for those marking points.
\end{attention}

\courssub{Negative Feedback: The Engine of Stability}

\textbf{Observation.} After your meal of fufu and eru, blood glucose climbs. Two hours later it is back to normal, and it does \emph{not} keep falling until you faint. Something switched the correction off once it had done its job. That "something" is the essence of negative feedback: the response cancels the stimulus that caused it.

\begin{definition}[Negative feedback]
\cle{Negative feedback} is a control mechanism in which a change in a controlled variable triggers a response that \textbf{reverses} (counteracts) the initial change, returning the variable towards its set point. Because the response switches off the stimulus, the system is self-limiting and stable. Negative feedback is the mechanism underlying almost all homeostasis.
\end{definition}

The general model can be summarised as a loop:
\[
\text{stimulus} \rightarrow \text{receptor} \rightarrow \text{control centre} \rightarrow \text{effector} \rightarrow \text{response} \rightarrow \text{stimulus reduced}
\]

\begin{popfigure}
\begin{center}
\begin{tikzpicture}[
box/.style={draw=popblue, thick, rounded corners, fill=popblueL, text width=2.6cm, align=center, minimum height=1cm, font=\small},
arr/.style={-{Stealth[length=3mm]}, thick, popdark}]
\node[box, align=center] (stim) at (0,0){\textbf{Stimulus}\\ variable deviates from set point};
\node[box, align=center] (rec) at (4.6,0){\textbf{Receptor}\\ detects the change};
\node[box, align=center] (cc) at (9.2,0){\textbf{Control centre}\\ compares with set point};
\node[box, align=center] (eff) at (9.2,-3){\textbf{Effector}\\ muscle or gland};
\node[box, align=center] (resp) at (4.6,-3){\textbf{Response}\\ corrective action};
\draw[arr] (stim) -- (rec);
\draw[arr] (rec) -- (cc);
\draw[arr] (cc) -- (eff);
\draw[arr] (eff) -- (resp);
\draw[arr, poporange] (resp) -- (stim) node[midway, left, font=\small\itshape, text=poporange, align=center]{negative feedback:\\ change is reversed};
\end{tikzpicture}
\end{center}
\legende{General model of a negative feedback loop. The response reduces the original stimulus, so the loop switches itself off.}
\end{popfigure}

\begin{exemplebox}[Two canonical examples]
\textbf{Blood glucose.} A rise in blood glucose after a meal is detected by the $\beta$-cells of the islets of Langerhans (pancreas), which secrete \textbf{insulin}. Insulin causes the liver and muscles to convert glucose to glycogen and increases glucose uptake by cells; blood glucose falls back to the set point, and insulin secretion then decreases. A fall in glucose triggers the $\alpha$-cells to secrete \textbf{glucagon}, which stimulates the liver to break down glycogen — the same loop working in the opposite direction.

\textbf{Body temperature.} A rise in core temperature is detected by thermoreceptors in the hypothalamus and skin; the hypothalamus (control centre) sends nerve impulses to effectors: sweat glands secrete more sweat and skin arterioles dilate (vasodilation), increasing heat loss until temperature returns to $37\ ^{\circ}\mathrm{C}$.
\end{exemplebox}

\textbf{Set point, normal range and oscillation.} Because negative feedback only switches on \emph{after} a deviation has been detected, and takes time to act, the controlled variable never sits exactly on the set point. Instead it \textbf{oscillates} around it within a \textbf{normal range}: a small rise triggers a correction that slightly overshoots, producing a small fall, which triggers the opposite correction, and so on. This is exactly what the graph below shows for blood glucose after a meal.

\begin{popfigure}
\begin{center}
\begin{tikzpicture}
\begin{axis}[
width=12cm, height=6.5cm,
xlabel={Time after meal (min)},
ylabel={Blood glucose (mg per $100\ \mathrm{cm^3}$)},
xmin=0, xmax=150, ymin=60, ymax=150,
xtick={0,30,60,90,120,150},
ytick={70,90,110,130},
grid=both, grid style={popline, very thin},
axis line style={popdark},
]
\addplot[domain=0:150, samples=200, thick, popblue] {90 + 40*exp(-((x-30)^2)/500) - 12*exp(-((x-95)^2)/300)};
\addplot[domain=0:150, dashed, poporange, thick] {90};
\node[font=\small, text=poporange, anchor=west] at (axis cs:100,92) {set point ($\approx 90$)};
\addplot[domain=0:150, dotted, popmuted] {110};
\addplot[domain=0:150, dotted, popmuted] {70};
\node[font=\scriptsize, text=popmuted, anchor=west] at (axis cs:120,112) {upper limit};
\node[font=\scriptsize, text=popmuted, anchor=west] at (axis cs:120,72) {lower limit};
\end{axis}
\end{tikzpicture}
\end{center}
\legende{Blood glucose oscillating around its set point after a meal. Insulin brings the peak down; a slight undershoot is corrected by glucagon. The variable stays within narrow limits, never fixed.}
\end{popfigure}

\courssub{Positive Feedback: Amplification, Not Stability}

\begin{definition}[Positive feedback]
\cle{Positive feedback} is a control mechanism in which a change in a variable triggers a response that \textbf{amplifies} (reinforces) the initial change, pushing the variable further in the same direction. The process accelerates until an external event ends it. Positive feedback therefore \textbf{cannot} maintain a stable internal environment; it is rare in the body and is used only when a rapid, self-accelerating outcome is needed.
\end{definition}

\begin{exemplebox}[Two classic examples of positive feedback]
\textbf{Blood clotting.} When a vessel is damaged, platelets adhere and release chemicals that attract and activate \emph{more} platelets, which release \emph{more} chemicals — the clot grows explosively until the wound is sealed.

\textbf{Oxytocin in labour.} During childbirth, stretching of the cervix stimulates the pituitary to release \textbf{oxytocin}, which strengthens uterine contractions; stronger contractions stretch the cervix further, releasing more oxytocin. The loop escalates until the baby is born, which removes the stimulus and breaks the loop.
\end{exemplebox}

Notice the pattern: in both cases the loop ends only when an \emph{external} event (sealed wound, birth) interrupts it. Left alone, positive feedback would run to destruction — which is precisely why it is unsuitable for homeostasis.

\courssub{Nervous versus Hormonal Control}

Homeostatic corrections reach effectors by two routes: \textbf{nervous} (impulses along neurones) and \textbf{hormonal} (chemical messengers in the blood). Many responses, such as temperature control, use both.

\begin{center}
\begin{tabularx}{\linewidth}{|l|X|X|}
\hline
\rowcolor{popblueL} \textbf{Feature} & \textbf{Nervous control} & \textbf{Hormonal control} \\
\hline
Pathway & Along neurones (reflex arcs) & In the blood plasma \\
\hline
Speed & Very fast (milliseconds) & Slower (seconds to hours) \\
\hline
Duration of response & Short-lived & Longer-lasting \\
\hline
Target & Precise, localised effectors & Often widespread (all cells with receptors) \\
\hline
Example & Sweating, vasodilation, shivering & Insulin and glucagon in glucose control; ADH in water balance \\
\hline
\end{tabularx}
\end{center}

\begin{exoresolu}[Analysing a control loop (GCE style)]
\textbf{Statement.} During a football match in Yaound\'e, a player's body temperature rises from $37\ ^{\circ}\mathrm{C}$ to $38.5\ ^{\circ}\mathrm{C}$. Describe, using the four-step method, how the body returns the temperature to normal, and state the type of feedback involved.
\tcblower
\textbf{Solution.}
\begin{enumerate}
\item \textbf{Variable and set point:} core body temperature, set point about $37\ ^{\circ}\mathrm{C}$.
\item \textbf{Receptor:} thermoreceptors in the \textbf{hypothalamus} detect the rise in blood temperature (skin thermoreceptors also contribute).
\item \textbf{Control centre and messenger:} the \textbf{thermoregulatory centre of the hypothalamus} compares the value with the set point and sends \textbf{nerve impulses} to effectors.
\item \textbf{Effectors and response:} \textbf{sweat glands} increase sweat secretion (evaporation removes heat) and \textbf{arterioles supplying the skin} dilate (vasodilation), increasing blood flow to the surface and heat loss by radiation. Skeletal muscles stop any shivering.
\end{enumerate}
As temperature falls back to $37\ ^{\circ}\mathrm{C}$, the stimulus disappears and the responses are switched off: this is \textbf{negative feedback}, because the response reverses the initial change.
\end{exoresolu}

\begin{exercice}[Check your understanding]
\begin{enumerate}
\item Define homeostasis and explain why the phrase "within narrow limits" is essential to your definition.
\item For the control of blood glucose after a meal, name the receptor, the control centre, the hormone, the effector and the response.
\item Explain, with one example, why positive feedback cannot maintain a constant internal environment.
\item Give two differences between nervous and hormonal control, and one homeostatic example of each.
\end{enumerate}
\end{exercice}

\begin{corrige}[Exercise 1]
\begin{enumerate}
\item Homeostasis is the maintenance of a relatively constant internal environment despite external changes. "Within narrow limits" is essential because variables oscillate around a set point; they are never exactly fixed.
\item Receptor: $\beta$-cells of the islets of Langerhans; control centre: the pancreas itself (it both detects and responds); hormone: insulin; effectors: liver and muscle cells; response: conversion of glucose to glycogen and increased glucose uptake, so blood glucose falls.
\item Positive feedback amplifies the initial change, driving the variable ever further from its starting value (e.g.\ oxytocin strengthening contractions until birth). Stability requires reversal of change, not amplification.
\item Nervous: fast, short-lived, via neurones (e.g.\ sweating). Hormonal: slower, longer-lasting, via blood (e.g.\ insulin in glucose control).
\end{enumerate}
\end{corrige}

\begin{aretenir}[Key points of Section 1]
\begin{itemize}
\item The internal environment = blood plasma, tissue fluid and lymph; cells exchange materials only with it.
\item Homeostasis keeps variables \textbf{within narrow limits around a set point}, never exactly constant.
\item Every control loop has: receptor $\rightarrow$ control centre $\rightarrow$ effector $\rightarrow$ response.
\item \textbf{Negative feedback} reverses the change and produces stability (glucose, temperature); \textbf{positive feedback} amplifies it (clotting, labour) and cannot stabilise.
\item Controlled variables oscillate around the set point because feedback acts only after a deviation.
\item Nervous control is fast and localised; hormonal control is slower but longer-lasting.
\end{itemize}
\end{aretenir}

\coursec{Section 2: The Liver and Its Role in Homeostasis}

In Section 1 we established the \cle{homeostatic principle}: the internal environment of the body (blood glucose, temperature, water potential, pH, toxin levels) must be kept within narrow limits, and this is achieved by \cle{negative feedback} loops involving receptors, control centres and effectors. We now meet the single most important \emph{effector organ} in the whole of chemical homeostasis: the \cle{liver}. Almost every variable we listed in Section 1 — blood glucose, amino acid levels, toxin concentration, blood temperature — is corrected by the liver. If the body were a factory, the liver would be its chemical processing plant, its warehouse, its waste-treatment unit and its central heating system, all in one.

\courssub{Position and Gross Structure of the Liver}

The \cle{liver} is the largest internal organ of the body (about $1.5\ \mathrm{kg}$ in an adult human). It lies in the upper right of the abdomen, just beneath the diaphragm, protected by the lower ribs. It is dark reddish-brown, soft, and divided into two main lobes (a large right lobe and a smaller left lobe). Tucked beneath its lower surface is the \cle{gall bladder}, a small greenish sac.

\begin{propriete}[The liver: the body's main metabolic and homeostatic organ]
The liver is the principal \cle{metabolic} and \cle{homeostatic} organ of the body. It regulates the blood concentrations of glucose, amino acids and lipids; removes toxic substances; stores essential reserves; produces bile and heat; and recycles the components of worn-out red blood cells. Because so many homeostatic corrections are executed by the liver, severe liver damage threatens the stability of the entire internal environment.
\end{propriete}

\textbf{The double blood supply.} Most organs receive blood from one artery and drain it through one vein. The liver is exceptional: it receives blood from \emph{two} vessels.

\begin{itemize}
\item The \cle{hepatic artery} (a branch of the aorta) brings \textbf{oxygenated blood} to supply the liver's own cells with oxygen for their intense respiration.
\item The \cle{hepatic portal vein} brings \textbf{deoxygenated blood rich in absorbed nutrients} directly from the small intestine (and also from the spleen and pancreas). This is crucial: everything absorbed from a meal — glucose, amino acids, but also any toxins or alcohol swallowed — passes through the liver \emph{before} reaching the general circulation. The liver is therefore the body's first checkpoint on all absorbed material.
\end{itemize}

Blood leaves the liver through the \cle{hepatic vein}, which joins the inferior vena cava. Thus the blood reaching the heart has already been "adjusted" by the liver.

\begin{attention}[Do not confuse the hepatic portal vein with the hepatic vein]
The \textbf{hepatic portal vein} carries blood \emph{into} the liver from the gut; the \textbf{hepatic vein} carries blood \emph{out of} the liver towards the heart. A portal vein is a vein that links two capillary networks (here, the capillaries of the intestine and the sinusoids of the liver) instead of returning blood directly to the heart.
\end{attention}

\courssub{Microscopic Structure: The Liver Lobule}

Under the microscope, the liver is seen to be built from about a million tiny cylindrical units called \cle{lobules}, each roughly $1$--$2\ \mathrm{mm}$ across. Each lobule has:

\begin{itemize}
\item \textbf{Hepatocytes} (\cle{liver cells}): arranged in rows (plates) radiating from the centre like spokes of a wheel. These are the working cells that carry out all the metabolic reactions.
\item \textbf{Sinusoids}: wide, leaky blood capillaries running between the rows of hepatocytes. They carry the mixed blood from branches of the hepatic artery and hepatic portal vein (at the edge of the lobule) towards the centre. Their walls are so permeable that plasma bathes the hepatocytes directly, allowing rapid exchange of substances.
\item \textbf{Central vein}: at the centre of the lobule; it collects blood from the sinusoids and drains into the hepatic vein.
\item \textbf{Kupffer cells}: specialised macrophages fixed in the walls of the sinusoids. They engulf and destroy old red blood cells, bacteria and debris — the liver's cleaning squad.
\item \textbf{Bile canaliculi}: tiny channels between hepatocytes that collect the \cle{bile} they secrete and carry it towards the \emph{edge} of the lobule (opposite direction to blood flow) into bile ducts, which lead to the gall bladder.
\end{itemize}

\begin{popfigure}
\begin{center}
\begin{tikzpicture}[scale=1.05]
% central vein
\draw[fill=popblueL, draw=popblue, thick] (0,0) circle (0.45);
\node[popdark, font=\small] at (0,0) {Central};
\node[popdark, font=\small] at (0,-0.28) {vein};
% hepatocyte plates (radiating rows)
\foreach \a in {0,60,120,180,240,300}{
  \foreach \r in {0.85,1.25,1.65}{
    \draw[fill=popgold, draw=poporange] (\a:\r) circle (0.17);
  }
}
% sinusoids between plates
\foreach \a in {30,90,150,210,270,330}{
  \draw[popblue, thick, ->] (\a:2.0) -- (\a:0.62);
}
% bile canaliculi (outward)
\foreach \a in {0,120,240}{
  \draw[poppurple, thick, ->] (\a+18:0.7) -- (\a+18:1.95);
}
% boundary
\draw[popmuted, dashed] (0,0) circle (2.15);
% labels
\node[popblue, font=\small, align=center] at (3.1,1.4) {Sinusoid\\(blood flows in)};
\draw[popblue, thin] (2.6,1.2) -- (1.45,0.85);
\node[poporange, font=\small, align=center] at (3.1,-0.6) {Hepatocyte\\(liver cell)};
\draw[poporange, thin] (2.55,-0.45) -- (1.35,-0.35);
\node[poppurple, font=\small, align=center] at (-3.2,1.3) {Bile canaliculus\\(bile flows out)};
\draw[poppurple, thin] (-2.6,1.1) -- (-1.5,0.75);
\node[popgreen, font=\small, align=center] at (-3.2,-1.1) {Kupffer cell\\(in sinusoid wall)};
\draw[popgreen, thin] (-2.55,-0.95) -- (-1.15,-0.55);
\draw[fill=popgreen, draw=popgreen] (30:1.35) circle (0.09);
\end{tikzpicture}
\end{center}
\legende{Transverse section of a liver lobule. Blood flows \emph{inwards} through sinusoids from the hepatic artery and portal vein branches to the central vein, while bile flows \emph{outwards} in canaliculi towards the bile duct.}
\end{popfigure}

\begin{attention}[Blood and bile flow in opposite directions]
In a lobule, blood in the sinusoids flows from the \textbf{outside towards the central vein}, while bile in the canaliculi flows from the \textbf{centre towards the bile duct at the edge}. Examiners reward candidates who state the two directions correctly.
\end{attention}

\courssub{Regulation of Blood Glucose by the Liver}

After a meal of \emph{fufu} and \emph{ndol\'e}, digested carbohydrate floods the blood as glucose via the hepatic portal vein. Left unchecked, blood glucose would shoot far above the norm (about $90\ \mathrm{mg}$ per $100\ \mathrm{cm^{3}}$ of blood); hours later, during fasting or exercise, it would crash. The liver, directed by two pancreatic hormones, irons out these fluctuations.

\begin{definition}[Glycogenesis, glycogenolysis and gluconeogenesis]
\begin{itemize}
\item \cle{Glycogenesis}: the conversion of excess blood glucose into the storage polysaccharide \textbf{glycogen} in liver (and muscle) cells. Stimulated by \textbf{insulin}.
\item \cle{Glycogenolysis}: the breakdown of liver glycogen back into glucose, which is released into the blood. Stimulated by \textbf{glucagon} (and adrenaline).
\item \cle{Gluconeogenesis}: the synthesis of \emph{new} glucose from non-carbohydrate sources (amino acids, glycerol, lactic acid) when glycogen reserves run low, e.g.\ during prolonged fasting. Stimulated by glucagon.
\end{itemize}
\end{definition}

The control works exactly like the feedback loops of Section 1:

\begin{itemize}
\item \textbf{Blood glucose rises} (after a meal) $\rightarrow$ $\beta$-cells of the islets of Langerhans (pancreas) secrete \cle{insulin} $\rightarrow$ liver increases glucose uptake, glycogenesis, and glucose respiration, while glycogenolysis is inhibited $\rightarrow$ blood glucose \textbf{falls} back to the norm.
\item \textbf{Blood glucose falls} (fasting, exercise) $\rightarrow$ $\alpha$-cells secrete \cle{glucagon} $\rightarrow$ liver carries out glycogenolysis and gluconeogenesis $\rightarrow$ blood glucose \textbf{rises} back to the norm.
\end{itemize}

Note that muscle glycogen, unlike liver glycogen, cannot release glucose into the blood: muscle keeps its glycogen for its own respiration. Only the liver can top up blood glucose directly.

\begin{popfigure}
\begin{center}
\begin{tikzpicture}[node distance=0.55cm, font=\small,
box/.style={draw=popblue, fill=popblueL, rounded corners, align=center, text width=3.4cm, inner sep=4pt},
horm/.style={draw=poporange, fill=popgold, rounded corners, align=center, text width=2.6cm, inner sep=3pt},
eff/.style={draw=popgreen, fill=popgreenL, rounded corners, align=center, text width=3.6cm, inner sep=4pt},
arr/.style={->, thick, popdark}]
\node[box] (high) {Blood glucose \textbf{rises} (after a meal)};
\node[horm, below=of high] (ins) {Pancreas secretes \textbf{insulin}};
\node[eff, below=of ins] (gly) {Liver: glycogenesis + increased respiration of glucose};
\node[box, below=of gly] (norm1) {Blood glucose returns to \textbf{norm}};
\node[box, right=4.6cm of high] (low) {Blood glucose \textbf{falls} (fasting, exercise)};
\node[horm, below=of low] (glu) {Pancreas secretes \textbf{glucagon}};
\node[eff, below=of glu] (gol) {Liver: glycogenolysis + gluconeogenesis};
\node[box, below=of gol] (norm2) {Blood glucose returns to \textbf{norm}};
\draw[arr] (high) -- (ins);
\draw[arr] (ins) -- (gly);
\draw[arr] (gly) -- (norm1);
\draw[arr] (low) -- (glu);
\draw[arr] (glu) -- (gol);
\draw[arr] (gol) -- (norm2);
\draw[arr, dashed, poppurple] (norm1.south) .. controls (0,-7.6) .. (low.south);
\draw[arr, dashed, poppurple] (norm2.north) .. controls (7.2,1.2) .. (high.north);
\node[poppurple, align=center, text width=3cm] at (3.6,-6.9) {negative feedback};
\end{tikzpicture}
\end{center}
\legende{Blood glucose regulation: insulin and glucagon act antagonistically on the liver in two complementary negative feedback loops.}
\end{popfigure}

\begin{exoresolu}[Explaining a glucose tolerance curve]
A student drinks a solution containing $50\ \mathrm{g}$ of glucose. Her blood glucose rises from $90$ to $150\ \mathrm{mg}/100\ \mathrm{cm^{3}}$ within 30 minutes, then falls back to $90\ \mathrm{mg}/100\ \mathrm{cm^{3}}$ after about 2 hours. Explain the fall.
\tcblower
\textbf{Solution.} The absorbed glucose raises blood glucose above the set point. This is detected by the $\beta$-cells of the islets of Langerhans in the pancreas, which secrete \textbf{insulin} into the blood. Insulin reaches the liver (and muscles) and stimulates: (i) increased uptake of glucose by cells; (ii) \textbf{glycogenesis} — conversion of glucose to glycogen for storage; (iii) increased oxidation of glucose in respiration. As glucose is removed from the blood, its concentration falls back towards the norm; the falling level then reduces insulin secretion — a negative feedback loop. In a diabetic person lacking effective insulin, the level would remain high and glucose would appear in the urine.
\end{exoresolu}

\courssub{Deamination of Amino Acids and Urea Formation}

Unlike glucose and fats, excess \cle{amino acids} cannot be stored in the body. After a protein-rich meal (think of a generous plate of \emph{eru} with plenty of meat), surplus amino acids arrive at the liver, which processes them in two steps.

\begin{definition}[Deamination]
\cle{Deamination} is the removal of the amino group ($-\mathrm{NH_2}$) from an amino acid in the liver. The amino group is converted first to \textbf{ammonia} ($\mathrm{NH_3}$), which is highly toxic; the remaining carbon skeleton (keto-acid) can be respired for energy or converted into glucose or fat.
\end{definition}

Because ammonia is so toxic, the liver immediately converts it, with carbon dioxide, into the much less toxic \cle{urea}, via a sequence of reactions called the \cle{ornithine cycle} (or urea cycle). In outline:
\[
2\,\mathrm{NH_3} + \mathrm{CO_2} \;\longrightarrow\; \mathrm{CO(NH_2)_2} \;(\text{urea}) \;+\; \mathrm{H_2O}
\]
Urea passes into the blood and is later removed by the kidneys and excreted in urine (Section 5). This is a beautiful example of homeostatic integration: one organ (liver) \emph{detoxifies}, another (kidney) \emph{eliminates}.

\courssub{Detoxification by the Liver}

\cle{Detoxification} is the chemical conversion of poisonous substances into harmless (or less harmful) ones, usually so they can be excreted. Hepatocytes are packed with enzymes for this:

\begin{itemize}
\item \textbf{Alcohol}: ethanol is oxidised by the enzyme \emph{alcohol dehydrogenase} to ethanal, then to ethanoic acid, which enters respiration. Chronic heavy drinking overworks and kills hepatocytes, leading to \cle{cirrhosis} — replacement of liver tissue by useless fibrous scar tissue.
\item \textbf{Drugs}: medicines (and poisons) are modified so the kidneys can excrete them. This is why doctors adjust drug doses for patients with liver disease.
\item \textbf{Hydrogen peroxide} ($\mathrm{H_2O_2}$): a toxic by-product of metabolism, split by the enzyme \cle{catalase}, extremely abundant in the liver:
\[
2\,\mathrm{H_2O_2} \xrightarrow{\text{catalase}} 2\,\mathrm{H_2O} + \mathrm{O_2}
\]
\end{itemize}

\begin{experience}[Catalase in liver tissue]
Drop a small cube of fresh raw liver into hydrogen peroxide solution in a test tube. Vigorous effervescence occurs as oxygen is released; the gas relights a glowing splint. Boiled liver produces no bubbles, because catalase — an enzyme, hence a protein — has been denatured by heat. This classic practical demonstrates both the presence of catalase in liver and the enzyme's role in detoxifying $\mathrm{H_2O_2}$.
\end{experience}

\begin{savaistu}[Liver disease in Cameroon]
Viral \cle{hepatitis} (especially hepatitis B) is a significant public-health problem in Cameroon and can progress to cirrhosis and liver cancer. Vaccination against hepatitis B, avoiding unsterilised needles and sharp instruments, and moderating alcohol consumption are effective protective measures. A damaged liver cannot regulate blood glucose, detoxify poisons or make urea properly — a direct demonstration of how one organ underpins the whole internal environment.
\end{savaistu}

\courssub{Bile Production, Storage Functions and Recycling of Red Blood Cells}

\textbf{Bile production.} Hepatocytes continuously manufacture \cle{bile}, a greenish alkaline fluid containing bile salts, bile pigments, cholesterol and water. Bile contains no digestive enzyme; its key role is the \cle{emulsification of fats} in the duodenum: bile salts break large fat globules into tiny droplets, enormously increasing the surface area on which lipase can act. Its alkalinity also helps neutralise the acidic chyme arriving from the stomach, providing the alkaline pH that pancreatic and intestinal enzymes require.

\begin{attention}[Common mistake: bile is \emph{made} by the liver, \emph{stored} in the gall bladder]
Candidates frequently write that bile is "produced by the gall bladder". This is wrong. Bile is \textbf{produced by hepatocytes (liver)}; the \textbf{gall bladder only stores and concentrates} it, releasing it into the duodenum when fatty food arrives. Keep the two verbs distinct in the examination.
\end{attention}

\textbf{Storage.} The liver is the body's warehouse, storing: \textbf{glycogen} (the glucose reserve), \textbf{iron} (from recycled haemoglobin), and the fat-soluble \textbf{vitamins A, D, E, K} (plus vitamin $\mathrm{B_{12}}$). These reserves are released when the blood level falls — again, homeostasis.

\textbf{Recycling of red blood cells.} Red blood cells live only about 120 days. Worn-out cells are engulfed by \cle{Kupffer cells} (and by the spleen). The haemoglobin is split: the \textbf{iron} is recovered and stored for making new haemoglobin; the remaining haem portion is converted into the green-yellow pigment \cle{bilirubin}, which is excreted in bile (it gives faeces their brown colour). Nothing useful is wasted.

\courssub{The Liver as a Heat-Producing Organ}

Every reaction described above — glycogenesis, deamination, detoxification — releases energy, and much of it appears as \cle{heat}. Because of its enormous metabolic rate and large mass, the liver is one of the body's chief \textbf{thermogenic organs}, acting like a permanent central-heating boiler. The warm blood leaving the liver via the hepatic vein distributes this heat around the body, contributing to the maintenance of the constant body temperature of about $37\ ^{\circ}\mathrm{C}$ studied in Section 1. During cold exposure, increased metabolic activity in the liver (alongside shivering in muscles) raises heat production.

\begin{methode}[GCE tip: link every liver function to the variable it keeps constant]
In structured questions, examiners reward candidates who connect each liver function to a \emph{specific homeostatic variable}. Train yourself to answer in pairs:
\begin{enumerate}
\item Glycogenesis / glycogenolysis / gluconeogenesis $\rightarrow$ constant \textbf{blood glucose}.
\item Deamination + urea formation $\rightarrow$ constant \textbf{amino acid level}; removal of toxic \textbf{ammonia}.
\item Detoxification (alcohol, drugs, $\mathrm{H_2O_2}$) $\rightarrow$ constant, non-toxic \textbf{chemical composition} of blood.
\item Bile production $\rightarrow$ efficient fat digestion (supporting nutrient supply).
\item Storage of glycogen, iron, vitamins $\rightarrow$ stable \textbf{reserves} released on demand.
\item Breakdown of red blood cells $\rightarrow$ recycling of \textbf{iron}, removal of cell debris.
\item High metabolic rate $\rightarrow$ constant \textbf{body temperature}.
\end{enumerate}
A sentence such as "the liver regulates blood glucose by converting excess glucose to glycogen under insulin control" scores far more than "the liver controls sugar".
\end{methode}

\begin{exercice}[Consolidation]
\begin{enumerate}
\item Draw and label a liver lobule, indicating the direction of blood flow and of bile flow.
\item Explain, using the terms glycogenesis and glycogenolysis, how the liver responds to (a) a carbohydrate-rich meal and (b) a 24-hour fast.
\item Outline the fate of an excess amino acid molecule from the moment it reaches the liver until its nitrogen leaves the body.
\item State two differences between the roles of the liver and the gall bladder in relation to bile.
\end{enumerate}
\end{exercice}

\begin{corrige}[Exercise 1]
\textbf{1.} The diagram should show a roughly hexagonal lobule with a central vein, radiating rows of hepatocytes, sinusoids carrying blood \emph{inwards} from branches of the hepatic artery and hepatic portal vein to the central vein, bile canaliculi carrying bile \emph{outwards} to a bile duct, and Kupffer cells in the sinusoid walls (see the figure above).
\par\textbf{2.} (a) After the meal, blood glucose rises; insulin is secreted; the liver converts glucose to glycogen (\textbf{glycogenesis}) and increases glucose respiration, lowering blood glucose. (b) During the fast, blood glucose falls; glucagon is secreted; the liver breaks glycogen down to glucose (\textbf{glycogenolysis}) and later makes glucose from amino acids and glycerol (\textbf{gluconeogenesis}), raising blood glucose.
\par\textbf{3.} The amino acid is \textbf{deaminated}: its amino group becomes ammonia; ammonia is combined with $\mathrm{CO_2}$ in the \textbf{ornithine cycle} to form \textbf{urea}; urea enters the blood, is filtered out by the \textbf{kidneys} and eliminated in \textbf{urine}. The carbon skeleton is respired or converted to glucose/fat.
\par\textbf{4.} The liver \emph{produces} bile (hepatocytes secrete it) and also excretes bile pigments into it; the gall bladder only \emph{stores and concentrates} bile and releases it into the duodenum. Bile production continues even if the gall bladder is surgically removed.
\end{corrige}

\begin{aretenir}[Key points of Section 2]
\begin{itemize}
\item The liver has a \textbf{double blood supply}: oxygenated blood via the hepatic artery, nutrient-rich blood via the hepatic portal vein; it drains via the hepatic vein.
\item Its functional unit is the \textbf{lobule}: hepatocytes, sinusoids, Kupffer cells, bile canaliculi, central vein.
\item It regulates blood glucose through \textbf{glycogenesis, glycogenolysis and gluconeogenesis}, under \textbf{insulin} and \textbf{glucagon} — a classic negative feedback system.
\item It \textbf{deaminates} excess amino acids and converts toxic ammonia to \textbf{urea} (ornithine cycle).
\item It \textbf{detoxifies} alcohol, drugs and hydrogen peroxide (catalase); damage leads to hepatitis and cirrhosis.
\item It \textbf{produces} bile (the gall bladder only stores it), stores glycogen, iron and vitamins, recycles haemoglobin into iron and bilirubin, and generates metabolic \textbf{heat} for thermoregulation.
\end{itemize}
\end{aretenir}

\coursec{Section 3: The Skin — Structure and Functions}

In the previous section we saw how the liver acts as the body's great internal chemical factory, regulating blood glucose, detoxifying poisons and producing heat through its intense metabolic activity. But the liver works inside a body that is constantly exchanging matter and energy with the outside world. The organ that stands at this frontier — the largest organ of the human body, covering about $1.8\ \mathrm{m^2}$ in an adult — is the \cle{skin}. The skin is far more than a simple wrapping: it is a sense organ, an excretory organ, a protective armour and, above all, the principal effector organ of \cle{thermoregulation}, one of the finest examples of homeostasis you will meet at Advanced Level.

\begin{popfigure}
\begin{tikzpicture}[scale=1.0, font=\small]
% ---- Epidermis layers ----
\fill[popgold!35] (0,4.6) rectangle (10,5.6);
\fill[poporange!30] (0,3.9) rectangle (10,4.6);
\fill[poppink!35] (0,3.0) rectangle (10,3.9);
% wavy top of cornified layer
\draw[thick, popdark] (0,5.6) .. controls (1.2,5.85) and (2.2,5.4) .. (3.4,5.65)
 .. controls (4.6,5.9) and (5.6,5.45) .. (6.8,5.7)
 .. controls (8.0,5.9) and (9.0,5.5) .. (10,5.75);
\draw[thick, popdark] (0,4.6) -- (10,4.6);
\draw[thick, popdark] (0,3.9) -- (10,3.9);
\draw[thick, popdark] (0,3.0) -- (10,3.0);
\node at (9.0,5.15) {\footnotesize cornified layer};
\node at (9.0,4.25) {\footnotesize granular layer};
\node at (9.0,3.45) {\footnotesize Malpighian layer};
% ---- Dermis ----
\fill[popblueL!60] (0,0.4) rectangle (10,3.0);
\draw[thick, popdark] (0,3.0) -- (10,3.0);
% ---- Subcutaneous fat ----
\fill[popgold!50] (0,-0.6) rectangle (10,0.4);
\draw[thick, popdark] (0,0.4) -- (10,0.4);
\foreach \x in {0.8,2.2,3.6,5.0,6.4,7.8,9.2}{\draw[popdark] (\x,-0.1) circle (0.28);}
\node at (5.0,-0.95) {\footnotesize subcutaneous adipose tissue (fat)};
% ---- Blood capillaries ----
\draw[very thick, poppink!80!black] (0.3,1.1) .. controls (2,1.5) and (3,0.7) .. (5,1.2) .. controls (7,1.7) and (8,0.8) .. (9.7,1.3);
\node[poppink!80!black] at (1.6,0.75) {\footnotesize blood capillaries};
% ---- Sweat gland ----
\draw[thick, poppurple] (2.0,3.0) -- (2.0,4.9);
\draw[thick, poppurple] (2.0,1.6) .. controls (1.6,1.3) and (2.4,1.1) .. (2.0,0.8) .. controls (1.6,0.5) and (2.4,0.3) .. (2.0,0.1);
\node[poppurple, align=center] at (0.9,2.3) {\footnotesize sweat gland\\ \footnotesize and duct};
\draw[->, poppurple] (1.4,2.1) -- (1.95,1.4);
% sweat pore
\fill[poppurple] (2.0,4.95) circle (0.06);
% ---- Hair follicle ----
\draw[very thick, popdark] (5.2,5.75) -- (5.7,1.2);
\draw[thick, popdark] (5.35,1.6) .. controls (5.0,1.1) and (6.2,1.0) .. (5.95,1.6);
\node[popdark] at (5.2,6.0) {\footnotesize hair};
\node[popdark, align=center] at (7.1,1.35) {\footnotesize hair follicle};
\draw[->, popdark] (6.6,1.3) -- (5.8,1.35);
% erector pili muscle
\draw[thick, popgreen!60!black] (5.55,1.9) -- (4.4,2.9);
\node[popgreen!60!black, align=center] at (3.4,2.6) {\footnotesize erector pili\\ \footnotesize muscle};
\draw[->, popgreen!60!black] (3.9,2.55) -- (4.75,2.5);
% sebaceous gland
\fill[poporange!70!black] (6.15,2.35) circle (0.22);
\fill[poporange!70!black] (6.35,2.15) circle (0.18);
\node[poporange!70!black, align=center] at (8.1,2.4) {\footnotesize sebaceous\\ \footnotesize gland};
\draw[->, poporange!70!black] (7.5,2.35) -- (6.4,2.3);
% ---- Sensory receptors ----
\fill[popblue!80!black] (8.3,2.9) circle (0.12);
\node[popblue!80!black] at (8.3,3.25) {\footnotesize touch receptor};
\draw[thick, popblue!80!black] (8.3,2.78) -- (8.3,2.0);
\draw[thick, popblue!80!black] (8.3,1.55) ellipse (0.18 and 0.3);
\node[popblue!80!black, align=center] at (8.3,0.9) {\footnotesize pressure\\ \footnotesize receptor};
\draw[thick, popblue!80!black] (8.3,1.25) -- (8.3,1.85);
% labels epidermis/dermis
\node[popdark] at (0.7,5.9) {\footnotesize \textbf{EPIDERMIS}};
\node[popdark] at (0.55,2.0) {\footnotesize \textbf{DERMIS}};
\end{tikzpicture}
\legende{Vertical section through mammalian skin, showing the three layers of the epidermis, the dermis with its structures, and the subcutaneous adipose tissue.}
\end{popfigure}

\courssub{Gross Structure of the Skin}

A vertical section through mammalian skin reveals two main layers, the \cle{epidermis} above and the \cle{dermis} below, resting on a layer of \cle{subcutaneous adipose tissue} (fat).

\textbf{The epidermis}\par
The epidermis is the thin outer layer, made of epithelial cells and containing no blood vessels. From the inside outwards, three layers are distinguished:

\begin{itemize}
\item The \cle{Malpighian layer} (germinative layer): the deepest layer, resting on the dermis. Its cells are alive and divide continuously by \textbf{mitosis}, constantly producing new cells that are pushed upwards. Its cells also contain the brown pigment \cle{melanin}, which absorbs harmful ultraviolet (UV) radiation and protects the underlying tissues — particularly important under the intense tropical sun of Cameroon.
\item The \cle{granular layer}: here the cells, pushed up from below, begin to die. Their cytoplasm fills with granules of \cle{keratin}, a tough, waterproof fibrous protein.
\item The \cle{cornified layer} (stratum corneum): the outermost layer, made of flat, dead, fully keratinised cells. It forms a tough, flexible, waterproof barrier and is continually worn away and replaced from below.
\end{itemize}

This continuous production of cells in the Malpighian layer, their keratinisation as they rise, and their final shedding is called \cle{keratinisation}. It explains how the skin repairs itself after a cut and why the surface is always renewed.

\textbf{The dermis}\par
The dermis is the thicker, living layer beneath the epidermis, made of connective tissue with strong collagen and elastic fibres. It contains:

\begin{itemize}
\item \textbf{Blood capillaries}, supplied by arterioles, which provide nutrients and, crucially, adjust heat loss by changing their diameter;
\item \textbf{Sweat glands}: coiled tubes whose ducts carry sweat to a pore on the surface;
\item \textbf{Hair follicles}, pits from which hairs grow, each with an \textbf{erector pili muscle} attached;
\item \textbf{Sebaceous glands}, which secrete \emph{sebum}, an oily substance that keeps hair and skin supple and waterproof and inhibits microbes;
\item \textbf{Sensory receptors} for touch, pressure, temperature (heat and cold) and pain, connected to sensory neurones;
\item Nerve endings and, beneath the dermis, the \textbf{subcutaneous adipose tissue}, which insulates the body and stores energy.
\end{itemize}

\courssub{Functions of the Skin}

\begin{itemize}
\item \textbf{Protection.} The cornified layer is a mechanical barrier against knocks and abrasion; it is waterproof, preventing both excessive water loss and the entry of pathogens; sebum has mild antiseptic properties; and melanin absorbs UV radiation, protecting the DNA of deeper cells.
\item \textbf{Sensitivity.} Its many receptors make the skin a major sense organ, informing the nervous system about touch, pressure, pain and temperature.
\item \textbf{Synthesis of vitamin D.} When UV light falls on the skin, a cholesterol derivative in the epidermis is converted into \cle{vitamin D}, essential for calcium absorption and healthy bones.
\item \textbf{Excretion.} Sweat contains water, salts (mainly sodium chloride) and small amounts of \cle{urea}; the skin is therefore a minor excretory organ.
\item \textbf{Thermoregulation} — the skin's key homeostatic role, developed below.
\end{itemize}

\begin{attention}[GCE tip — structures versus functions]
Examiners frequently ask candidates to \emph{relate skin structures to their functions}. Do not confuse the two. Learn this table:
\begin{center}
\begin{tabularx}{\linewidth}{|l|X|}
\hline
\rowcolor{popblueL} \textbf{Structure} & \textbf{Function} \\
\hline
Cornified layer & Mechanical and microbial protection; waterproofing \\
\hline
Melanin (Malpighian layer) & Absorbs UV radiation \\
\hline
Blood capillaries & Vasodilation/vasoconstriction to control heat loss \\
\hline
Sweat glands & Secrete sweat; evaporation cools the body \\
\hline
Hair + erector pili muscle & Trap an insulating layer of air when erected \\
\hline
Sebaceous gland & Sebum waterproofs and protects against microbes \\
\hline
Sensory receptors & Detect touch, pressure, temperature, pain \\
\hline
Adipose tissue & Insulation; energy store \\
\hline
\end{tabularx}
\end{center}
\end{attention}

\courssub{Thermoregulation: the Skin's Homeostatic Role}

Mammals are \cle{endotherms}: they generate their own heat by respiration and must keep their deep (core) body temperature close to $37\ ^\circ\mathrm{C}$, whatever the external conditions, because enzymes work best within a narrow temperature range. A fall of a few degrees slows metabolism dangerously; a rise of a few degrees denatures enzymes.

\begin{definition}[Thermoregulation]
\cle{Thermoregulation} is the maintenance of a constant internal body temperature, independent of fluctuations in the external temperature, by balancing heat production and heat loss. In mammals it is controlled by the \cle{hypothalamus} of the brain, which acts as the thermoregulatory centre (the body's ``thermostat'').
\end{definition}

The hypothalamus receives information about blood temperature from its own thermoreceptors and about skin temperature from receptors in the dermis. It then sends nerve impulses (and triggers hormonal responses) to effectors in the skin and muscles. Two key effector responses involve the diameter of arterioles supplying the skin capillaries:

\begin{definition}[Vasodilation and vasoconstriction]
\cle{Vasodilation} is the widening of the arterioles supplying the skin capillaries: more warm blood flows near the surface, so more heat is lost by radiation. \cle{Vasoconstriction} is the narrowing of these arterioles: blood is diverted away from the surface, reducing heat loss.
\end{definition}

\textbf{Responses to cold}\par
\begin{itemize}
\item \textbf{Vasoconstriction} reduces blood flow near the skin surface, cutting heat loss;
\item \textbf{Shivering}: rapid, involuntary contractions of skeletal muscles release heat from respiration;
\item \textbf{Erection of hairs} (erector pili muscles contract): the raised hairs trap a thicker layer of \emph{still air}, and it is this trapped air — a poor conductor — that insulates the body (in humans, with our sparse hair, this response is largely vestigial, visible only as ``goose pimples'');
\item \textbf{Increased metabolic rate}: the hormones \cle{adrenaline} and \cle{thyroxine} increase the rate of respiration in cells, generating more heat.
\end{itemize}

\textbf{Responses to heat}\par
\begin{itemize}
\item \textbf{Vasodilation} brings more blood to the surface, increasing heat loss by radiation;
\item \textbf{Increased sweating}: sweat evaporates from the skin surface, and evaporation removes heat (see method box below);
\item \textbf{Relaxation of erector pili muscles}: hairs lie flat, the insulating air layer is lost, and heat escapes more easily;
\item \textbf{Reduced metabolic rate}: less thyroxine is secreted, so less heat is generated by respiration.
\end{itemize}

\begin{methode}[Explaining how sweating cools the body]
To earn full marks, explain sweating in four logical steps:
\begin{enumerate}
\item Sweat glands secrete sweat (mainly water) onto the skin surface;
\item The most energetic water molecules escape as vapour — this is \textbf{evaporation};
\item Evaporation requires \cle{latent heat of vaporisation}, which is taken \emph{from the skin and the blood in its capillaries};
\item The skin and blood therefore lose heat and the body is cooled.
\end{enumerate}
The essential phrase is: \emph{``evaporation of sweat takes latent heat from the body''}. Simply writing ``sweat cools the body'' scores little.
\end{methode}

\begin{attention}[Common mistakes]
\begin{itemize}
\item Sweat glands do \textbf{not} ``open and close'' — they secrete more or less sweat. (Pores are passive openings.)
\item Hairs do \textbf{not} ``trap heat''. They trap a layer of \emph{still air}, which is a poor conductor of heat and therefore reduces heat loss. Always mention the insulating air layer.
\item Vasoconstriction does not stop blood flow to the skin; it \emph{reduces} it. And blood itself is not ``cooled in the skin and sent back'' as the main mechanism — the point is how much heat is \emph{lost} at the surface.
\item Do not confuse the \emph{receptor} (skin thermoreceptors, hypothalamus) with the \emph{effector} (muscles, glands, arterioles).
\item Sweating only cools the body if the sweat \emph{evaporates}; in very humid air (as in Douala during the rainy season) evaporation is slow and cooling is less efficient.
\end{itemize}
\end{attention}

\begin{popfigure}
\begin{tikzpicture}[font=\small, node distance=0.55cm,
 box/.style={draw=popdark, rounded corners=2pt, fill=popblueL, align=center, text width=3.1cm, minimum height=0.9cm},
 stim/.style={draw=popdark, rounded corners=2pt, fill=popgold!40, align=center, text width=2.6cm, minimum height=0.9cm},
 eff/.style={draw=popdark, rounded corners=2pt, fill=popgreenL, align=center, text width=3.4cm}]
% central detector/coordinator
\node[box] (hypo) {\textbf{Hypothalamus} (thermoregulatory centre)};
\node[stim, above left=0.7cm and 1.6cm of hypo] (hot) {Body temperature \textbf{rises} above $37^\circ$C};
\node[stim, below left=0.7cm and 1.6cm of hypo] (cold) {Body temperature \textbf{falls} below $37^\circ$C};
\node[eff, above right=0.7cm and 1.4cm of hypo] (hoteff) {Vasodilation; sweating increases; hairs lie flat};
\node[eff, below right=0.7cm and 1.4cm of hypo] (coldeff) {Vasoconstriction; shivering; hairs erect; more thyroxine/adrenaline};
\node[box, right=1.5cm of hypo, fill=poppink!30] (norm) {Temperature returns to $37^\circ$C};
\draw[->, thick, popdark] (hot) -- (hypo);
\draw[->, thick, popdark] (cold) -- (hypo);
\draw[->, thick, popdark] (hypo) -- (hoteff);
\draw[->, thick, popdark] (hypo) -- (coldeff);
\draw[->, thick, popdark] (hoteff.east) -- ++(0.5,0) |- (norm.north east);
\draw[->, thick, popdark] (coldeff.east) -- ++(0.5,0) |- (norm.south east);
\draw[->, thick, poppurple, dashed] (norm.west) -- (hypo.east) node[midway, above, poppurple]{\footnotesize negative feedback};
\end{tikzpicture}
\legende{Negative feedback loop of thermoregulation. A rise or fall in body temperature is detected by the hypothalamus, which triggers effector responses that return the temperature to the norm — the response cancels the stimulus.}
\end{popfigure}

This is a classic \cle{negative feedback} loop: any deviation from the set point ($37\ ^\circ\mathrm{C}$) triggers responses that \emph{reverse} the deviation. The skin provides both the peripheral receptors and most of the effectors.

\begin{exoresolu}[Explaining a response to heat]
\textbf{Statement.} A student in Douala plays football at midday in full sun. Describe and explain the changes that occur in his skin to prevent his body temperature from rising dangerously.

\tcblower
\textbf{Solution.}
\begin{enumerate}
\item \textbf{Detection:} as his blood temperature rises above $37\ ^\circ\mathrm{C}$, thermoreceptors in the hypothalamus (and heat receptors in the skin) detect the rise; the hypothalamus sends nerve impulses to skin effectors.
\item \textbf{Vasodilation:} arterioles supplying the skin capillaries widen, so more warm blood flows close to the surface and more heat is lost by radiation to the surroundings.
\item \textbf{Sweating:} sweat glands secrete more sweat onto the skin surface; the sweat evaporates, and evaporation takes latent heat of vaporisation from the skin and blood, cooling the body.
\item \textbf{Hairs:} the erector pili muscles relax, the hairs lie flat, the insulating layer of trapped air is removed, and heat loss increases.
\item \textbf{Feedback:} as the body temperature falls back to $37\ ^\circ\mathrm{C}$, the hypothalamus reduces these responses — negative feedback.
\end{enumerate}
\end{exoresolu}

\begin{exercice}[Cold weather in the highlands]
During a cold night on the slopes of Mount Cameroon, a hiker's body temperature begins to fall. (a) Name the organ that detects this fall. (b) Describe four responses of the body to the cold, stating for each one how it conserves or produces heat. (c) Explain why this control mechanism is described as negative feedback.
\end{exercice}

\begin{corrige}[Exercise 1]
(a) The \textbf{hypothalamus} (with cold receptors in the skin).
(b) \textbf{Vasoconstriction} reduces blood flow near the skin surface, reducing heat loss by radiation. \textbf{Shivering} — rapid muscle contractions — releases heat from respiration. \textbf{Erection of hairs} traps a layer of still air, a poor conductor, which insulates the body. \textbf{Increased secretion of adrenaline and thyroxine} raises the metabolic (respiration) rate, producing more heat.
(c) It is negative feedback because the responses triggered by the fall in temperature act to \emph{reverse} that fall, bringing the temperature back to the set point; the response cancels the stimulus that caused it.
\end{corrige}

\begin{aretenir}[Key points]
\begin{itemize}
\item Skin = epidermis (Malpighian, granular, cornified layers) + dermis (capillaries, sweat glands, hair follicles, sebaceous glands, receptors) + subcutaneous fat.
\item Keratinisation continually renews the protective cornified layer; melanin protects against UV.
\item Skin functions: protection, sensitivity, vitamin D synthesis, minor excretion, and thermoregulation.
\item The hypothalamus is the thermoregulatory centre; the skin provides receptors and effectors.
\item Cold: vasoconstriction, shivering, hair erection (trapped air), more thyroxine/adrenaline. Heat: vasodilation, sweating (evaporation removes latent heat), hairs flat.
\item Temperature control is a negative feedback loop: the response cancels the deviation.
\end{itemize}
\end{aretenir}

\coursec{Section 4: Excretory Products and Excretory Structures in Animals}

In Section 3 we saw that the skin is far more than a protective envelope: it is also an \emph{excretory} organ, since sweat carries away water, salts and small amounts of urea. This observation raises a wider question. Sweat is only one small stream of waste leaving the body. What are all the wastes that a living organism must eliminate, where do they come from, and through which structures do they leave? That is the subject of this section. We will discover a beautiful logic: the chemical nature of an animal's nitrogenous waste is dictated largely by the availability of water in its habitat, and evolution has produced a remarkable diversity of excretory structures, from the tiny contractile vacuole of \emph{Amoeba} to the kidneys of mammals.

\courssub{Excretion, Egestion and Secretion: Three Words You Must Never Confuse}

Consider what happens after a meal of \emph{ndol\'e} and \emph{miondo}. The digested nutrients are absorbed into the blood and used by cells; the undigested fibres pass out through the anus; meanwhile, every cell in the body is continuously producing carbon dioxide, water and nitrogenous wastes from respiration and protein metabolism. Only some of these materials count as \emph{excretory products}.

\begin{definition}[Excretion, egestion and secretion]
\begin{itemize}
\item \cle{Excretion} is the removal from the body of the \textbf{waste products of metabolism} (substances produced by chemical reactions inside cells), together with substances present in excess, before they accumulate to toxic levels.
\item \cle{Egestion} (defaecation) is the elimination of \textbf{undigested food material} from the gut through the anus. This material has never entered the body cells or taken part in metabolism.
\item \cle{Secretion} is the production and release of a \textbf{useful substance} (e.g.\ hormones, enzymes, mucus) by a cell or gland.
\end{itemize}
\end{definition}

\begin{attention}[Faeces are NOT an excretory product]
The commonest examination error is to list faeces among excretory products. Faeces consist mainly of undigested food that passed through the alimentary canal \emph{without ever entering the cells}; its removal is \cle{egestion}, not excretion. Note, however, that bile pigments in faeces (breakdown products of haemoglobin, excreted by the liver into the gut) \emph{are} genuine excretory products.
\end{attention}

\courssub{The Main Excretory Products of Animals}

Metabolism generates several categories of waste that must be removed:

\begin{center}
\begin{tabularx}{\linewidth}{|l|X|X|}
\hline
\rowcolor{popblueL} \textbf{Excretory product} & \textbf{Source (metabolic origin)} & \textbf{Main route of removal} \\
\hline
Carbon dioxide & Aerobic respiration in all cells & Lungs (gills in fish) \\
\hline
Water & Respiration; also excess from food and drink & Kidneys, skin, lungs \\
\hline
Nitrogenous wastes (ammonia, urea, uric acid) & Deamination of excess amino acids (mainly in the liver) & Kidneys (or equivalent organs) \\
\hline
Mineral salts (e.g.\ Na$^+$, Cl$^-$) & Excess intake in the diet & Kidneys, skin (sweat) \\
\hline
Bile pigments & Breakdown of haemoglobin from old red blood cells & Liver, via bile into the gut \\
\hline
\end{tabularx}
\end{center}

Of these, the \cle{nitrogenous wastes} deserve special attention, because the \emph{form} in which nitrogen is excreted varies dramatically between animal groups — and for a very good reason.

\courssub{Nitrogenous Wastes: Ammonia, Urea and Uric Acid}

\textbf{Where nitrogenous waste comes from.} The body cannot store excess amino acids. In the liver, surplus amino acids undergo \cle{deamination}: the amino group ($-\mathrm{NH_2}$) is removed, forming \textbf{ammonia} (\ce{NH3}), while the remaining carbon skeleton is respired or converted to carbohydrate or fat. Ammonia is extremely toxic — even small concentrations damage cells — so it must either be eliminated immediately or converted into a safer form.

The three nitrogenous excretory products form a logical series:

\begin{propriete}[The toxicity--water cost relationship]
The more \textbf{toxic} a nitrogenous waste is, the more \textbf{water} is needed to dilute and excrete it safely, and the \textbf{less energy} is needed to produce it:
\begin{itemize}
\item \cle{Ammonia}: highly toxic, extremely soluble; must be diluted in large volumes of water and excreted continuously. Little energy needed to produce it (direct product of deamination).
\item \cle{Urea}: moderately toxic, quite soluble; can be tolerated at low concentration, so less water is needed. Energy is required to convert ammonia into urea in the liver (ornithine cycle).
\item \cle{Uric acid}: almost non-toxic, almost insoluble; excreted as a semi-solid paste or crystals, so very little water is lost. It costs the most energy to synthesise.
\end{itemize}
The conversion of ammonia into urea or uric acid is therefore an \textbf{energy--water trade-off}: the animal spends metabolic energy to save water.
\end{propriete}

\begin{popfigure}
\begin{center}
\begin{tikzpicture}[
  box/.style={draw=popblue, thick, rounded corners, fill=popblueL, text width=3.6cm, align=center, minimum height=4.6cm, inner sep=4pt},
  hdr/.style={font=\bfseries, text=popdark}
]
\node[box, align=center] (am) at (0,0){\textcolor{popdark}{\textbf{Ammonia}}\\[2pt]
Toxicity: \textbf{very high}\\
Solubility: \textbf{very high}\\
Water needed: \textbf{very much}\\
Energy cost: \textbf{low}\\[3pt]
\emph{Examples:} bony fish, tadpoles, aquatic invertebrates};
\node[box, align=center] (ur) at (4.7,0){\textcolor{popdark}{\textbf{Urea}}\\[2pt]
Toxicity: \textbf{moderate}\\
Solubility: \textbf{high}\\
Water needed: \textbf{moderate}\\
Energy cost: \textbf{moderate}\\[3pt]
\emph{Examples:} mammals, adult amphibians, sharks};
\node[box, align=center] (ua) at (9.4,0){\textcolor{popdark}{\textbf{Uric acid}}\\[2pt]
Toxicity: \textbf{very low}\\
Solubility: \textbf{very low}\\
Water needed: \textbf{very little}\\
Energy cost: \textbf{high}\\[3pt]
\emph{Examples:} birds, insects, reptiles, land snails};
\draw[-{Stealth}, very thick, poporange] (-1.9,2.9) -- (11.3,2.9)
  node[midway, above, text=popdark, align=center] {decreasing toxicity, solubility and water loss};
\draw[-{Stealth}, very thick, popgreen] (-1.9,-2.9) -- (11.3,-2.9)
  node[midway, below, text=popdark, align=center] {increasing energy cost of synthesis};
\end{tikzpicture}
\end{center}
\legende{Comparison of the three nitrogenous excretory products. Moving from ammonia to uric acid, toxicity, solubility and water requirement fall, while the energy cost of synthesis rises.}
\end{popfigure}

\courssub{Nitrogenous Waste and Habitat: An Evolutionary Trade-off}

The pattern above only makes sense when linked to the animal's environment:

\begin{itemize}
\item \textbf{Ammonotelic animals} excrete \cle{ammonia}. They are almost all \textbf{aquatic} (bony fish, tadpoles, many aquatic invertebrates): surrounded by water, they can afford to flush ammonia out continuously across the gills or body surface.
\item \textbf{Ureotelic animals} excrete \cle{urea}. \textbf{Mammals} and \textbf{adult amphibians} convert toxic ammonia into urea in the liver, an adaptation to life where water must be conserved but is not critically scarce.
\item \textbf{Uricotelic animals} excrete \cle{uric acid}. \textbf{Birds, insects and reptiles} live where water is precious (flight, deserts, life inside an egg shell). Excreting an insoluble paste conserves nearly all body water. In birds, this also suits development inside the \emph{cleidoic egg}, where toxic soluble wastes could not be diluted away.
\end{itemize}

\begin{savaistu}[Why bird droppings are white]
The white part of a bird's droppings on a car windscreen in Yaound\'e is not faeces — it is \textbf{uric acid}, the bird's nitrogenous waste, excreted as a thick paste through the cloaca together with the darker faecal material. A flying bird cannot afford to carry litres of diluting water, so evolution chose the expensive but water-saving option.
\end{savaistu}

\begin{methode}[Deducing an animal's habitat from its nitrogenous waste]
\begin{enumerate}
\item \textbf{Identify the excretory product} named or described in the question (ammonia, urea or uric acid).
\item \textbf{Apply the toxicity--water rule}: highly toxic and soluble waste $\Rightarrow$ much water needed $\Rightarrow$ aquatic habitat; insoluble waste $\Rightarrow$ water conservation $\Rightarrow$ terrestrial or aerial habitat.
\item \textbf{Name the strategy}: ammonotelic, ureotelic or uricotelic.
\item \textbf{Justify with the trade-off}: state that converting ammonia to urea or uric acid costs energy but saves water, which is advantageous in the stated habitat.
\item \textbf{Give one named example} of an animal matching the strategy.
\end{enumerate}
\end{methode}

\begin{exoresolu}[GCE-style question on nitrogenous excretion]
\textit{An animal excretes its nitrogenous waste as a whitish, semi-solid paste and loses almost no water in the process. (a) Name the nitrogenous waste and the type of excretory strategy. (b) Give one named example of such an animal. (c) Explain the advantage and the cost of this strategy.}
\tcblower
\textbf{(a)} The waste is \textbf{uric acid}; the animal is \textbf{uricotelic}.\\
\textbf{(b)} Example: a bird (e.g.\ the domestic fowl), an insect (e.g.\ the cockroach) or a reptile (e.g.\ a lizard).\\
\textbf{(c)} \emph{Advantage:} uric acid is almost insoluble and virtually non-toxic, so it can be excreted with minimal water loss — a decisive advantage for terrestrial life, flight, or development inside a shelled egg. \emph{Cost:} synthesising uric acid from ammonia requires considerable metabolic energy. The strategy is therefore an energy--water trade-off: energy is spent to conserve water.
\end{exoresolu}

\courssub{Excretory Structures in Invertebrates}

Long before kidneys evolved, invertebrates developed a variety of structures to get rid of wastes and control water balance.

\textbf{Contractile vacuoles in protozoa.} Freshwater protozoa such as \emph{Amoeba} and \emph{Paramecium} live in a hypotonic medium: water constantly enters by osmosis. The \cle{contractile vacuole} collects this excess water (with some dissolved wastes) and periodically contracts, expelling the fluid through the cell surface. In \emph{Paramecium} there are two fixed contractile vacuoles with radiating canals; in \emph{Amoeba} a single vacuole moves with the cytoplasm. Without it, the cell would swell and burst.

\textbf{Flame cells in flatworms.} Flatworms such as \emph{Planaria} and the tapeworm possess a network of fine tubules ending in hollow cells called \cle{flame cells}. Each flame cell bears a tuft of cilia whose beating (resembling a flickering flame) drives fluid along the tubules to excretory pores on the body surface. The system removes water and nitrogenous waste (mainly ammonia).

\textbf{Nephridia in annelids.} In the earthworm, almost every segment contains a pair of coiled tubes called \cle{nephridia}. Each nephridium has a ciliated funnel, the \emph{nephrostome}, opening into the coelom of the preceding segment; cilia draw coelomic fluid into a long coiled tubule surrounded by capillaries, where useful substances are reabsorbed into the blood, and the remaining waste fluid leaves through a ventral pore, the \emph{nephridiopore}. Note the principle of \emph{filtration followed by reabsorption} — the same logic used by the vertebrate kidney.

\textbf{Malpighian tubules in insects.} Insects such as the cockroach possess numerous slender blind-ending tubes, the \cle{Malpighian tubules}, which float in the haemolymph of the body cavity and open into the gut at the junction of midgut and hindgut. Their cells actively transport potassium ions, uric acid and other wastes from the haemolymph into the tubule lumen; water follows by osmosis. The fluid passes into the hindgut and rectum, where most of the water and valuable salts are \textbf{reabsorbed}, so that nearly dry uric acid is voided with the faeces — a superb water-conserving device.

\begin{popfigure}
\begin{center}
\begin{tikzpicture}[scale=1.0, every node/.style={align=center}]
% Gut tube
\draw[very thick, popdark, rounded corners] (0.5,0) -- (3.2,0) -- (3.2,-0.9) -- (6.4,-0.9) -- (6.4,-1.8) -- (9.6,-1.8);
\draw[very thick, popdark, rounded corners] (0.5,-0.55) -- (3.2,-0.55) -- (3.2,-1.45) -- (6.4,-1.45) -- (6.4,-2.35) -- (9.6,-2.35);
\node[text=popdark] at (1.8,0.35) {\small midgut};
\node[text=popdark] at (4.8,-1.85) {\small hindgut};
\node[text=popdark] at (8.2,-2.75) {\small rectum};
% Malpighian tubules (wavy blind tubes above)
\foreach \x/\h in {3.6/2.6, 4.2/3.1, 4.8/2.8, 5.4/3.3, 6.0/2.7}{
  \draw[thick, poporange] (\x,-0.1) to[out=90,in=-90] (\x-0.25,\h*0.5) to[out=90,in=-90] (\x+0.15,\h);
}
\node[text=poporange, font=\small] at (4.8,3.75) {Malpighian tubules (blind ends)};
% Haemolymph label
\node[text=popblue, font=\small, align=center] at (1.6,2.2){haemolymph\\ (body cavity)};
% Arrows: uptake from haemolymph, flow to gut
\draw[-{Stealth}, thick, popblue] (2.3,1.9) -- (3.6,1.5);
\draw[-{Stealth}, thick, popgreen] (5.6,1.2) -- (5.6,-0.35);
\node[text=popgreen, font=\small, align=center] at (7.9,0.6){wastes + water\\ enter gut};
% Reabsorption arrow in rectum
\draw[-{Stealth}, thick, poppurple] (8.2,-2.1) -- (8.2,-0.6);
\node[text=poppurple, font=\small, text width=3.2cm] at (11.6,-1.1) {water and salts\\ reabsorbed\\ in rectum};
\end{tikzpicture}
\end{center}
\legende{The Malpighian tubule system of an insect. Blind-ending tubules bathed in haemolymph extract uric acid and salts; the fluid drains into the gut, and the rectum reabsorbs most of the water, so nearly dry waste is eliminated.}
\end{popfigure}

\courssub{Excretory Organs in Vertebrates: An Overview}

In vertebrates, excretion is shared between several organs, each handling particular products:

\begin{center}
\begin{tabularx}{\linewidth}{|l|X|}
\hline
\rowcolor{popblueL} \textbf{Organ} & \textbf{Excretory role} \\
\hline
Kidneys & Remove nitrogenous waste (urea), excess water and mineral salts as urine; central to osmoregulation (Section 6). \\
\hline
Skin & Sweat glands excrete water, salts and traces of urea (Section 3). \\
\hline
Lungs & Excrete carbon dioxide and water vapour during expiration. \\
\hline
Liver & \emph{Produces} urea by deamination and the ornithine cycle; excretes bile pigments from haemoglobin breakdown into the gut; detoxifies many poisons and drugs. \\
\hline
\end{tabularx}
\end{center}

The \cle{liver} thus plays a double role: it manufactures the very urea that the kidneys then remove, and it excretes bile pigments — the only major excretory products that leave the body via the gut.

\begin{aretenir}[Key points of Section 4]
\begin{itemize}
\item Excretion removes \emph{metabolic} wastes; egestion removes undigested food (faeces are \emph{not} excretory); secretion releases useful substances.
\item The three nitrogenous wastes form a series: \textbf{ammonia} (very toxic, very soluble, much water, low energy cost) $\rightarrow$ \textbf{urea} (intermediate) $\rightarrow$ \textbf{uric acid} (almost insoluble, minimal water, high energy cost).
\item Habitat dictates strategy: aquatic animals are ammonotelic, mammals and adult amphibians ureotelic, birds/insects/reptiles uricotelic.
\item Invertebrate excretory structures: contractile vacuoles (protozoa), flame cells (flatworms), nephridia (annelids), Malpighian tubules (insects).
\item Vertebrate excretory organs: kidneys, skin, lungs and liver.
\end{itemize}
\end{aretenir}

\begin{savaistu}[GCE tip: one example per strategy]
Examiners frequently ask for a \emph{named} example. Memorise at least one for each: \textbf{ammonotelic} — bony fish (e.g.\ tilapia from Lake Nyos); \textbf{ureotelic} — human (or any mammal); \textbf{uricotelic} — domestic fowl (or cockroach). Writing "aquatic animals" or "birds" without a named species can cost the example mark.
\end{savaistu}

\begin{exercice}[Checking your understanding]
\begin{enumerate}
\item Distinguish clearly between excretion, egestion and secretion, giving one example of each.
\item A tadpole excretes ammonia, but the adult frog excretes urea. Explain this change in terms of habitat and the toxicity--water relationship.
\item Copy and complete: flame cells are to flatworms as \ldots\ are to insects, and as \ldots\ are to earthworms.
\item Explain why excreting uric acid is advantageous to a bird developing inside a shelled egg, and state the metabolic cost of this strategy.
\end{enumerate}
\end{exercice}

\begin{corrige}[Exercise of Section 4]
\begin{enumerate}
\item \textbf{Excretion}: removal of metabolic wastes, e.g.\ urea removed by the kidneys. \textbf{Egestion}: removal of undigested food via the anus, e.g.\ faeces. \textbf{Secretion}: release of a useful substance, e.g.\ insulin from the pancreas.
\item The tadpole is fully aquatic: water is abundant, so highly toxic ammonia can be diluted and flushed out continuously at little energy cost. The adult frog spends much time on land where water must be conserved; it therefore pays the energy cost of converting ammonia into the less toxic urea, which needs far less water for its elimination.
\item Flame cells are to flatworms as \textbf{Malpighian tubules} are to insects, and as \textbf{nephridia} are to earthworms.
\item Inside the shelled (cleidoic) egg, no water can be exchanged with the outside, so a soluble toxic waste like ammonia would accumulate and kill the embryo. Uric acid is almost insoluble and non-toxic, so it is stored safely as crystals until hatching, with virtually no water loss. The cost is the large amount of metabolic energy needed to synthesise uric acid from ammonia.
\end{enumerate}
\end{corrige}

Having surveyed the excretory products and the structures that handle them across the animal kingdom, we are now ready to examine in detail the most sophisticated excretory machine of all: the vertebrate \textbf{renal system}, the subject of Section 5.

\coursec{Section 5: The Renal System — Structure, Urine Formation and Elimination}

In Section 4 we surveyed the variety of excretory products (ammonia, urea, uric acid) and the excretory structures found across the animal kingdom — from flame cells in flatworms to Malpighian tubules in insects. We now turn to the \emph{mammalian} excretory system, which is the model required for the GCE Advanced Level. The mammalian kidney is a masterpiece of biological engineering: it filters about \SI{180}{L} of blood plasma every day, yet returns \SI{99}{\%} of the filtrate to the bloodstream, producing only \SI{1}{--}\SI{2}{L} of urine. Understanding how this happens — ultrafiltration, selective reabsorption, and tubular secretion — is central to homeostasis.

\courssub{The Mammalian Urinary System: Organs and Blood Supply}

The urinary system consists of two kidneys, two ureters, the urinary bladder, and the urethra. The kidneys are bean-shaped organs lying against the posterior abdominal wall, one on each side of the vertebral column, at the level of the first to third lumbar vertebrae. The right kidney sits slightly lower than the left because of the liver above it.

\begin{popfigure}
\begin{tikzpicture}[scale=0.7, every node/.style={font=\small}]
% Kidneys
\draw[popdark, thick, fill=poporange!10] (1,5) .. controls (1,6.5) and (3,7) .. (4,6.5) .. controls (5,6) and (5,4) .. (4,3.5) .. controls (3,3) and (1,3.5) .. (1,5);
\draw[popdark, thick, fill=poporange!10] (10,5) .. controls (10,6.5) and (12,7) .. (13,6.5) .. controls (14,6) and (14,4) .. (13,3.5) .. controls (12,3) and (10,3.5) .. (10,5);
% Renal arteries and veins
\draw[popred, thick] (4,5) -- (5.5,5) -- (5.5,4.5);
\draw[popblue, thick] (4,4.5) -- (5.5,4.5) -- (5.5,5);
\draw[popred, thick] (10,5) -- (8.5,5) -- (8.5,4.5);
\draw[popblue, thick] (10,4.5) -- (8.5,4.5) -- (8.5,5);
% Aorta and vena cava
\draw[popred, very thick] (5.5,4.5) -- (5.5,0.5);
\draw[popblue, very thick] (8.5,4.5) -- (8.5,0.5);
% Ureters
\draw[popdark, thick] (4,3.5) .. controls (4.5,2) and (5.5,1) .. (5.5,0.5);
\draw[popdark, thick] (13,3.5) .. controls (12.5,2) and (11.5,1) .. (11.5,0.5);
% Bladder
\draw[popdark, thick, fill=popyellow!20] (5.5,0.5) .. controls (5.5,0) and (8.5,0) .. (8.5,0.5) .. controls (9,0.5) and (9,-0.5) .. (8.5,-1.5) .. controls (8,-2) and (6,-2) .. (6,-1.5) .. controls (5.5,-0.5) and (5.5,0) .. (5.5,0.5);
% Urethra
\draw[popdark, thick] (7,-1.5) -- (7,-2.5);
% Labels
\node[popred, left] at (4,5) {Renal artery};
\node[popblue, left] at (4,4.5) {Renal vein};
\node[popdark, above] at (2.5,6.5) {Left kidney};
\node[popdark, above] at (11.5,6.5) {Right kidney};
\node[popdark, right] at (10,5) {Renal artery};
\node[popdark, right] at (10,4.5) {Renal vein};
\node[popred, right] at (5.5,2) {Aorta};
\node[popblue, left] at (8.5,2) {Vena cava};
\node[popdark] at (7,0) {Urinary bladder};
\node[popdark] at (7,-2) {Urethra};
\node[popdark] at (5,1.5) {Ureter};
\node[popdark] at (12,1.5) {Ureter};
\end{tikzpicture}
\legende{The human urinary system showing kidneys, ureters, bladder, urethra and major blood vessels.}
\end{popfigure}

Each kidney receives blood via the \cle{renal artery} (a direct branch of the abdominal aorta) and returns filtered blood via the \cle{renal vein} to the inferior vena cava. The high blood pressure in the renal artery is essential for the first stage of urine formation: ultrafiltration. The renal artery enters the kidney at the \cle{hilum}, branching into segmental, interlobar, arcuate, and interlobular arteries before reaching the \cle{afferent arterioles} that supply individual nephrons.

\courssub{Gross Structure of the Kidney}

A longitudinal section of the kidney reveals three distinct regions:

\begin{itemize}
\item \cle{Cortex} — the outer, darker region containing the glomeruli, Bowman's capsules, and the proximal and distal convoluted tubules.
\item \cle{Medulla} — the inner, paler region consisting of conical \cle{renal pyramids} whose bases face the cortex and apexes (renal papillae) project into the pelvis. The pyramids contain loops of Henlé and collecting ducts.
\item \cle{Renal pelvis} — a funnel-shaped cavity that collects urine from the papillae via minor and major calyces and drains into the ureter.
\end{itemize}

\begin{popfigure}
\begin{tikzpicture}[scale=0.8, every node/.style={font=\small}]
% Kidney outline
\draw[popdark, thick, fill=poporange!5] (0,0) .. controls (0,2) and (1,3) .. (3,3) .. controls (5,3) and (6,2) .. (6,0) .. controls (6,-1) and (5,-1.5) .. (3,-1.5) .. controls (1,-1.5) and (0,-1) .. (0,0);
% Cortex
\draw[popdark, fill=poporange!20] (0,0) .. controls (0,2) and (1,3) .. (3,3) .. controls (5,3) and (6,2) .. (6,0) .. controls (5.5,0) and (4.5,0.5) .. (3,0.5) .. controls (1.5,0.5) and (0.5,0) .. (0,0);
% Medulla (pyramids)
\foreach \x in {0.8, 1.8, 2.8, 3.8, 4.8} {
  \draw[popdark, fill=popyellow!30] (\x,0.5) -- (\x+0.4,0.5) -- (3,-1) -- (3,-1.2) -- cycle;
}
% Pelvis
\draw[popdark, fill=popblue!20] (3,-1.2) .. controls (3,-1.5) and (4,-1.5) .. (4,-1.2) .. controls (4,-1) and (3,-1) .. (3,-1.2);
% Ureter
\draw[popdark, thick] (4,-1.2) -- (5,-1.2) -- (5,-2);
% Labels
\node[popdark] at (1.5,2.5) {Cortex};
\node[popdark] at (3,-0.5) {Medulla (pyramids)};
\node[popdark] at (4.5,-1) {Renal pelvis};
\node[popdark] at (5.5,-1.5) {Ureter};
\node[popdark] at (0.5,1) {Renal capsule};
% Blood vessels (simplified)
\draw[popred, thick] (3,3) -- (3,2.5);
\draw[popblue, thick] (3,2.5) -- (3,2);
\end{tikzpicture}
\legende{Longitudinal section of a mammalian kidney.}
\end{popfigure}

\courssub{The Nephron: Functional Unit of the Kidney}

Each kidney contains about one million \cle{nephrons}. A nephron is a microscopic tubule closed at one end (Bowman's capsule) and opening at the other into a collecting duct. It is intimately associated with a capillary network. There are two types of nephrons: \cle{cortical nephrons} (85\%, glomeruli in outer cortex, short loops of Henlé) and \cle{juxtamedullary nephrons} (15\%, glomeruli near medulla, long loops of Henlé extending deep into medulla — crucial for concentrating urine).

\begin{propriete}[The Nephron as the Functional Unit]
The nephron is the structural and functional unit of the kidney. All the processes of urine formation — ultrafiltration, selective reabsorption, and tubular secretion — occur within a single nephron and its associated blood vessels. A mammalian nephron consists of:
\begin{enumerate}
\item \cle{Bowman's capsule} — a cup-shaped dilation surrounding the \cle{glomerulus} (a knot of capillaries). The space between the visceral and parietal layers is the \cle{capsular space} (Bowman's space).
\item \cle{Proximal convoluted tubule (PCT)} — highly coiled, lined with a simple cuboidal epithelium with dense \cle{microvilli} (brush border) giving an enormous surface area for reabsorption.
\item \cle{Loop of Henlé} — a U-shaped tube with a \cle{descending limb} (thin, permeable to water) and an \cle{ascending limb} (thin then thick, impermeable to water, actively transports Na$^+$, Cl$^-$).
\item \cle{Distal convoluted tubule (DCT)} — coiled, lined with simple cuboidal epithelium with fewer microvilli; site of hormonal fine-tuning.
\item \cle{Collecting duct} — receives filtrate from several nephrons, runs through the medulla to the pelvis; permeability to water regulated by ADH.
\end{enumerate}
The associated blood supply: \cle{afferent arteriole} $\to$ \cle{glomerulus} $\to$ \cle{efferent arteriole} $\to$ \cle{peritubular capillaries} (around cortical nephrons) and \cle{vasa recta} (around juxtamedullary loops of Henlé) $\to$ \cle{renal vein}.
\end{propriete}

\begin{popfigure}
\begin{tikzpicture}[scale=0.65, every node/.style={font=\small}]
% Afferent arteriole
\draw[popred, thick] (0,5) -- (2,5) node[midway, above] {Afferent arteriole};
% Glomerulus
\draw[popred, thick, fill=popred!20] (2,4) .. controls (2,5.5) and (4,5.5) .. (4,4) .. controls (4,2.5) and (2,2.5) .. (2,4);
\node[popred] at (3,4) {Glomerulus};
% Bowman's capsule
\draw[popblue, thick, fill=popblue!10] (1.5,4.5) .. controls (1.5,6) and (4.5,6) .. (4.5,4.5) .. controls (4.5,3) and (1.5,3) .. (1.5,4.5);
\node[popblue] at (3,5.5) {Bowman's capsule};
% Efferent arteriole
\draw[popblue, thick] (4,4) -- (6,4) node[midway, above] {Efferent arteriole};
% Peritubular capillaries
\draw[poppurple, thick] (6,4) .. controls (7,4) and (8,3) .. (8,2) .. controls (8,1) and (7,0) .. (6,0) .. controls (5,0) and (4,1) .. (4,2) .. controls (4,3) and (5,3.5) .. (6,3.5);
% PCT
\draw[popdark, thick, fill=popgreen!10] (3,3) .. controls (2,2) and (1,1) .. (1,0) .. controls (1,-0.5) and (2,-0.5) .. (2,0) .. controls (2,1) and (3,2) .. (3,3);
\node[popdark] at (1.5,1.5) {PCT};
% Loop of Henle
\draw[popdark, thick] (2,0) -- (2,-2) -- (4,-2) -- (4,0);
\node[popdark] at (3,-1.5) {Loop of Henle};
% DCT
\draw[popdark, thick, fill=poporange!10] (4,0) .. controls (5,1) and (6,2) .. (6,3) .. controls (6,3.5) and (5,3.5) .. (5,3) .. controls (5,2) and (4,1) .. (4,0);
\node[popdark] at (5.5,2) {DCT};
% Collecting duct
\draw[popdark, thick, fill=poppink!10] (6,3) -- (6,5) -- (7,5) -- (7,-2.5);
\node[popdark] at (7.5,3) {Collecting duct};
% Vasa recta
\draw[poppurple, thick, dashed] (6,2) -- (6,-2) -- (7,-2) -- (7,2);
% Labels for stages
\node[popred, rotate=90] at (0.5,4) {\textbf{Ultrafiltration}};
\node[popgreen, rotate=90] at (0.5,1) {\textbf{Selective reabsorption}};
\node[poporange, rotate=90] at (5.5,1.5) {\textbf{Secretion \& adjustment}};
\end{tikzpicture}
\legende{A nephron with its blood supply, annotated with the three stages of urine formation.}
\end{popfigure}

\courssub{Stage 1: Ultrafiltration in the Bowman's Capsule}

\begin{definition}[Ultrafiltration]
Ultrafiltration is the process by which blood plasma is forced under high pressure through the filtration barrier of the glomerulus into the lumen of Bowman's capsule, producing \cle{glomerular filtrate}. The barrier consists of three layers: (1) the \cle{fenestrated endothelium} of glomerular capillaries (pores \SI{70}{--}\SI{100}{nm}), (2) the \cle{basement membrane} (acellular meshwork of collagen and glycoproteins, negatively charged, acts as the main molecular sieve), and (3) the \cle{podocytes} (specialised epithelial cells of Bowman's capsule with interdigitating foot processes leaving \cle{filtration slits} of \SI{25}{--}\SI{30}{nm} bridged by a \cle{slit diaphragm}).
\end{definition}

\begin{aretenir}[Why is pressure high in the glomerulus?]
\begin{itemize}
\item The \cle{afferent arteriole} is wider than the \cle{efferent arteriole} $\to$ resistance to outflow $>$ resistance to inflow.
\item Blood enters the glomerulus at systemic arterial pressure (\SI{\sim 60}{mmHg}) and leaves at lower pressure.
\item Net filtration pressure (NFP) $\approx$ \SI{10}{--}\SI{15}{mmHg}, calculated as:
\[
\text{NFP} = P_{\text{GC}} - P_{\text{BC}} - \pi_{\text{GC}}
\]
where $P_{\text{GC}}$ = glomerular capillary hydrostatic pressure (\SI{\sim 55}{mmHg}), $P_{\text{BC}}$ = Bowman's capsule hydrostatic pressure (\SI{\sim 15}{mmHg}), $\pi_{\text{GC}}$ = glomerular capillary oncotic pressure (\SI{\sim 30}{mmHg}).
\end{itemize}
\end{aretenir}

\begin{propriete}[Composition of Glomerular Filtrate vs Blood Plasma]
The filtrate is essentially \textbf{plasma without proteins}. Because the basement membrane and podocytes retain molecules larger than \SI{\sim 68}{kDa} (and negatively charged molecules are further repelled), the filtrate contains:
\begin{itemize}
\item Water, inorganic ions (Na$^+$, K$^+$, Cl$^-$, HCO$_3^-$, Ca$^{2+}$, Mg$^{2+}$, PO$_4^{3-}$), glucose, amino acids, urea, uric acid, creatinine — \emph{all at concentrations nearly identical to plasma}.
\item \textbf{No (or trace) proteins} (albumin, globulins), \textbf{no blood cells}, \textbf{no platelets}.
\end{itemize}
The \cle{glomerular filtration rate (GFR)} in a healthy adult is \SI{\sim 125}{mL/min} (\SI{\sim 180}{L/day}), tightly regulated by autoregulation (myogenic response, tubuloglomerular feedback), sympathetic nerves, and hormones (angiotensin II, ANP).
\end{propriete}

\courssub{Stage 2: Selective Reabsorption}

\begin{definition}[Selective Reabsorption]
Selective reabsorption is the process by which useful substances are transported from the tubular fluid back into the peritubular capillaries. It is \emph{selective} because the tubule epithelium chooses what to reabsorb and in what quantities, according to the body's needs. Reabsorption can be \cle{active} (requires ATP, against electrochemical gradient) or \cle{passive} (diffusion, osmosis, facilitated diffusion, down gradients).
\end{definition}

\begin{itemize}
\item \textbf{Proximal Convoluted Tubule (PCT)} — reabsorbs \textbf{all glucose, all amino acids}, \textbf{$\sim$65--70\% of filtered water}, \textbf{$\sim$65--70\% of Na$^+$, Cl$^-$, HCO$_3^-$, K$^+$}, and variable amounts of Ca$^{2+}$, Mg$^{2+}$, PO$_4^{3-}$. Mechanisms:
  \begin{itemize}
  \item \cle{Active transport} of Na$^+$ out of the tubular cell across the basolateral membrane via \cle{Na$^+$/K$^+$-ATPase} (3 Na$^+$ out, 2 K$^+$ in per ATP) creates a low intracellular Na$^+$ concentration.
  \item \cle{Co-transport (symport)} of glucose/amino acids with Na$^+$ across the apical membrane (microvilli increase surface area $\sim$40-fold). SGLT2 for glucose, various amino acid transporters.
  \item \cle{Counter-transport (antiport)}: Na$^+$/H$^+$ exchanger (NHE3) secretes H$^+$ for Na$^+$ reabsorption; H$^+$ combines with HCO$_3^-$ to form H$_2$CO$_3$ $\to$ CO$_2$ + H$_2$O (catalysed by carbonic anhydrase), CO$_2$ diffuses into cell, reforms HCO$_3^-$ which exits basolaterally.
  \item Water follows by \cle{osmosis} (obligatory water reabsorption) through aquaporin-1 channels and paracellularly.
  \item Mitochondria in PCT cells provide ATP for active transport; basal infoldings increase basolateral surface area.
  \end{itemize}
\item \textbf{Loop of Henlé} — creates a \cle{medullary osmotic gradient} (counter-current multiplier) enabling water reabsorption in the collecting duct under ADH control.
  \begin{itemize}
  \item \cle{Descending limb}: thin epithelium, highly permeable to water (aquaporin-1), impermeable to solutes. As filtrate descends, water leaves by osmosis into hypertonic medulla $\to$ tubular fluid becomes progressively more concentrated (up to \SI{1200}{mOsm/L} at bend).
  \item \cle{Ascending limb}: thin then thick segment; \textbf{impermeable to water}. Thick ascending limb actively reabsorbs Na$^+$, K$^+$, 2Cl$^-$ via \cle{NKCC2} co-transporter (furosemide target). This dilutes tubular fluid (down to \SI{100}{mOsm/L}) and adds solutes to medullary interstitium.
  \end{itemize}
\item \textbf{Distal Convoluted Tubule (DCT)} — reabsorbs \SI{\sim 5}{\%} of filtered Na$^+$ (via NCC co-transporter, thiazide target), Cl$^-$, Ca$^{2+}$ (regulated by parathyroid hormone). Impermeable to water without ADH.
\item \textbf{Collecting Duct} — final adjustment under hormonal control:
  \begin{itemize}
  \item \cle{ADH (vasopressin)}: binds V2 receptors $\to$ cAMP $\to$ insertion of \cle{aquaporin-2} water channels into apical membrane $\to$ water reabsorption into hypertonic medulla $\to$ concentrated urine.
  \item \cle{Aldosterone}: acts on principal cells $\to$ increases ENaC (epithelial Na$^+$ channel) and Na$^+$/K$^+$-ATPase $\to$ Na$^+$ reabsorption, K$^+$ and H$^+$ secretion.
  \item \cle{ANP (atrial natriuretic peptide)}: inhibits Na$^+$ reabsorption $\to$ natriuresis, diuresis.
  \end{itemize}
\end{itemize}

\begin{attention}[Common Mistake]
Do \textbf{not} say ``urea is filtered out of the kidney''. Urea is filtered \textbf{out of the blood in the glomerulus} into the nephron. Much of it stays in the tubule and is excreted; some is reabsorbed passively (especially in the collecting duct and thin ascending limb). The kidney \emph{excretes} urea; it does not filter urea out of the kidney.
\end{attention}

\courssub{Stage 3: Tubular Secretion and Final Adjustment}

In the DCT and collecting duct, certain substances are actively transported \emph{from} the peritubular capillaries \emph{into} the tubular lumen. This is \cle{tubular secretion}. It serves to:
\begin{itemize}
\item Eliminate excess H$^+$ (acid-base balance) via H$^+$-ATPase and H$^+$/K$^+$-ATPase in intercalated cells; also NH$_4^+$ secretion (ammoniagenesis in PCT).
\item Eliminate excess K$^+$ (via ROMK channels in principal cells, driven by aldosterone).
\item Remove certain drugs, toxins, and metabolic wastes (e.g., creatinine, penicillin, urate) that were not filtered efficiently or need additional clearance.
\item Fine-tune Na$^+$/water balance under aldosterone and ADH.
\end{itemize}
The final urine flows into the renal pelvis, down the ureters (propelled by \cle{peristalsis} of smooth muscle in ureter walls), and is stored in the bladder until \cle{micturition} (voluntary relaxation of the external urethral sphincter, somatic control via pudendal nerve).

\courssub{Composition of Normal Urine}

Normal urine is about \SI{95}{\%} water and \SI{5}{\%} solutes. Typical daily output: \SI{1}{--}\SI{2}{L} (minimum \SI{\sim 500}{mL} to excrete solutes). Key constituents:
\begin{itemize}
\item \cle{Urea} — main nitrogenous waste (\SI{20}{--}\SI{30}{g/day}, \SI{\sim 300}{--}\SI{600}{mmol/L}).
\item \cle{Inorganic salts} — NaCl, phosphates, sulphates, K$^+$, Mg$^{2+}$, Ca$^{2+}$.
\item \cle{Creatinine, uric acid} — minor nitrogenous wastes.
\item \textbf{No glucose, no protein} (or trace amounts below detection: protein < \SI{150}{mg/day}, glucose < \SI{0.8}{mmol/day}).
\item pH typically \SI{4.5}{--}\num{8.0} (average \num{6.0}), specific gravity \SI{1.003}{--}\num{1.030}.
\end{itemize}

\begin{exemplebox}[Interpreting a Plasma–Filtrate–Urine Table]
A typical GCE question gives a table like:

\begin{center}
\begin{tabularx}{\linewidth}{|l|X|X|X|}\hline
\rowcolor{popblueL} Substance & Blood plasma & Glomerular filtrate & Urine \\ \hline
Water & 90\% & 90\% & 95\% \\ \hline
Proteins & 7\% & 0\% & 0\% \\ \hline
Glucose & 0.1\% & 0.1\% & 0\% \\ \hline
Urea & 0.03\% & 0.03\% & 2\% \\ \hline
Na$^+$ & 0.3\% & 0.3\% & 0.6\% \\ \hline
\end{tabularx}
\end{center}

\textbf{Explanation:}
\begin{itemize}
\item Proteins: too large to filter $\to$ absent in filtrate and urine.
\item Glucose: freely filtered, but \emph{completely reabsorbed} in PCT $\to$ absent in urine (unless plasma glucose exceeds \emph{renal threshold} $\approx$ \SI{10}{mmol/L}).
\item Urea: filtered, partially reabsorbed passively (about 50\%), but concentrated because water is reabsorbed $\to$ higher \% in urine.
\item Na$^+$: filtered, mostly reabsorbed, but some excreted to regulate blood volume/pressure.
\end{itemize}
\end{exemplebox}

\courssub{The Counter-Current Multiplier and Medullary Gradient}

\begin{experience}[Understanding the Counter-Current Multiplier]
\textbf{Observation:} The medullary interstitial fluid becomes progressively more hypertonic from cortex (\SI{\sim 300}{mOsm/L}) to papilla (\SI{\sim 1200}{mOsm/L} in humans, up to \SI{5000}{mOsm/L} in desert rodents). This gradient is generated by the loop of Henlé and maintained by the vasa recta.

\textbf{Mechanism (simplified):}
\begin{enumerate}
\item Active transport of Na$^+$, K$^+$, 2Cl$^-$ out of the \textbf{thick ascending limb} (impermeable to water) lowers interstitial water potential.
\item Water leaves the \textbf{descending limb} (permeable to water) by osmosis into the interstitium.
\item The filtrate becomes concentrated at the bend of the loop.
\item As fluid ascends, solutes are pumped out but water cannot follow $\to$ fluid becomes dilute, interstitium becomes more concentrated.
\item This \emph{single effect} multiplied along the length of the loop creates the axial gradient.
\item \cle{Vasa recta} (counter-current exchanger): blood flows down into medulla (gains solutes, loses water) and up (loses solutes, gains water) $\to$ preserves the gradient while supplying nutrients/O$_2$.
\end{enumerate}
\textbf{Significance:} The gradient allows the collecting duct to reabsorb water under ADH control, producing urine hypertonic to plasma.
\end{experience}

\courssub{Hormonal Regulation of Water and Salt Balance}

\begin{aretenir}[Key Hormones Acting on the Nephron]
\begin{tabularx}{\linewidth}{|l|X|X|X|}\hline
\rowcolor{popblueL} Hormone & Source & Target & Main Action \\ \hline
ADH (vasopressin) & Posterior pituitary & Collecting duct (principal cells) & Inserts aquaporin-2 $\to$ $\uparrow$ water permeability $\to$ concentrated urine, $\downarrow$ volume \\ \hline
Aldosterone & Adrenal cortex (zona glomerulosa) & DCT, collecting duct (principal cells) & $\uparrow$ ENaC, Na$^+$/K$^+$-ATPase $\to$ $\uparrow$ Na$^+$ reabsorption, $\uparrow$ K$^+$/H$^+$ secretion \\ \hline
ANP & Heart atria & Collecting duct, afferent arteriole & $\downarrow$ Na$^+$ reabsorption, $\uparrow$ GFR $\to$ natriuresis, diuresis \\ \hline
Parathyroid hormone (PTH) & Parathyroid glands & DCT, PCT & $\uparrow$ Ca$^{2+}$ reabsorption, $\downarrow$ PO$_4^{3-}$ reabsorption \\ \hline
Angiotensin II & Systemic (RAAS) & Afferent/efferent arterioles, PCT & Vasoconstriction (efferent $>$ afferent) $\to$ maintains GFR; $\uparrow$ Na$^+$ reabsorption in PCT \\ \hline
\end{tabularx}
\end{aretenir}

\courssub{Adaptations of the Nephron}

\begin{aretenir}[Key Adaptations]
\begin{tabularx}{\linewidth}{|l|X|}\hline
\rowcolor{popblueL} Structure & Adaptation for Function \\ \hline
Glomerulus & Fenestrated endothelium, high pressure (afferent $>$ efferent diameter) $\to$ high filtration rate. \\ \hline
Basement membrane & Negative charge $\&$ pore size $\to$ size- and charge-selective barrier. \\ \hline
Podocytes & Foot processes with filtration slits + slit diaphragm $\to$ prevent protein loss. \\ \hline
PCT cells & Dense microvilli (brush border) $\uparrow$ surface area $\sim$40$\times$; numerous mitochondria $\to$ ATP for active transport; basal infoldings $\uparrow$ basolateral area. \\ \hline
Loop of Henlé & Long loops in desert mammals $\to$ steep medullary gradient $\to$ concentrated urine; short loops in freshwater mammals $\to$ dilute urine. \\ \hline
Collecting duct & Aquaporin-2 vesicles (ADH-regulated) $\to$ variable water permeability. \\ \hline
Peritubular capillaries / vasa recta & Dense network, low pressure, fenestrated $\to$ efficient uptake of reabsorbed substances without washing out medullary gradient. \\ \hline
Juxtaglomerular apparatus & Modified DCT (macula densa) + JG cells (modified afferent/efferent arterioles) $\to$ renin secretion, tubuloglomerular feedback. \\ \hline
\end{tabularx}
\end{aretenir}

\courssub{Micturition, Kidney Failure and Dialysis (GCE Extension)}

\begin{itemize}
\item \textbf{Micturition}: Bladder fills (\SI{150}{--}\SI{300}{mL} $\to$ urge) $\to$ stretch receptors $\to$ spinal reflex (parasympathetic pelvic nerves, S2--S4) $\to$ detrusor muscle contracts, internal urethral sphincter (smooth muscle, involuntary) relaxes. Voluntary control (somatic pudendal nerve) of external urethral sphincter (skeletal muscle) allows postponement until appropriate.
\item \textbf{Kidney failure}: Loss of filtration capacity (GFR $<$ \SI{15}{mL/min} = stage 5 CKD) $\to$ accumulation of urea, creatinine, K$^+$, H$^+$, water, phosphate $\to$ uraemia, metabolic acidosis, hyperkalaemia, pulmonary oedema, anaemia (low EPO), renal osteodystrophy.
\item \textbf{Dialysis} (artificial kidney):
  \begin{itemize}
  \item \cle{Haemodialysis}: Blood pumped through a dialyser (semi-permeable hollow fibres) against dialysate fluid (controlled electrolyte composition); diffusion removes wastes, excess ions; ultrafiltration (hydrostatic pressure gradient) removes water. Typically \SI{4}{h}, 3 times/week. Requires vascular access (AV fistula, graft, or catheter).
  \item \cle{Peritoneal dialysis}: Peritoneum as natural membrane; dialysate infused into peritoneal cavity (dwell \SI{4}{--}\SI{6}{h}), diffusion/osmosis occur, drain. Can be continuous ambulatory (CAPD) or automated (APD).
  \end{itemize}
  \emph{Kidney transplant} is the definitive treatment; lifelong immunosuppression required (tacrolimus, mycophenolate, prednisolone). Living donor or deceased donor.
\end{itemize}

\begin{methode}[Explaining Glucose in Urine of a Diabetic Person]
\begin{enumerate}
\item State that in diabetes mellitus, blood glucose concentration exceeds the \cle{renal threshold} (\SI{\sim 10}{mmol/L} or \SI{\sim 180}{mg/dL}).
\item Glucose is freely filtered at the glomerulus (concentration in filtrate = plasma concentration).
\item The PCT reabsorbs glucose via Na$^+$/glucose co-transporters (SGLT2 in early PCT, SGLT1 in late PCT). These carriers have a maximum transport rate (\cle{Tm}, transport maximum $\approx$ \SI{375}{mg/min} in adults).
\item When filtered load (GFR $\times$ plasma [glucose]) $>$ Tm, the carriers are saturated $\to$ glucose remains in the tubule $\to$ appears in urine (\cle{glycosuria}).
\item Consequence: glucose in tubule lowers water potential (osmotic effect) $\to$ less water reabsorbed $\to$ osmotic diuresis $\to$ polyuria $\to$ polydipsia (thirst).
\end{enumerate}
\end{methode}

\begin{exoresolu}[Comparing Plasma, Filtrate and Urine Concentrations]
\textbf{Statement:} The table below shows concentrations (in mmol/L) of three substances in blood plasma, glomerular filtrate, and urine of a healthy person.

\begin{center}
\begin{tabular}{|c|c|c|c|}\hline
\rowcolor{popblueL} Substance & Plasma & Filtrate & Urine \\ \hline
Urea & 5 & 5 & 300 \\ \hline
Glucose & 5 & 5 & 0 \\ \hline
Protein & 70 & 0 & 0 \\ \hline
\end{tabular}
\end{center}

Explain the values for each substance.

\textbf{Solution:}
\begin{itemize}
\item \textbf{Protein (70, 0, 0)}: Proteins are too large to pass the filtration barrier (basement membrane + podocytes + charge barrier). Hence they are absent from filtrate and urine. Presence of protein in urine (\cle{proteinuria}) indicates glomerular damage.
\item \textbf{Glucose (5, 5, 0)}: Glucose is small and freely filtered, so filtrate concentration equals plasma concentration. In the PCT, \emph{all} glucose is actively reabsorbed (via Na$^+$/glucose co-transport SGLT2/SGLT1) provided plasma glucose is below the renal threshold (\SI{\sim 10}{mmol/L}). Hence urine contains no glucose.
\item \textbf{Urea (5, 5, 300)}: Urea is small and freely filtered (filtrate = plasma). About 50\% of filtered urea is passively reabsorbed (mainly in collecting duct and thin ascending limb), but because \emph{large volumes of water} are reabsorbed (obligatory in PCT, facultative in collecting duct under ADH), the remaining urea becomes highly concentrated. The urine concentration (300 mmol/L) is 60$\times$ plasma, reflecting the kidney's concentrating ability. Urea also contributes to the medullary osmotic gradient (urea recycling).
\end{itemize}
\end{exoresolu}

\begin{exoresolu}[Calculating Filtration Fraction and Reabsorption Rate]
\textbf{Statement:} In a healthy adult, renal plasma flow (RPF) = \SI{600}{mL/min}, GFR = \SI{125}{mL/min}, plasma [urea] = \SI{5}{mmol/L}, urine flow rate = \SI{1}{mL/min}, urine [urea] = \SI{300}{mmol/L}. Calculate: (a) filtration fraction, (b) filtered load of urea, (c) urea excretion rate, (d) fraction of filtered urea excreted.

\textbf{Solution:}
\begin{align*}
\text{(a) Filtration fraction} &= \frac{\text{GFR}}{\text{RPF}} = \frac{125}{600} = 0.208 \approx 20.8\% \\[4pt]
\text{(b) Filtered load of urea} &= \text{GFR} \times \text{plasma [urea]} = 125\ \mathrm{mL/min} \times 5\ \mathrm{mmol/L} \\
&= 0.125\ \mathrm{L/min} \times 5\ \mathrm{mmol/L} = 0.625\ \mathrm{mmol/min} \\[4pt]
\text{(c) Urea excretion rate} &= \text{urine flow} \times \text{urine [urea]} = 1\ \mathrm{mL/min} \times 300\ \mathrm{mmol/L} \\
&= 0.001\ \mathrm{L/min} \times 300\ \mathrm{mmol/L} = 0.3\ \mathrm{mmol/min} \\[4pt]
\text{(d) Fraction excreted} &= \frac{\text{excretion rate}}{\text{filtered load}} = \frac{0.3}{0.625} = 0.48 = 48\%
\end{align*}
\textit{Interpretation:} About 52\% of filtered urea is reabsorbed; the rest is excreted. This illustrates that urea handling is not all-or-none like glucose.
\end{exoresolu}

\begin{exercice}[GCE-Style Questions]
\begin{enumerate}
\item Draw a labelled diagram of a mammalian nephron with its associated blood vessels. Indicate the direction of blood flow and the sites of ultrafiltration, selective reabsorption, and tubular secretion.
\item A student writes: ``The kidney filters urea out of the blood and excretes it.'' Criticise this statement and rewrite it correctly.
\item Explain why the loop of Henlé is longer in desert-dwelling mammals (e.g., camel) than in freshwater mammals (e.g., beaver).
\item In a urinalysis, glucose is detected. Give two possible physiological explanations.
\item Describe the role of ADH in the production of concentrated urine.
\item A patient has a GFR of \SI{120}{mL/min} and plasma creatinine of \SI{100}{\mu mol/L}. Urine flow is \SI{1.5}{mL/min} and urine creatinine is \SI{12}{mmol/L}. Calculate the creatinine clearance. What does this value indicate about GFR? (Creatinine is freely filtered and not reabsorbed, but \SI{\sim 10}{\%} is secreted).
\item Explain how the vasa recta helps maintain the medullary osmotic gradient.
\end{enumerate}
\end{exercice}

\begin{corrige}[Exercise 5]
\begin{enumerate}
\item Diagram should show: afferent arteriole $\to$ glomerulus $\to$ efferent arteriole $\to$ peritubular capillaries/vasa recta $\to$ renal vein; Bowman's capsule, PCT, loop of Henlé, DCT, collecting duct. Annotate: ultrafiltration at glomerulus/Bowman's capsule; selective reabsorption mainly PCT (also loop, DCT, collecting duct); secretion at DCT/collecting duct.
\item Criticism: ``filters urea out of the kidney'' is wrong — urea is filtered out of the \emph{blood} (in the glomerulus) into the nephron. Corrected: ``The kidney filters urea out of the blood in the glomerulus; most remains in the tubule and is excreted in urine.''
\item Longer loop of Henlé $\to$ greater counter-current multiplication $\to$ steeper medullary osmotic gradient $\to$ more water can be reabsorbed from collecting duct under ADH $\to$ more concentrated urine $\to$ water conservation in arid environments.
\item (1) Diabetes mellitus: blood glucose $>$ renal threshold $\to$ saturation of PCT glucose transporters $\to$ glycosuria. (2) Renal glycosuria: genetic defect in PCT glucose transporters (SGLT2) $\to$ glucose not reabsorbed despite normal blood glucose.
\item ADH (vasopressin) released from posterior pituitary when blood osmolarity rises (detected by hypothalamic osmoreceptors) or blood volume falls (baroreceptors). Acts on collecting duct principal cells via V2 receptors $\to$ cAMP $\to$ insertion of aquaporin-2 water channels into apical membrane $\to$ increased water permeability $\to$ water moves out of duct into hypertonic medulla by osmosis $\to$ concentrated urine, reduced volume.
\item Creatinine clearance $C_{\text{Cr}} = \frac{U_{\text{Cr}} \times V}{P_{\text{Cr}}} = \frac{12\ \mathrm{mmol/L} \times 1.5\ \mathrm{mL/min}}{0.1\ \mathrm{mmol/L}} = \frac{18}{0.1} = 180\ \mathrm{mL/min}$. Since creatinine is secreted (~10\%), clearance overestimates true GFR by ~10--15\%. True GFR $\approx$ \SI{160}{mL/min} (slightly high, possibly due to hyperfiltration or measurement error).
\item Vasa recta are counter-current exchangers: blood descends into medulla (gains solutes by diffusion, loses water by osmosis) and ascends (loses solutes, gains water). This U-shaped flow minimises net solute/water removal from medulla, preserving the gradient generated by loops of Henlé.
\end{enumerate}
\end{corrige}

\begin{savaistu}[Cameroon Context: Kidney Disease and Dialysis]
In Cameroon, chronic kidney disease (CKD) is a growing public health concern, often linked to hypertension, diabetes, and glomerulonephritis (sometimes post-infectious, e.g., after streptococcal infections or malaria). Access to dialysis is limited — major centres are in Yaoundé (Hôpital Central, Hôpital Général) and Douala (Laquintinie, Hôpital Général). A session of haemodialysis costs \SI{30\,000}{--}\SI{50\,000}{FCFA} (often not fully covered by insurance), making it unaffordable for many. Kidney transplantation is performed occasionally (e.g., at Hôpital Général de Yaoundé) but requires lifelong immunosuppression. Prevention — control of blood pressure, blood sugar, avoidance of nephrotoxic herbal concoctions — remains the most effective strategy. The \emph{National Programme for the Fight against Non-Communicable Diseases} includes CKD screening in high-risk groups.
\end{savaistu}

\begin{popfigure}
\begin{tikzpicture}
\begin{axis}[
    ybar,
    bar width=0.35cm,
    width=0.9\linewidth,
    height=8cm,
    enlarge x limits=0.25,
    ylabel={Concentration (arbitrary units)},
    symbolic x coords={Urea, Glucose, Protein},
    xtick=data,
    nodes near coords,
    nodes near coords align={vertical},
    legend style={at={(0.5,-0.15)}, anchor=north, legend columns=3},
    ymin=0,
    ymax=350,
    tick label style={font=\small},
    label style={font=\small},
    legend style={font=\small},
]
\addplot[fill=popblue!70] coordinates {(Urea,5) (Glucose,5) (Protein,70)};
\addplot[fill=popgreen!70] coordinates {(Urea,5) (Glucose,5) (Protein,0)};
\addplot[fill=poporange!70] coordinates {(Urea,300) (Glucose,0) (Protein,0)};
\legend{Blood plasma, Glomerular filtrate, Urine}
\end{axis}
\end{tikzpicture}
\legende{Comparison of urea, glucose and protein concentrations in blood plasma, glomerular filtrate and urine (typical values for a healthy adult).}
\end{popfigure}

\courssub{Summary of Key Points}

\begin{aretenir}[Section 5 at a Glance]
\begin{itemize}
\item The urinary system: 2 kidneys, 2 ureters, bladder, urethra. Renal artery/vein connect to aorta/vena cava.
\item Kidney: cortex (glomeruli, tubules), medulla (pyramids, loops, collecting ducts), pelvis.
\item Nephron = functional unit: Bowman's capsule, PCT, loop of Henlé, DCT, collecting duct + capillary network. Cortical vs juxtamedullary nephrons.
\item \textbf{Ultrafiltration}: high pressure in glomerulus (afferent $>$ efferent) $\to$ filtrate = plasma minus proteins. Barrier: fenestrated endothelium, basement membrane, podocytes with slit diaphragm. GFR $\sim$ \SI{125}{mL/min}.
\item \textbf{Selective reabsorption}: PCT reabsorbs all glucose/amino acids, \SI{\sim 65}{--}\SI{70}{\%} water/ions (active Na$^+$ transport + osmosis). Loop creates medullary gradient via counter-current multiplier. DCT/collecting duct fine-tune under ADH/aldosterone/ANP/PTH.
\item \textbf{Tubular secretion}: H$^+$, K$^+$, NH$_4^+$, drugs, toxins $\to$ tubular lumen (DCT, collecting duct).
\item Urine: 95\% water, urea, salts, creatinine; \textbf{no glucose, no protein} normally. pH 4.5--8.0.
\item Adaptations: microvilli, mitochondria, long loops (desert mammals), aquaporins, fenestrated capillaries, vasa recta counter-current exchange.
\item Diabetes $\to$ glycosuria (filtered load $>$ Tm). Kidney failure $\to$ dialysis/transplant.
\item Hormonal control: ADH (water), aldosterone (Na$^+$/K$^+$), ANP (natriuresis), RAAS (GFR/Na$^+$), PTH (Ca$^{2+}$).
\end{itemize}
\end{aretenir}

\coursec{Section 6: Osmoregulation -- Kidneys and Animal Adaptations}

Having explored the structure of the renal system and the detailed mechanisms of urine formation in Section 5, we now turn to the physiological \emph{control} of that process. The nephron is a powerful filtration and reabsorption machine, but it does not operate at a fixed rate. The volume and concentration of the final urine are continuously adjusted to match the body's immediate needs. This dynamic regulation of the internal water and solute balance is the essence of \cle{osmoregulation}.

\courssub{What Is Osmoregulation?}

\begin{definition}[Osmoregulation]
Osmoregulation is the process by which organisms control the \cle{water potential} ($\Psi_w$) of their body fluids (blood, lymph, interstitial fluid, cytoplasm) by regulating the gain and loss of water and solutes (mainly salts), thereby maintaining a stable internal environment suitable for cellular metabolism.
\end{definition}

\begin{propriete}[Water Potential and Osmosis in the Body]
Water moves by osmosis from a region of higher (less negative) water potential to a region of lower (more negative) water potential across a selectively permeable membrane.
\begin{itemize}
    \item Pure water: $\Psi_w = 0$ kPa (highest potential).
    \item Adding solutes lowers $\Psi_w$ (makes it more negative).
    \item Blood plasma $\Psi_w$ is normally maintained around \textbf{-300 kPa to -350 kPa} (approx. 300 mOsm/L) in mammals.
\end{itemize}
If blood $\Psi_w$ falls (becomes more negative $\rightarrow$ dehydration), water leaves cells $\rightarrow$ crenation. If blood $\Psi_w$ rises (becomes less negative $\rightarrow$ overhydration), water enters cells $\rightarrow$ lysis. Osmoregulation prevents both extremes.
\end{propriete}

\begin{aretenir}[Key Variables in Osmoregulation]
\begin{itemize}
    \item \textbf{Input:} Water and salts from food, drink, metabolic water (respiration), and (in some animals) environment.
    \item \textbf{Output:} Urine, faeces, sweat, exhaled air, (eggs/milk).
    \item \textbf{The Kidney:} The main effector organ in vertebrates; varies urine volume and concentration to balance inputs and outputs.
\end{itemize}
\end{aretenir}

\courssub{The Kidney as an Osmoregulatory Organ}

The mammalian kidney can produce urine ranging from very dilute (\textbf{50--100 mOsm/L}, $\Psi_w \approx 0$ kPa) to very concentrated (\textbf{1200--1400 mOsm/L}, $\Psi_w \approx -3500$ kPa in humans; up to -5000 kPa in rodents). This flexibility depends on two key factors:
\begin{enumerate}
    \item \textbf{The Medullary Osmotic Gradient:} Established and maintained by the \cle{Loop of Henlé} acting as a \cle{counter-current multiplier}.
    \item \textbf{Hormonal Control of Permeability:} \cle{Antidiuretic Hormone (ADH)} (also called \cle{vasopressin}) regulates the permeability of the \cle{distal convoluted tubule (DCT)} and \cle{collecting duct} to water.
\end{enumerate}

\begin{definition}[Antidiuretic Hormone (ADH / Vasopressin)]
ADH is a peptide hormone \textbf{synthesised in the neurosecretory cells of the \cle{hypothalamus}}, transported down axons to the \cle{posterior pituitary gland} where it is \textbf{stored and released} into the bloodstream. Its target cells are the principal cells of the \textbf{distal convoluted tubule and collecting duct} in the kidney nephron.
\end{definition}

\begin{attention}[Common GCE Mistake: Site of ADH Production]
\textbf{WRONG:} ``ADH is produced by the pituitary gland.'' \\
\textbf{CORRECT:} ADH is \textbf{produced (synthesised) by the hypothalamus} and only \textbf{stored and released} by the posterior pituitary. The posterior pituitary is neural tissue (axon terminals), not a true glandular epithelium. This distinction is frequently tested in multiple-choice and short-answer questions.
\end{attention}

\begin{propriete}[Mechanism of ADH Action]
\begin{enumerate}
    \item ADH binds to \cle{V2 receptors} on the basolateral membrane of DCT/collecting duct cells.
    \item This triggers a cAMP second messenger cascade.
    \item Vesicles containing \cle{aquaporin-2 (AQP2) water channels} fuse with the \textbf{apical (luminal) membrane}.
    \item Water permeability increases $\rightarrow$ water moves out of the filtrate into the hypertonic medullary interstitium by osmosis.
    \item When ADH levels fall, aquaporins are removed by endocytosis $\rightarrow$ membrane becomes impermeable to water $\rightarrow$ dilute urine is produced.
\end{enumerate}
\end{propriete}

\begin{savaistu}[The Counter-Current Multiplier (Outline for GCE)]
The Loop of Henlé creates a corticomedullary osmotic gradient (cortex $\approx$ 300 mOsm/L $\rightarrow$ inner medulla $\approx$ 1200+ mOsm/L).
\begin{itemize}
    \item \textbf{Descending limb:} Highly permeable to water (AQP1), impermeable to salts. Water leaves by osmosis $\rightarrow$ filtrate becomes concentrated.
    \item \textbf{Ascending limb:} Impermeable to water. \textbf{Thick segment:} Active transport of Na$^+$, K$^+$, 2Cl$^-$ out (NKCC2 cotransporter). Dilutes filtrate, adds solutes to interstitium.
    \item \textbf{Vasa recta:} Counter-current exchanger blood vessels that preserve the gradient by slow flow and U-shape.
\end{itemize}
\textit{Examinable:} You must be able to describe the permeability differences and the role of active transport in the thick ascending limb. You do \emph{not} need to derive the mathematics of the multiplier.
\end{savaistu}

\courssub{The ADH Negative Feedback Loop}

This is the core homeostatic mechanism for blood water potential. It operates in two opposing directions.

\begin{popfigure}
\centering
\begin{tikzpicture}[
    node distance=1.8cm and 2.2cm,
    box/.style={draw=popdark, thick, rounded corners, fill=popblueL, text width=3.2cm, align=center, minimum height=1.2cm},
    arrow/.style={->, thick, popdark},
    plus/.style={font=\bfseries\large, poporange},
    minus/.style={font=\bfseries\large, popgreen}
]
% Dehydration pathway (left)
\node[box, align=center] (stim1){Dehydration\\Blood $\Psi_w$ falls\\(more negative)};
\node[box, below=of stim1, align=center] (det1){Hypothalamic\\osmoreceptors\\stimulated};
\node[box, below=of det1, align=center] (eff1){Posterior pituitary\\releases \textbf{more ADH}};
\node[box, below=of eff1, align=center] (resp1){DCT \& Collecting duct\\become \textbf{permeable} to water\\(AQP2 insertion)};
\node[box, below=of resp1, align=center] (out1){Water reabsorbed\\by osmosis into blood\\(medullary gradient)};
\node[box, below=of out1, align=center] (res1){Small volume of\\ \textbf{concentrated urine}\\(high $\Psi_s$, low $\Psi_w$)};
\node[box, below=of res1, align=center] (fb1){Blood $\Psi_w$ rises\\(returns to normal)};

% Overhydration pathway (right)
\node[box, right=of stim1, align=center] (stim2){Overhydration\\Blood $\Psi_w$ rises\\(less negative)};
\node[box, below=of stim2, align=center] (det2){Hypothalamic\\osmoreceptors\\inhibited};
\node[box, below=of det2, align=center] (eff2){Posterior pituitary\\releases \textbf{less ADH}};
\node[box, below=of eff2, align=center] (resp2){DCT \& Collecting duct\\become \textbf{impermeable} to water\\(AQP2 removal)};
\node[box, below=of resp2, align=center] (out2){Water \textbf{not} reabsorbed\\\textit{remains in tubule}};
\node[box, below=of out2, align=center] (res2){Large volume of\\ \textbf{dilute urine}\\(low $\Psi_s$, high $\Psi_w$)};
\node[box, below=of res2, align=center] (fb2){Blood $\Psi_w$ falls\\(returns to normal)};

% Arrows
\draw[arrow] (stim1) -- (det1);
\draw[arrow] (det1) -- (eff1);
\draw[arrow] (eff1) -- (resp1);
\draw[arrow] (resp1) -- (out1);
\draw[arrow] (out1) -- (res1);
\draw[arrow] (res1) -- (fb1);
\draw[arrow, bend left=40] (fb1) to node[plus, pos=0.5, xshift=-1.5cm] {+} (stim1);

\draw[arrow] (stim2) -- (det2);
\draw[arrow] (det2) -- (eff2);
\draw[arrow] (eff2) -- (resp2);
\draw[arrow] (resp2) -- (out2);
\draw[arrow] (out2) -- (res2);
\draw[arrow] (res2) -- (fb2);
\draw[arrow, bend right=40] (fb2) to node[minus, pos=0.5, xshift=1.5cm] {-} (stim2);

% Cross labels
\node[align=center, font=\small\bfseries, poppurple] at ($(stim1)!0.5!(stim2) + (0,0.5)$) {STIMULUS};
\node[align=center, font=\small\bfseries, poppurple] at ($(det1)!0.5!(det2) + (0,-0.3)$) {RECEPTOR};
\node[align=center, font=\small\bfseries, poppurple] at ($(eff1)!0.5!(eff2) + (0,-0.3)$) {EFFECTOR};
\node[align=center, font=\small\bfseries, poppurple] at ($(resp1)!0.5!(resp2) + (0,-0.3)$) {RESPONSE};
\node[align=center, font=\small\bfseries, poppurple] at ($(out1)!0.5!(out2) + (0,-0.3)$) {OUTCOME};
\node[align=center, font=\small\bfseries, poppurple] at ($(fb1)!0.5!(fb2) + (0,-0.3)$) {FEEDBACK};
\end{tikzpicture}
\legende{ADH negative feedback loops for dehydration (left) and overhydration (right). Note the opposing responses restoring blood water potential to the set point.}
\end{popfigure}

\begin{methode}[Predicting Urine Volume and Concentration from Blood Water Potential]
\textbf{Step 1:} Identify the stimulus: Is blood $\Psi_w$ \textbf{low} (dehydration) or \textbf{high} (overhydration)?
\textbf{Step 2:} Determine hypothalamus response: Low $\Psi_w \rightarrow$ \textbf{stimulate} osmoreceptors $\rightarrow$ more ADH. High $\Psi_w \rightarrow$ \textbf{inhibit} osmoreceptors $\rightarrow$ less ADH.
\textbf{Step 3:} Determine kidney tubule response: More ADH $\rightarrow$ \textbf{insert AQP2} $\rightarrow$ high permeability. Less ADH $\rightarrow$ \textbf{remove AQP2} $\rightarrow$ low permeability.
\textbf{Step 4:} Determine urine output: High permeability + medullary gradient $\rightarrow$ water leaves tubule $\rightarrow$ \textbf{small volume, concentrated urine}. Low permeability $\rightarrow$ water stays in tubule $\rightarrow$ \textbf{large volume, dilute urine}.
\textbf{Step 5:} State the feedback: Urine output corrects the blood $\Psi_w$ back toward the set point.
\end{methode}

\begin{exoresolu}[Interpreting a Clinical Scenario]
\textbf{Statement:} A patient is admitted after 48 hours lost in the desert without water. Blood tests show plasma osmolarity = 320 mOsm/L (normal $\approx$ 290 mOsm/L). Predict and explain the state of their ADH secretion, kidney tubule permeability, and urine characteristics.

\tcblower
\textbf{Detailed Solution:}
\begin{enumerate}
    \item \textbf{Blood Water Potential:} Plasma osmolarity is \textbf{elevated} (320 > 290 mOsm/L). Therefore, blood $\Psi_w$ is \textbf{lower (more negative)} than normal $\rightarrow$ \textbf{dehydration}.
    \item \textbf{ADH Secretion:} Hypothalamic osmoreceptors detect low $\Psi_w$ $\rightarrow$ \textbf{stimulated} $\rightarrow$ posterior pituitary releases \textbf{high levels of ADH} (maximal secretion).
    \item \textbf{Tubule Permeability:} High ADH $\rightarrow$ abundant AQP2 channels inserted into apical membranes of DCT and collecting duct $\rightarrow$ \textbf{high water permeability}.
    \item \textbf{Urine Characteristics:}
        \begin{itemize}
            \item \textbf{Volume: Very low} (oliguria, $< 0.5$ mL/min) to conserve water.
            \item \textbf{Concentration: Very high} (hypertonic, up to 1200--1400 mOsm/L), dark colour.
            \item Water moves osmotically from filtrate $\rightarrow$ hypertonic medulla $\rightarrow$ blood.
        \end{itemize}
    \item \textbf{Feedback:} Water reabsorption raises blood volume and $\Psi_w$ toward normal.
\end{enumerate}
\textbf{GCE Answer Format:} ``Blood water potential is low. Osmoreceptors in hypothalamus stimulated. More ADH released from posterior pituitary. Collecting duct becomes permeable to water. Water reabsorbed by osmosis into blood. Small volume of concentrated urine produced. Blood water potential rises back to normal.''
\end{exoresolu}

\courssub{Osmoregulation in Different Environments}

Animals face vastly different osmotic challenges depending on their habitat. The kidney (or equivalent organ) is adapted accordingly.

\begin{popfigure}
\centering
\begin{tikzpicture}[
    node distance=1.2cm and 1.5cm,
    envbox/.style={draw=popdark, thick, rounded corners, fill=popgreenL, text width=3.5cm, align=center, minimum height=1.5cm},
    fishbox/.style={draw=popdark, thick, rounded corners, fill=popblueL, text width=3.8cm, align=center, minimum height=1.5cm},
    arrow/.style={->, thick, popdark},
    warrow/.style={->, thick, popblue, bend left=20},
    sarrow/.style={->, thick, poporange, bend right=20}
]
% Freshwater Fish (Left)
\node[envbox, align=center] (fwtitle){\textbf{Freshwater Bony Fish}\\(e.g. Tilapia, Trout)};
\node[fishbox, below=of fwtitle, align=center] (fwprob){\textbf{Problem:} Water enters body by osmosis\\ (environment $\Psi_w = 0$ > body $\Psi_w$).\\Salts diffuse out.};
\node[fishbox, below=of fwprob, align=center] (fwkid){\textbf{Kidney:} Large glomeruli, high filtration rate.\\Produces \textbf{large volume of dilute urine}.};
\node[fishbox, below=of fwkid, align=center] (fwgill){\textbf{Gills:} Active uptake of salts (Na$^+$, Cl$^-$)\\via mitochondria-rich cells (chloride cells).\\No drinking.};
\node[fishbox, below=of fwgill] (fwbal) {\textbf{Net:} Excrete water, conserve salts.};

% Marine Bony Fish (Right)
\node[envbox, right=of fwtitle, align=center] (mwtitle){\textbf{Marine Bony Fish}\\(e.g. Mackerel, Cod)};
\node[fishbox, below=of mwtitle, align=center] (mwprob){\textbf{Problem:} Water leaves body by osmosis\\ (environment $\Psi_w$ very negative).\\Salts diffuse in.};
\node[fishbox, below=of mwprob, align=center] (mwkid){\textbf{Kidney:} Small/absent glomeruli (aglomerular),\\low filtration rate.\\Produces \textbf{small volume of isotonic urine}.};
\node[fishbox, below=of mwkid, align=center] (mwgill){\textbf{Gills:} Active \textbf{excretion} of excess salts\\ (Na$^+$, Cl$^-$) via chloride cells.\\ \textbf{Drinks seawater} to gain water.};
\node[fishbox, below=of mwgill] (mwbal) {\textbf{Net:} Conserve water, excrete salts.};

% Arrows for water/salt movement
\draw[warrow] (fwprob.west) -- ++(-1,0) node[left, popblue] {Water IN};
\draw[sarrow] (fwprob.east) -- ++(1,0) node[right, poporange] {Salts OUT};
\draw[warrow] (mwprob.west) -- ++(-1,0) node[left, popblue] {Water OUT};
\draw[sarrow] (mwprob.east) -- ++(1,0) node[right, poporange] {Salts IN};

% Drinking arrow
\draw[arrow, popblue] (mwgill.south west) -- ++(-0.5,-0.5) node[below left, popblue, font=\small] {Drinks seawater};

% Urine arrows
\draw[arrow, popmuted] (fwkid.south) -- ++(0,-0.3) node[below, font=\small] {Large volume, dilute urine};
\draw[arrow, popmuted] (mwkid.south) -- ++(0,-0.3) node[below, font=\small] {Small volume, isotonic urine};
\end{tikzpicture}
\legende{Comparative osmoregulation in freshwater vs. marine bony fish. Note the opposite directions of water and salt fluxes and the contrasting roles of the kidney and gills. Tilapia, common in Cameroonian rivers and fish farms, is a classic freshwater example.}
\end{popfigure}

\begin{attention}[GCE Exam Tip: Arrows in Fish Osmoregulation Diagrams]
In GCE practical or theory questions asking you to draw or annotate diagrams of fish osmoregulation:
\begin{itemize}
    \item Use \textbf{blue arrows} for \textbf{water movement} (always follow water potential gradient: high $\Psi_w \rightarrow$ low $\Psi_w$).
    \item Use \textbf{red/orange arrows} for \textbf{active salt transport} (against concentration gradient, requires ATP).
    \item Label the gills clearly: ``Active salt uptake'' (freshwater) vs ``Active salt excretion'' (marine).
    \item Label the kidney: ``Dilute urine'' vs ``Concentrated/Isotonic urine''.
\end{itemize}
\end{attention}

\textbf{Freshwater Animals (General)}\par
\begin{itemize}
    \item \textbf{Protozoa (Amoeba, Paramecium):} Contractile vacuole collects excess water (from osmosis) and expels it. Energy (ATP) required.
    \item \textbf{Freshwater Invertebrates:} Malpighian tubules (insects) or nephridia produce dilute urine; active salt uptake across body surface/gills.
    \item \textbf{Amphibians:} Permeable skin gains water; produce large volumes of dilute urine; active salt transport across skin.
\end{itemize}

\textbf{Marine Animals (General)}\par
\begin{itemize}
    \item \textbf{Marine Invertebrates (many):} \textbf{Osmoconformers} -- body fluids iso-osmotic with seawater (no energy spent on osmoregulation, but limited habitat range).
    \item \textbf{Marine Bony Fish (Teleosts):} As described above -- drink seawater, excrete salts at gills, small urine volume.
    \item \textbf{Marine Elasmobranchs (Sharks, Rays):} Retain urea and TMAO in blood to raise $\Psi_s$ (lower $\Psi_w$) to match or exceed seawater $\rightarrow$ slight water gain $\rightarrow$ excreted by rectal gland (salt) and kidney.
\end{itemize}

\textbf{Terrestrial Animals: The Challenge of Desiccation}\par
Air has very low water potential. Water loss occurs via respiration, excretion, cutaneous evaporation. Adaptations are behavioural and physiological.

\begin{experience}[Case Study: The Kangaroo Rat (\textit{Dipodomys}) -- Master of Water Conservation]
The kangaroo rat (North American desert rodent) survives \textbf{without drinking water}. Similar principles apply to Sahelian rodents in Cameroon's Far North region (e.g., gerbils, jerboas).
\begin{center}
\begin{tabularx}{\linewidth}{|l|X|}\hline
\textbf{Adaptation} & \textbf{Function} \\\hline
\textbf{Long Loop of Henlé} (relative medullary thickness $\approx$ 5--6$\times$ cortex) & Creates extremely steep medullary gradient $\rightarrow$ urine up to \textbf{5000--6000 mOsm/L} (5--6$\times$ plasma). \\\hline
\textbf{Nasal counter-current heat exchange} & Cools expired air $\rightarrow$ water condenses in nasal passages $\rightarrow$ reabsorbed. \\\hline
\textbf{Nocturnal / Burrowing} & Avoids daytime heat $\rightarrow$ reduces evaporative cooling needs. \\\hline
\textbf{Metabolic water} & Oxidation of dry seeds (carbohydrates/fats) yields $\approx$ 1.1 g water/g fat. Meets all water needs. \\\hline
\textbf{Dry faeces} & Rectal water reabsorption highly efficient. \\\hline
\end{tabularx}
\end{center}
\textbf{Key GCE Point:} The length of the Loop of Henlé correlates with the maximum urine concentration ability. Desert mammals have long loops; beavers (aquatic) have short loops.
\end{experience}

\begin{savaistu}[The Camel -- Not Just a ``Water Tank'']
Contrary to popular belief, the camel's hump stores \textbf{fat}, not water. Its osmoregulatory adaptations include:
\begin{itemize}
    \item Tolerates body temperature fluctuation (34--41$^\circ$C) $\rightarrow$ reduces sweating.
    \item Can lose up to \textbf{30\% body water} (vs 12\% fatal for humans) without circulatory collapse.
    \item Oval red blood cells (non-lyse on rapid rehydration).
    \item Very long Loops of Henlé $\rightarrow$ highly concentrated urine (syrup-like).
    \item Dry faeces (can be used as fuel immediately).
\end{itemize}
Camels are well adapted to the Sahelian zones of northern Cameroon.
\end{savaistu}

\courssub{Osmoconformers vs. Osmoregulators}

\begin{definition}[Osmoconformers and Osmoregulators]
\begin{itemize}
    \item \textbf{Osmoconformers:} Maintain body fluid osmolarity \textbf{equal to} the external medium (iso-osmotic). No active regulation of water potential. Most marine invertebrates (jellyfish, starfish, mussels). \textit{Advantage:} Low energy cost. \textit{Disadvantage:} Cannot tolerate salinity changes (stenohaline).
    \item \textbf{Osmoregulators:} Maintain body fluid osmolarity \textbf{different from} (usually lower than) the external medium. Active regulation via excretory organs. Most vertebrates, freshwater/marine fish, terrestrial animals. \textit{Advantage:} Wider habitat range (euryhaline potential). \textit{Disadvantage:} High energy cost (up to 5--10\% basal metabolic rate).
\end{itemize}
\end{definition}

\begin{aretenir}[Summary Table: Osmoregulatory Strategies]
\begin{center}
\begin{tabularx}{\linewidth}{|l|X|X|X|}\hline
\textbf{Group} & \textbf{Environment} & \textbf{Strategy} & \textbf{Key Structures} \\\hline
Marine Invertebrates & Stable seawater & Osmoconformers & Body surface \\\hline
Freshwater Protozoa & Hypotonic & Osmoregulators & Contractile vacuole \\\hline
Freshwater Fish & Hypotonic & Osmoregulators & Kidney (dilute urine), Gills (salt uptake) \\\hline
Marine Bony Fish & Hypertonic & Osmoregulators & Kidney (little urine), Gills (salt excretion), Drink seawater \\\hline
Marine Cartilaginous Fish & Hypertonic & Osmoregulators (urea retention) & Rectal gland, Kidney \\\hline
Terrestrial Vertebrates & Air (desiccating) & Osmoregulators & Kidney (concentrated urine), Skin, Lungs, Behaviour \\\hline
\end{tabularx}
\end{center}
\end{aretenir}

\courssub{Worked Example: Comparing Kidney Function}

\begin{exoresolu}[Data Analysis: Urine Concentration in Different Mammals]
\textbf{Statement:} The table shows maximum urine osmolarity and relative medullary thickness (medulla:cortex ratio) for four mammals.

\begin{center}
\begin{tabular}{|c|c|c|}\hline
\textbf{Species} & \textbf{Max Urine Osmolarity (mOsm/L)} & \textbf{Medulla:Cortex Ratio} \\\hline
Beaver & 600 & 1.2 \\\hline
Human & 1200 & 3.0 \\\hline
Rat & 2500 & 5.5 \\\hline
Kangaroo Rat & 5500 & 6.0 \\\hline
\end{tabular}
\end{center}

\textbf{Questions:}
\begin{enumerate}
    \item Describe the relationship between medullary thickness and maximum urine concentration.
    \item Explain the structural basis for this relationship.
    \item Predict the likely habitat of the beaver and the kangaroo rat based on this data.
\end{enumerate}

\tcblower
\textbf{Detailed Solution:}
\begin{enumerate}
    \item \textbf{Relationship:} There is a \textbf{positive correlation}. As the medulla:cortex ratio increases (1.2 $\rightarrow$ 6.0), the maximum urine osmolarity increases (600 $\rightarrow$ 5500 mOsm/L). The kangaroo rat produces urine $\approx$ 9$\times$ more concentrated than the beaver.
    \item \textbf{Structural Basis:} A thicker medulla houses \textbf{longer Loops of Henlé}. The counter-current multiplier operates over a longer distance, allowing a steeper corticomedullary osmotic gradient to be established. A steeper gradient means a lower water potential in the inner medulla, creating a stronger osmotic driving force for water reabsorption from the collecting duct under ADH influence.
    \item \textbf{Habitat Prediction:}
        \begin{itemize}
            \item \textbf{Beaver:} Low concentrating ability (600 mOsm/L), short loops $\rightarrow$ \textbf{aquatic/semi-aquatic} (freshwater). Water is abundant; no need to conserve.
            \item \textbf{Kangaroo Rat:} Extreme concentrating ability (5500 mOsm/L), very long loops $\rightarrow$ \textbf{arid desert}. Must survive without free water; maximal water conservation essential.
        \end{itemize}
\end{enumerate}
\textbf{Examiner's Comment:} This is a classic GCE data interpretation question. Always link \textbf{Loop of Henlé length} $\rightarrow$ \textbf{Medullary gradient steepness} $\rightarrow$ \textbf{Max urine concentration} $\rightarrow$ \textbf{Habitat water availability}.
\end{exoresolu}

\courssub{Integration: The Big Picture}

\begin{aretenir}[Osmoregulation: Integrated Summary]
\begin{enumerate}
    \item \textbf{Universal Principle:} Maintain blood/tissue $\Psi_w$ within narrow limits for enzyme function, cell volume, blood pressure.
    \item \textbf{Mammalian Kidney:} The effector organ. Flexibility comes from \textbf{Loop of Henlé (gradient builder)} + \textbf{ADH (permeability switch)}.
    \item \textbf{ADH Feedback:} Hypothalamus (sensor) $\rightarrow$ Posterior Pituitary (release) $\rightarrow$ Collecting Duct (effector) $\rightarrow$ Urine concentration adjusted $\rightarrow$ Blood $\Psi_w$ corrected.
    \item \textbf{Environmental Adaptation:} Structure matches function. Long loops $\leftrightarrow$ desert. Large glomeruli/dilute urine $\leftrightarrow$ freshwater. Drinking + salt glands $\leftrightarrow$ marine.
    \item \textbf{Energy Cost:} Osmoregulation is expensive (active transport). Osmoconformers avoid this cost but lose ecological flexibility.
\end{enumerate}
\end{aretenir}

\begin{exercice}[GCE-Style Practice Questions]
\begin{enumerate}
    \item \textbf{Define} the term \emph{osmoregulation}. State the organ responsible in mammals. [2 marks]
    \item \textbf{Distinguish} between the site of synthesis and the site of release of ADH. [2 marks]
    \item A person drinks 1 litre of pure water rapidly. Describe the sequence of events in the ADH negative feedback loop that restores blood water potential. [5 marks]
    \item \textbf{Compare} the osmoregulatory challenges and kidney/gill adaptations of a freshwater bony fish and a marine bony fish. [6 marks]
    \item \textbf{Explain} why the kangaroo rat can survive without drinking water, referring to its kidney structure and metabolic physiology. [4 marks]
    \item \textbf{Distinguish} between an osmoconformer and an osmoregulator. Give one example of each. [3 marks]
    \item The graph below shows the relationship between plasma ADH concentration and urine osmolarity in a human. (Imagine a sigmoidal curve: low ADH $\rightarrow$ 50 mOsm/L; high ADH $\rightarrow$ 1200 mOsm/L).
    \begin{enumerate}
        \item Describe the shape of the curve. [1 mark]
        \item Explain the physiological significance of the plateau at high ADH concentrations. [2 marks]
    \end{enumerate}
\end{enumerate}
\end{exercice}

\begin{corrige}[Exercise 6]
\begin{enumerate}
    \item Osmoregulation: Control of water potential of body fluids by regulating water/solute balance. Organ: Kidney.
    \item Synthesis: Hypothalamus (neurosecretory cells). Release: Posterior pituitary gland.
    \item Sequence: Water absorbed $\rightarrow$ blood $\Psi_w$ rises (less negative) $\rightarrow$ hypothalamic osmoreceptors inhibited $\rightarrow$ less ADH released $\rightarrow$ collecting duct impermeable (AQP2 removed) $\rightarrow$ large volume dilute urine produced $\rightarrow$ blood $\Psi_w$ falls back to normal.
    \item Freshwater: Water in, salts out. Large glomeruli, dilute urine, gills uptake salts. Marine: Water out, salts in. Small/absent glomeruli, little urine, drinks seawater, gills excrete salts.
    \item Long Loop of Henlé $\rightarrow$ steep gradient $\rightarrow$ very concentrated urine (5500 mOsm/L). Metabolic water from seed oxidation. Nocturnal/burrowing reduces loss. Dry faeces.
    \item Osmoconformer: Body fluids iso-osmotic with environment (e.g. jellyfish). Osmoregulator: Body fluids regulated different from environment (e.g. mammal/freshwater fish).
    \item (a) Sigmoidal: steep rise then plateau. (b) Plateau = maximum concentrating ability reached; all aquaporins inserted; limited by medullary gradient steepness.
\end{enumerate}
\end{corrige}

\begin{attention}[Final Exam Traps to Avoid]
\begin{itemize}
    \item \textbf{ADH source:} Hypothalamus (make), Posterior Pituitary (release). Never ``pituitary produces ADH''.
    \item \textbf{Water potential direction:} Water moves \textbf{from high $\Psi_w$ to low $\Psi_w$}. In dehydration, blood $\Psi_w$ is \textbf{low}. In freshwater, environment $\Psi_w$ (0) is \textbf{higher} than fish blood $\Psi_w$ $\rightarrow$ water enters fish.
    \item \textbf{Marine fish drink water:} They \textbf{drink seawater} to replace osmotic loss. Freshwater fish \textbf{do not drink}.
    \item \textbf{Urine concentration limit:} Determined by \textbf{length of Loop of Henlé} (medullary gradient), not just ADH amount. ADH allows you to \emph{use} the gradient; the gradient \emph{sets the maximum}.
    \item \textbf{Contractile vacuole:} Expels water (osmoregulation), not nitrogenous waste (excretion) primarily.
\end{itemize}
\end{attention}

\coursec{Section 7: Excretory Products of Plants, Integration and Evaluation}

Having explored the sophisticated renal osmoregulatory mechanisms of animals in Section 6, we now turn to a fundamentally different strategy for maintaining internal constancy. Plants, as sessile photosynthetic organisms, face distinct metabolic challenges and have evolved a remarkably economical approach to waste management. This final section examines plant excretory products, their elimination pathways, and their surprising economic value, before integrating all homeostatic concepts into a unified framework and preparing you for GCE examination success.

\courssub{Why Plants Have Fewer Excretory Problems Than Animals}

\begin{propriete}[Plants as Metabolic Recyclers]
Unlike animals, plants do not possess specialised excretory organs (kidneys, Malpighian tubules, nephridia). This is because:
\begin{enumerate}
\item \textbf{Lower metabolic rate:} Plants are largely sessile and do not maintain high body temperatures; their catabolic activity is lower, producing fewer nitrogenous wastes per unit biomass.
\item \textbf{Reuse of primary metabolic wastes:} The two major gaseous by-products of metabolism, carbon dioxide ($\ce{CO2}$) and oxygen ($\ce{O2}$), are not waste in the absolute sense. $\ce{CO2}$ from respiration is reused in photosynthesis; $\ce{O2}$ from photosynthesis is reused in respiration. This internal cycling drastically reduces the need for active excretion.
\item \textbf{Storage over elimination:} Many potentially toxic metabolic by-products (alkaloids, tannins, calcium oxalate) are sequestered in vacuoles, old leaves, bark, or heartwood — tissues that are eventually shed or are metabolically inactive.
\item \textbf{Ammonia detoxification:} The small amounts of ammonia produced during protein deamination are immediately incorporated into amino acids (glutamine, asparagine) via the GS/GOGAT cycle, avoiding the need for urea or uric acid synthesis pathways.
\end{enumerate}
\end{propriete}

\begin{attention}[Common Misconception: Oxygen as a "Waste Product"]
\textbf{Do not state categorically that "oxygen is a waste product of plants".} Oxygen is a by-product of \emph{photosynthesis} (light-dependent reactions). For the plant cell, it is a valuable substrate for aerobic respiration. Only when photosynthetic rate far exceeds respiratory rate (bright light) does $\ce{O2}$ diffuse out as a net excretory product. In darkness, plants are net consumers of $\ce{O2}$. Always qualify your statement: \emph{"Oxygen is a by-product of photosynthesis which may be excreted when produced in excess of respiratory requirements."}
\end{attention}

\courssub{Gaseous Excretory Products and Their Exchange}

\begin{definition}[Guttation]
The exudation of liquid water (xylem sap) from hydathodes at leaf margins, driven by positive root pressure when transpiration is low (high humidity, night) and soil water potential is high. It is \emph{not} dew (which condenses from the atmosphere).
\end{definition}

\begin{definition}[Exudation]
The passive or active discharge of substances (gums, resins, latex, nectar) from plant tissues, often in response to injury or through specialised structures (nectaries, laticifers).
\end{definition}

Plants exchange gases primarily through \cle{stomata} (mainly on leaves) and \cle{lenticels} (on woody stems and roots). The direction of net flux depends on the balance between photosynthesis and respiration:

\begin{center}
\begin{tikzpicture}[scale=0.9, every node/.style={font=\small}]
% Leaf cross-section
\draw[popdark, thick] (0,0) -- (8,0); % lower epidermis
\draw[popdark, thick] (0,2) -- (8,2); % upper epidermis
% Mesophyll
\fill[popgreenL!30] (0,0) rectangle (8,2);
% Vascular bundle
\draw[poporange, thick] (4,0) -- (4,2);
\draw[poporange, thick] (3.8,0) -- (3.8,2);
\draw[poporange, thick] (4.2,0) -- (4.2,2);
% Stomata
\draw[popblue, thick] (2,0) -- (2.5,0); % closed
\draw[popblue, thick] (2,0.05) -- (2.5,0.05);
\draw[popblue, thick] (5.5,0) -- (6,0); % open
\draw[popblue, thick] (5.5,0.15) -- (6,0.15);
\draw[popblue, thick] (5.5,0) -- (5.5,0.15);
\draw[popblue, thick] (6,0) -- (6,0.15);
% Air spaces
\fill[popblue!10] (1,0.5) circle (0.4);
\fill[popblue!10] (7,0.5) circle (0.4);
% Arrows with integrated nodes
\draw[->, poporange, thick] (2.25,-0.5) -- (2.25,0) node[below] {$\ce{CO2}$ in (day)};
\draw[->, poporange, thick] (5.75,0.15) -- (5.75,0.6) node[right] {$\ce{O2}$ out (day)};
\draw[->, poppurple, thick] (5.75,-0.5) -- (5.75,0) node[below] {$\ce{O2}$ in (night)};
\draw[->, poppurple, thick] (2.25,0.05) -- (2.25,0.5) node[left] {$\ce{CO2}$ out (night)};
% Lenticel on stem (inset)
\begin{scope}[shift={(9,0.5)}, scale=0.6]
\draw[popdark, thick] (-0.5,0) -- (0.5,0);
\draw[popdark, thick] (-0.5,1) -- (0.5,1);
\fill[poporange!30] (-0.5,0) rectangle (0.5,1);
\draw[poporange, thick] (-0.2,0.3) -- (0.2,0.3);
\draw[poporange, thick] (-0.2,0.7) -- (0.2,0.7);
\draw[->, popblue] (0.5,0.5) -- (1,0.5) node[right] {Gas exchange};
\end{scope}
% Labels
\node[align=center] at (1,1.5) {Upper epidermis};
\node[align=center] at (1,0.2) {Lower epidermis};
\node[align=center] at (4,1) {Vascular\\bundle};
\node[align=center] at (1,0.5) {Spongy\\mesophyll};
\node[align=center] at (7,0.5) {Air space};
\node[align=center] at (2.25,-0.8) {Stoma (closed)};
\node[align=center] at (5.75,-0.8) {Stoma (open)};
\node[align=center] at (9.5,1.2) {Lenticel};
\end{tikzpicture}
\legende{Leaf cross-section showing stomatal gas exchange (day vs night) and a lenticel on a woody stem.}
\end{center}

\begin{aretenir}[Gaseous Exchange Summary]
\begin{itemize}
\item \textbf{Day (light):} Net $\ce{CO2}$ uptake, net $\ce{O2}$ release (photosynthesis $>$ respiration). Water vapour lost by transpiration.
\item \textbf{Night (dark):} Net $\ce{O2}$ uptake, net $\ce{CO2}$ release (respiration only). Transpiration reduced (stomata often closed).
\item \textbf{Stomatal regulation:} Guard cells integrate light, $\ce{CO2}$ concentration, humidity, and hormonal signals (ABA) to balance gas exchange with water conservation — a homeostatic feedback loop.
\end{itemize}
\end{aretenir}

\courssub{Non-Gaseous Excretory Products: Secondary Metabolites}

Many plant "wastes" are \cle{secondary metabolites} — compounds not directly involved in primary growth but crucial for defence, signalling, and ecological interactions. Their accumulation in vacuoles or cell walls effectively removes them from metabolic pools.

\begin{center}
\begin{tabularx}{\linewidth}{|l|X|X|}\hline
\rowcolor{popblueL}
\textbf{Class} & \textbf{Examples} & \textbf{Storage / Elimination} \\ \hline
\textbf{Tannins} & Hydrolysable (gallotannins), condensed (proanthocyanidins) & Vacuoles of bark, leaves, unripe fruits; astringent taste deters herbivores. \\ \hline
\textbf{Resins} & Diterpene acids (pine resin), essential oils & Secreted into resin canals/ducts; exude on injury to seal wounds. \\ \hline
\textbf{Gums} & Polysaccharides (gum arabic from \emph{Acacia senegal}) & Exuded from bark wounds; seal injury, prevent pathogen entry. \\ \hline
\textbf{Latex} & Complex emulsion: rubber (polyisoprene), alkaloids, proteins & Laticifers (articulated or non-articulated); exuded on cutting. \\ \hline
\textbf{Alkaloids} & Quinine (\emph{Cinchona}), caffeine (\emph{Coffea}), nicotine (\emph{Nicotiana}), morphine (\emph{Papaver}) & Vacuoles; often in roots, bark, seeds. Potent physiological effects on animals. \\ \hline
\textbf{Essential oils} & Menthol, eucalyptol, limonene & Oil glands/trichomes on leaf surface; volatile, deter herbivores, attract pollinators. \\ \hline
\textbf{Calcium oxalate crystals} & Raphides (needles), druses (clusters), prismatic crystals & Idioblasts (specialised cells); regulate $\ce{Ca^{2+}}$ homeostasis, deter herbivory (raphides cause mechanical irritation). \\ \hline
\end{tabularx}
\end{center}

\begin{savaistu}[Cameroon's Green Pharmacy and Economy]
Cameroon's biodiversity is a reservoir of valuable plant exudates:
\begin{itemize}
\item \textbf{Rubber latex:} \emph{Hevea brasiliensis} (introduced) and wild \emph{Landolphia} spp. (historically). Latex coagulation yields natural rubber — vital for tyres, medical gloves.
\item \textbf{Gum arabic:} \emph{Acacia senegal} and \emph{A. seyal} in the Far North (Sahel zone). Major export; used as emulsifier (E414) in food, pharmaceuticals, printing.
\item \textbf{Medicinal alkaloids:} \emph{Cinchona} (quinine — antimalarial, crucial in Cameroon), \emph{Rauwolfia vomitoria} (reserpine — antihypertensive), \emph{Voacanga africana} (voacangine — precursor for semi-synthetic drugs).
\item \textbf{Essential oils:} \emph{Eucalyptus} (cineole), \emph{Cymbopogon} (citronella) — cultivated for local and export markets.
\item \textbf{Tannins:} Bark of \emph{Terminalia} spp., \emph{Anogeissus leiocarpus} — traditional leather tanning, potential for modern adhesives.
\end{itemize}
These "wastes" sustain livelihoods and illustrate the principle: \emph{one organism's excretion is another's resource}.
\end{savaistu}

\courssub{Methods of Elimination in Plants}

Plants lack a circulatory system dedicated to waste transport to excretory organs. Elimination is largely passive and structural:

\begin{enumerate}
\item \textbf{Diffusion through stomata and lenticels:} Gaseous wastes ($\ce{CO2}$, $\ce{O2}$, volatile oils) and water vapour.
\item \textbf{Storage in senescing organs:} Before leaf abscission or bark sloughing, nutrients (N, P, K) are actively retranslocated (phloem mobility), while wastes (tannins, heavy metals, calcium oxalate) are left behind. The shed organ carries the waste away.
\item \textbf{Exudation:}
\begin{itemize}
\item \emph{Guttation:} Hydathodes release xylem sap (water + minerals) under root pressure.
\item \emph{Gum/resin exudation:} Often triggered by mechanical damage; seals wounds.
\item \emph{Latex flow:} Pressure-driven from laticifers upon incision.
\end{itemize}
\item \textbf{Heartwood formation:} In woody perennials, the inner xylem (heartwood) becomes metabolically inactive and is impregnated with resins, tannins, oils, and crystals — a massive, long-term waste sink.
\end{enumerate}

\begin{exemplebox}[Calcium Oxalate as a Calcium Homeostat]
Calcium enters the xylem as $\ce{Ca^{2+}}$ and moves passively with the transpiration stream. Excess $\ce{Ca^{2+}}$ in leaf cells would precipitate phosphate, disrupt signalling, and cause toxicity. Plants solve this by precipitating $\ce{Ca^{2+}}$ as insoluble calcium oxalate ($\ce{CaC2O4}$) crystals in vacuoles of specialised \cle{idioblasts}. This:
\begin{itemize}
\item Removes toxic free $\ce{Ca^{2+}}$ from cytoplasm.
\item Provides a calcium reserve that can be remobilised (via oxalate oxidase) if needed.
\item Forms raphides that deter herbivores (mechanical + chemical irritation).
\end{itemize}
This is a beautiful example of \emph{homeostasis via controlled precipitation and storage}.
\end{exemplebox}

\courssub{Integration Activity: A Homeostatic Case Study}

\begin{experience}[Integrated Homeostasis: The Marathon Runner in Douala]
\textbf{Scenario:} A 25-year-old athlete runs a 42 km marathon in Douala (temperature \SI{32}{\ensuremath{{}^{\circ}\mathrm{C}}}, humidity \SI{80}{\percent}). Analyse the coordinated homeostatic responses.

\textbf{Task:} For each challenge below, identify the \emph{sensor}, \emph{control centre}, \emph{effector}, \emph{feedback type}, and the \emph{organ(s)} involved. Explain the mechanism.

\begin{enumerate}
\item \textbf{Core temperature rises to \SI{39.5}{\ensuremath{{}^{\circ}\mathrm{C}}}.}
\item \textbf{Blood glucose falls to \SI{3.2}{\mmol\per\liter} at 30 km.}
\item \textbf{Blood $\ce{Na+}$ concentration increases due to sweat loss.}
\item \textbf{Blood pH drops to 7.25 (lactic acidosis).}
\item \textbf{Blood pressure drops due to cutaneous vasodilation and fluid loss.}
\end{enumerate}

\textbf{Guideline for solution:} Use the \begin{methode}Answering GCE Essay Questions on Homeostasis\end{methode} below. A worked solution follows in the \begin{corrige}\end{corrige} box.
\end{experience}

\begin{methode}[Answering GCE Essay Questions on Homeostasis]
\textbf{Step 1: Define the key term.} (e.g., "Homeostasis is the maintenance of a constant internal environment...")

\textbf{Step 2: Identify the controlled variable and its set point.} (e.g., "Blood glucose set point $\approx \SI{5}{\mmol\per\liter}$.")

\textbf{Step 3: Describe the sensor/receptor location and stimulus.}

\textbf{Step 4: Name the control centre (hypothalamus, medulla, pancreatic islets, juxtaglomerular apparatus).}

\textbf{Step 5: Detail the effector response (organ, hormone, mechanism).}

\textbf{Step 6: Explicitly state the feedback loop:} "This is \emph{negative feedback} because the response \emph{reduces} the original stimulus, restoring the set point."

\textbf{Step 7: Mention antagonistic mechanisms if applicable (e.g., insulin vs glucagon; sympathetic vs parasympathetic).}

\textbf{Step 8: Conclude with physiological significance.}
\end{methode}

\begin{corrige}[Marathon Case Study]
\begin{enumerate}
\item \textbf{Thermoregulation:} Hypothalamic thermoreceptors (and skin receptors) $\to$ hypothalamus $\to$ \textbf{effectors:} sweat glands (secretion), cutaneous arterioles (vasodilation), erector pili muscles (relaxation), increased respiratory rate. \textbf{Negative feedback:} Evaporative cooling lowers blood temperature $\to$ hypothalamus reduces signals. \textbf{Note:} High humidity impairs evaporation $\to$ risk of heat stroke (positive feedback danger).
\item \textbf{Blood glucose:} Pancreatic $\alpha$-cells detect low glucose $\to$ secrete \textbf{glucagon} $\to$ liver: glycogenolysis + gluconeogenesis $\to$ glucose released into blood. Adrenal medulla: adrenaline amplifies glycogenolysis. \textbf{Negative feedback:} Rising glucose inhibits $\alpha$-cells, stimulates $\beta$-cells (insulin) later.
\item \textbf{Osmoregulation ($\ce{Na+}$):} Hypothalamic osmoreceptors detect increased plasma osmolarity $\to$ posterior pituitary releases \textbf{ADH} $\to$ collecting ducts become water-permeable (aquaporins) $\to$ water reabsorbed $\to$ plasma diluted. Simultaneously, low BP $\to$ JG apparatus releases \textbf{renin} $\to$ angiotensin II $\to$ adrenal cortex releases \textbf{aldosterone} $\to$ distal tubule/collecting duct reabsorbs $\ce{Na+}$ (water follows). \textbf{Negative feedback:} Restored osmolarity/BP inhibits ADH/renin.
\item \textbf{pH regulation:} Central/peripheral chemoreceptors detect low pH/high $\ce{H+}$ $\to$ respiratory centre $\to$ increased ventilation (blow off $\ce{CO2}$) $\to$ $\ce{H2CO3}$ decreases $\to$ pH rises. Renal: proximal tubule secretes $\ce{H+}$, reabsorbs $\ce{HCO3-}$; collecting duct $\ce{H+}$ secretion (titratable acidity, $\ce{NH4+}$). \textbf{Negative feedback.}
\item \textbf{Blood pressure:} Baroreceptors (carotid sinus, aortic arch) detect low BP $\to$ vasomotor centre $\to$ sympathetic stimulation $\to$ vasoconstriction (skin, splanchnic), increased heart rate/contractility. ADH + aldosterone conserve volume. \textbf{Negative feedback.}
\end{enumerate}
\end{corrige}

\courssub{Chapter Synthesis: Concept Map of Homeostasis}

The following concept map integrates all major homeostatic organs, controlled variables, hormones, and feedback loops covered in this chapter. Study it carefully — it is your mental framework for the examination.

\begin{center}
\begin{tikzpicture}[
  concept/.style={draw=popblue, thick, fill=popblueL, rounded corners, minimum width=2.2cm, minimum height=1cm, align=center, font=\small},
  organ/.style={draw=popgreen, thick, fill=popgreenL, rounded corners, minimum width=2cm, minimum height=0.9cm, align=center, font=\small},
  hormone/.style={draw=poporange, thick, fill=poporange!20, ellipse, minimum width=2cm, minimum height=0.8cm, align=center, font=\small},
  varbox/.style={draw=poppurple, thick, fill=poppurple!20, rectangle, minimum width=2cm, minimum height=0.8cm, align=center, font=\small},
  arrow/.style={->, thick, popdark},
  lbl/.style={font=\tiny, popmuted, midway, sloped, allow upside down}
]
% Organs
\node[organ] (liver) at (-5, 2) {Liver};
\node[organ] (skin) at (0, 3.5) {Skin};
\node[organ] (kidney) at (5, 2) {Kidneys};
\node[organ] (hypoth) at (0, 0.5) {Hypothalamus};
\node[organ] (pancreas) at (-2.5, -1.5) {Pancreas};
\node[organ] (adrenal) at (2.5, -1.5) {Adrenal Glands};
\node[organ] (pituitary) at (0, -0.5) {Pituitary};
\node[organ] (heart) at (3.5, 3.5) {Heart};

% Variables
\node[varbox] (temp) at (-5, 4) {Core Temp};
\node[varbox] (gluc) at (-5, 0) {Blood Glucose};
\node[varbox] (osmo) at (5, 4) {Blood Osmolarity};
\node[varbox] (bp) at (5, 0) {Blood Pressure};
\node[varbox] (ph) at (0, -3) {Blood pH};
\node[varbox] (nitro) at (-2.5, 3) {Nitrogenous Waste};

% Hormones
\node[hormone] (insulin) at (-4, -1) {Insulin};
\node[hormone] (glucagon) at (-1, -1) {Glucagon};
\node[hormone] (adh) at (1, -1) {ADH};
\node[hormone] (aldo) at (4, -1) {Aldosterone};
\node[hormone] (adren) at (0, -2) {Adrenaline};
\node[hormone] (thyrox) at (-2.5, -2.5) {Thyroxine};

% Connections
\draw[arrow] (temp) -- (hypoth) node[lbl, above] {thermoreceptors};
\draw[arrow] (hypoth) -- (skin) node[lbl, above] {sympathetic/parasympathetic};
\draw[arrow] (skin) -- (temp) node[lbl, below] {vasodilation, sweat};

\draw[arrow] (gluc) -- (pancreas) node[lbl, left] {$\alpha/\beta$ cells};
\draw[arrow] (pancreas) -- (insulin) node[lbl, above] {high glucose};
\draw[arrow] (pancreas) -- (glucagon) node[lbl, above] {low glucose};
\draw[arrow] (insulin) -- (liver) node[lbl, above] {glycogen synthesis};
\draw[arrow] (glucagon) -- (liver) node[lbl, above] {glycogenolysis};
\draw[arrow] (liver) -- (gluc) node[lbl, left] {glucose release/uptake};

\draw[arrow] (osmo) -- (hypoth) node[lbl, right] {osmoreceptors};
\draw[arrow] (hypoth) -- (pituitary) node[lbl, left] {neurosecretion};
\draw[arrow] (pituitary) -- (adh) node[lbl, above] {release};
\draw[arrow] (adh) -- (kidney) node[lbl, above] {collecting duct};
\draw[arrow] (kidney) -- (osmo) node[lbl, right] {water reabsorption};

\draw[arrow] (bp) -- (kidney) node[lbl, right] {baroreceptors/JG app.};
\draw[arrow] (kidney) -- (aldo) node[lbl, above] {renin-angiotensin};
\draw[arrow] (adrenal) -- (aldo) node[lbl, left] {cortex};
\draw[arrow] (aldo) -- (kidney) node[lbl, below] {Na$^{+}$ reabsorption};
\draw[arrow] (kidney) -- (bp) node[lbl, right] {volume/pressure};

\draw[arrow] (ph) -- (hypoth) node[lbl, left] {chemoreceptors};
\draw[arrow] (hypoth) -- (kidney) node[lbl, right] {ventilation/renal};
\draw[arrow] (kidney) -- (ph) node[lbl, right] {H$^{+}$ secretion, HCO$_{3}^{-}$ reabs.};

\draw[arrow] (nitro) -- (liver) node[lbl, left] {deamination};
\draw[arrow] (liver) -- (kidney) node[lbl, above] {urea via blood};
\draw[arrow] (kidney) -- (nitro) node[lbl, right] {excretion};

\draw[arrow] (adren) -- (liver) node[lbl, above] {glycogenolysis};
\draw[arrow] (adren) -- (skin) node[lbl, above] {vasoconstriction};
\draw[arrow] (adren) -- (heart) node[lbl, right] {HR/contractility};
\draw[arrow] (heart) -- (bp) node[lbl, below] {cardiac output};

% Feedback loop labels
\node[font=\tiny, poporange, rotate=45] at (-3, 3) {Negative Feedback};
\node[font=\tiny, poporange, rotate=-45] at (3, 3) {Negative Feedback};
\node[font=\tiny, poporange] at (0, -2.5) {Negative Feedback};
\end{tikzpicture}
\legende{Integrated concept map: Homeostatic organs (green), controlled variables (purple), hormones (orange), and negative feedback loops (orange labels). Arrows show stimulus $\to$ sensor $\to$ control centre $\to$ effector $\to$ response.}
\end{center}

\courssub{Evaluation: GCE-Style Practice Questions}

\begin{exercice}[Structured Questions]
\begin{enumerate}
\item \textbf{(a)} Define the term \emph{homeostasis}. [2 marks]
\item \textbf{(b)} Explain how the liver maintains a constant blood glucose concentration. [6 marks]
\item \textbf{(c)} Describe the role of the skin in thermoregulation when body temperature rises above the set point. [5 marks]
\item \textbf{(d)} Outline the hormonal control of water reabsorption in the kidney collecting duct. [4 marks]
\item \textbf{(e)} Explain why plants do not need specialised excretory organs. [3 marks]
\end{enumerate}
\end{exercice}

\begin{exercice}[Essay Questions]
\begin{enumerate}
\item \textbf{``Homeostasis in mammals is achieved through negative feedback mechanisms involving the nervous and endocrine systems.''} Discuss this statement with reference to \textbf{three} different homeostatic processes. [15 marks]
\item Compare and contrast the excretory strategies of a mammal and a flowering plant. In your answer, refer to the nature of waste products, organs/structures involved, and the fate of metabolic wastes. [15 marks]
\item Describe the formation of urine in the mammalian nephron, explaining how the composition of the filtrate is modified to produce a final urine that is hypertonic to blood plasma. [15 marks]
\end{enumerate}
\end{exercice}

\begin{corrige}[Structured Questions - Marking Guidance]
\begin{enumerate}
\item \textbf{(a)} Maintenance of a constant internal environment (1) despite external fluctuations (1). Key variables: temperature, pH, glucose, water potential, $\ce{CO2}$, $\ce{O2}$.
\item \textbf{(b)} Pancreatic $\beta$-cells detect high glucose $\to$ insulin $\to$ liver: glycogenesis, glycolysis, lipogenesis (3); $\alpha$-cells detect low glucose $\to$ glucagon $\to$ liver: glycogenolysis, gluconeogenesis (3). Mention negative feedback.
\item \textbf{(c)} Hypothalamic thermoreceptors $\to$ heat loss centre $\to$ sympathetic inhibition of vasoconstrictors $\to$ cutaneous vasodilation (1); sweat glands activated $\to$ evaporation (1); erector pili relax $\to$ hair lies flat (1); behavioural responses (1); negative feedback restores set point (1).
\item \textbf{(d)} Increased plasma osmolarity $\to$ hypothalamic osmoreceptors $\to$ posterior pituitary releases ADH $\to$ ADH binds V2 receptors on collecting duct principal cells $\to$ cAMP $\to$ aquaporin-2 vesicles fuse with apical membrane $\to$ water reabsorbed by osmosis into hypertonic medulla $\to$ concentrated urine. Negative feedback.
\item \textbf{(e)} Lower metabolic rate (1); reuse of $\ce{CO2}$/$\ce{O2}$ (1); storage in vacuoles/senescent organs (1); no nitrogenous waste accumulation (ammonia reassimilated) (1).
\end{enumerate}
\end{corrige}

\begin{corrige}[Essay Question 1 - Outline Plan]
\textbf{Introduction:} Define homeostasis, negative feedback, roles of nervous (fast) and endocrine (slow, sustained) systems.

\textbf{Process 1: Thermoregulation}
\begin{itemize}
\item Sensors: hypothalamic + peripheral thermoreceptors.
\item Control: hypothalamus (heat gain/loss centres).
\item Effectors: skin (vasomotion, sweating, piloerection), skeletal muscle (shivering), liver (metabolic rate), behaviour.
\item Hormones: thyroxine (long-term), adrenaline (acute).
\item Feedback: Negative — response opposes stimulus.
\end{itemize}

\textbf{Process 2: Blood Glucose Regulation}
\begin{itemize}
\item Sensors: pancreatic $\alpha/\beta$ cells (also hypothalamic glucoreceptors).
\item Control: Islets of Langerhans (endocrine pancreas).
\item Effectors: Liver (glycogenolysis/genesis, gluconeogenesis), adipose, muscle.
\item Hormones: Insulin (lowers), Glucagon (raises), Adrenaline, Cortisol, Growth Hormone.
\item Feedback: Negative — antagonistic hormones.
\end{itemize}

\textbf{Process 3: Osmoregulation (Water Balance)}
\begin{itemize}
\item Sensors: Hypothalamic osmoreceptors, baroreceptors (carotid/aortic, JG apparatus).
\item Control: Hypothalamus (ADH synthesis), JG apparatus (renin).
\item Effectors: Kidney collecting ducts (ADH), distal tubule (Aldosterone), thirst centre.
\item Hormones: ADH, Aldosterone, Angiotensin II, ANP (atrial natriuretic peptide — opposes).
\item Feedback: Negative — restored osmolarity/volume inhibits hormone release.
\end{itemize}

\textbf{Conclusion:} Integration — e.g., exercise challenges all three; nervous system provides rapid response, endocrine sustains it; failure leads to pathology (diabetes, dehydration, heat stroke).
\end{corrige}

\begin{attention}[GCE Examination Traps to Avoid]
\begin{itemize}
\item \textbf{Confusing positive and negative feedback:} Positive feedback \emph{amplifies} the stimulus (e.g., oxytocin in childbirth, action potential depolarisation). Homeostasis relies almost entirely on \textbf{negative feedback}. Never write "positive feedback restores the set point".
\item \textbf{Omitting the feedback loop:} Describing the pathway (stimulus $\to$ response) without explicitly stating "this is negative feedback because the response reduces the stimulus" loses 1–2 marks per essay.
\item \textbf{Vague hormone actions:} "ADH increases water reabsorption" is insufficient. Say: "ADH inserts aquaporins into the apical membrane of collecting duct cells, increasing water permeability."
\item \textbf{Mixing up glucagon/insulin triggers:} $\alpha$-cells $\to$ glucagon $\to$ \textbf{low} glucose; $\beta$-cells $\to$ insulin $\to$ \textbf{high} glucose.
\item \textbf{Forgetting the liver's role in deamination/urea synthesis:} The liver \emph{produces} urea; the kidney \emph{excretes} it. They are not the same step.
\item \textbf{Plant excretion:} Do not list $\ce{O2}$ as a waste product without the photosynthesis context. Do not invent "plant kidneys".
\end{itemize}
\end{attention}

\begin{aretenir}[Chapter 7 — Key Takeaways for the Exam]
\begin{enumerate}
\item Plants \textbf{recycle} $\ce{CO2}$ and $\ce{O2}$; they \textbf{store} most other wastes (tannins, alkaloids, crystals) in vacuoles or senescent tissues.
\item Gaseous exchange occurs via \textbf{stomata} (leaves) and \textbf{lenticels} (woody stems); direction depends on photosynthesis/respiration balance.
\item \textbf{Guttation} = root pressure + hydathodes (liquid water); \textbf{Transpiration} = evaporation from stomata (water vapour).
\item Plant "wastes" (latex, gums, alkaloids, essential oils) have immense \textbf{economic and medicinal value} in Cameroon and globally.
\item \textbf{Integration:} Every homeostatic response follows: \textbf{Stimulus $\to$ Receptor $\to$ Control Centre $\to$ Effector $\to$ Response $\to$ Negative Feedback}.
\item \textbf{Essay technique:} Define $\to$ Variable/Set point $\to$ Sensor $\to$ Control Centre $\to$ Effector/Mechanism $\to$ Hormones/Nerves $\to$ \textbf{Explicit Negative Feedback Statement} $\to$ Significance.
\end{enumerate}
\end{aretenir}

\begin{savaistu}[Final Thought: Homeostasis as the Essence of Life]
From the calcium oxalate crystal in a \emph{Dieffenbachia} leaf to the aquaporin channel in your collecting duct; from the tannin-laden bark of a Mount Cameroon \emph{Prunus africana} to the glucagon secreted by your pancreatic $\alpha$-cells — homeostasis is the universal dialogue between an organism and its environment. It is not static perfection, but dynamic resilience. As you sit for your GCE, remember: the examiner rewards \textbf{mechanistic clarity} and \textbf{feedback logic} above rote memorisation. You have the map; now navigate with confidence.
\end{savaistu}

\newpage

\coursec{Integration activity}

\coursec{Integration and Evaluation}

\courssub{Aims of the activity}

By the end of this activity, you should be able to:

\begin{itemize}
\item \cle{Integrate} your knowledge of the \cle{liver}, the \cle{skin}, the \cle{kidneys} and \cle{feedback mechanisms} to explain how the internal environment of a whole organism is kept constant during intense exercise in a hot climate.
\item \cle{Interpret} and \cle{plot} physiological data (tables and graphs) as required in GCE Advanced Level Paper 2.
\item \cle{Explain} each observed change by naming the correct \cle{receptor}, \cle{coordination centre}, \cle{effector} and \cle{response} (the components of a feedback loop).
\item \cle{Predict} the consequences of a disturbance (drinking excess plain water) using the principle of \cle{negative feedback}.
\item \cle{Communicate} a complete, annotated homeostatic diagram in the style rewarded by GCE marking schemes.
\end{itemize}

\begin{aretenir}[Skills mobilised]
Data handling (plotting, reading values, calculating rates and percentage changes); biological explanation (thermoregulation, blood glucose regulation, osmoregulation by ADH); prediction and evaluation of hypotheses; scientific drawing and annotation.
\end{aretenir}

\begin{savaistu}[The Mount Cameroon Race of Hope]
The \cle{Mount Cameroon Race of Hope} is an annual international mountain race held in February. Athletes run from Molyko Stadium in Buea (approx. \SI{200}{m} altitude) to the summit of Mount Cameroon (\SI{4095}{m}) and back, a round trip of about \SI{38}{km}. The lower slopes are tropical (hot, humid), while the summit can be cool and windy. This extreme variation makes it a remarkable natural laboratory for studying human homeostasis under thermal, metabolic and osmotic stress.
\end{savaistu}

\courssub{Situation}

\begin{experience}[Case study: The Marathon Runner of Mount Cameroon]
Every year in February, athletes from all over the world take part in the \cle{Mount Cameroon Race of Hope}, running from Molyko Stadium in Buea up to the summit of Mount Cameroon (about \SI{4100}{m}) and back. The lower slopes are hot and humid: on race day the air temperature in Buea is about \SI{30}{\ensuremath{{}^{\circ}\mathrm{C}}} with high humidity, conditions in which the body gains heat much faster than it can lose it by radiation.

Ekane, a 24-year-old athlete from Buea, was monitored by a team of sports physiologists during the race. A small cannula allowed blood samples to be taken, sweat was collected from a patch on the forearm, and urine was collected before and after the race. The results are summarised below.

\textbf{Table 1: Physiological measurements before, during and after the race}

\begin{center}
\begin{tabularx}{\linewidth}{|l|X|X|X|}
\hline
\rowcolor{popblueL} \textbf{Measurement} & \textbf{Before race (07:00)} & \textbf{During race (at summit)} & \textbf{After race (2 h rest)} \\
\hline
Core body temperature (\si{\ensuremath{{}^{\circ}\mathrm{C}}}) & 36.8 & 39.2 & 37.1 \\
\hline
Blood glucose (\si{mg/100\,cm^{3}}) & 85 & 110 & 88 \\
\hline
Sweat rate (\si{dm^{3}.h^{-1}}) & 0.1 & 1.4 & 0.3 \\
\hline
Urine production (\si{cm^{3}.h^{-1}}) & 60 & 15 & 35 \\
\hline
Urine concentration (arbitrary units) & 1.0 & 3.2 & 1.8 \\
\hline
Blood ADH concentration (arbitrary units) & 1.0 & 4.5 & 1.6 \\
\hline
\end{tabularx}
\end{center}

Additional observations made by the medical team:
\begin{itemize}
\item During the climb, Ekane's skin was flushed (reddened) and dripping with sweat.
\item He drank only \SI{500}{cm^{3}} of water during the whole race.
\item After the race, a fellow athlete who had drunk \SI{4}{dm^{3}} of plain water in one hour felt dizzy, nauseous and confused, and had to be treated by the medical team.
\end{itemize}
\end{experience}

\courssub{Tasks}

Work in groups of four. Answer the following questions in order; each one builds on the previous one.

\begin{enumerate}
\item \textbf{Data handling.} Using graph paper, plot a bar chart (or a line graph with time on the $x$-axis) showing how \emph{sweat rate}, \emph{urine production} and \emph{blood ADH concentration} vary before, during and after the race. State the relationship you observe between blood ADH concentration and urine production.

\item \textbf{Thermoregulation.} Ekane's core temperature rose from \SI{36.8}{\ensuremath{{}^{\circ}\mathrm{C}}} to \SI{39.2}{\ensuremath{{}^{\circ}\mathrm{C}}}. 
\begin{enumerate}
\item Explain why muscular exercise in a hot, humid climate raises body temperature.
\item Describe the \cle{negative feedback loop} by which the body responds, naming the \cle{receptors}, the \cle{coordination centre} and the \cle{effectors} in the skin, and explain how each response (vasodilation, sweating) reduces body temperature.
\item Explain why the high humidity in Buea makes sweating less effective at cooling the body.
\end{enumerate}

\item \textbf{Blood glucose regulation.} Blood glucose rose from 85 to \SI{110}{mg/100\,cm^{3}} during the race.
\begin{enumerate}
\item Explain why the muscles need an increased supply of glucose during the race.
\item Name the hormone(s) responsible for this rise and describe, step by step, how the \cle{liver} brings it about (glycogenolysis and gluconeogenesis).
\item Predict, with reasons, what would happen to Ekane's blood glucose if the race lasted several more hours without food, and name the hormone that would then become dominant.
\end{enumerate}

\item \textbf{Osmoregulation.} 
\begin{enumerate}
\item Using the data, explain why urine production fell and urine concentration rose during the race.
\item Describe the mechanism by which \cle{ADH} acts on the kidney (name the part of the nephron concerned and the effect on water permeability).
\item Draw the negative feedback loop linking blood water potential, the hypothalamus, the pituitary gland, ADH and the kidney.
\end{enumerate}

\item \textbf{Prediction.} The athlete who drank \SI{4}{dm^{3}} of plain water in one hour became ill. Using your knowledge of osmoregulation, explain:
\begin{enumerate}
\item what happened to the water potential of his blood;
\item how ADH secretion and the kidneys responded;
\item why excessive dilution of the blood (hyponatraemia) can cause dizziness and confusion.
\end{enumerate}

\item \textbf{Synthesis.} On a large sheet of paper, draw a single annotated diagram showing how the \cle{skin}, \cle{liver} and \cle{kidneys}, coordinated by the nervous and endocrine systems, work together to maintain homeostasis during the race. Your diagram must show at least three complete feedback loops, each with receptor, coordination centre, effector and response clearly labelled.

\item \textbf{Evaluation.} Suggest two limitations of the data collected and one additional measurement the physiologists should have made to improve the study.
\end{enumerate}

\begin{methode}[Drawing a negative feedback loop for GCE Biology]
\begin{enumerate}
\item Identify the \textbf{controlled variable} (e.g. blood temperature, blood glucose, blood water potential).
\item Draw a \textbf{circle or box} for the \textbf{receptor} (e.g. thermoreceptors in skin/hypothalamus, $\alpha$-cells of pancreas, osmoreceptors in hypothalamus).
\item Draw a \textbf{box} for the \textbf{coordination centre} (e.g. hypothalamus, pancreas, hypothalamus/posterior pituitary).
\item Draw a \textbf{box} for the \textbf{effector} (e.g. skin arterioles/sweat glands, liver hepatocytes, collecting duct cells).
\item Draw an \textbf{arrow} from effector to response (e.g. vasodilation + sweating, glycogenolysis/gluconeogenesis, increased water permeability).
\item Draw a \textbf{return arrow} from the response back to the controlled variable, labelled \textbf{negative feedback} (the response opposes the initial change).
\item Use \textbf{distinct colours} or line styles for each loop if combining several on one diagram.
\end{enumerate}
\end{methode}

\courssub{Model answer}

\begin{corrige}[Integration activity — model answer and marking grid]

\textbf{Task 1.} The graph shows sweat rate rising sharply during the race (0.1 to \SI{1.4}{dm^{3}.h^{-1}}) while urine production falls (60 to \SI{15}{cm^{3}.h^{-1}}) and blood ADH rises (1.0 to 4.5 arbitrary units). \textbf{Relationship:} as blood ADH concentration increases, urine production decreases (an inverse relationship); both return towards normal after the race. \emph{(2 marks: correctly labelled axes and plotted points; 1 mark for the stated relationship.)}

\textbf{Task 2(a).} Muscle contraction releases energy: only about a quarter of the energy from respiration is converted into movement, the rest is released as \cle{heat}. At \SI{30}{\ensuremath{{}^{\circ}\mathrm{C}}} the air is close to skin temperature, so little heat is lost by radiation and conduction; heat therefore accumulates and core temperature rises.

\textbf{Task 2(b).} The feedback loop:
\begin{itemize}
\item \textbf{Receptors:} thermoreceptors in the skin detect the high external temperature, and the \cle{hypothalamus} detects the temperature of the blood flowing through it.
\item \textbf{Coordination centre:} the heat-loss centre of the \cle{hypothalamus}.
\item \textbf{Effectors and responses:} (i) arterioles supplying the skin \cle{vasodilate}, so more warm blood flows near the surface and more heat is lost by radiation — this causes the flushed skin; (ii) \cle{sweat glands} secrete more sweat, and the \cle{evaporation} of sweat removes heat from the skin (evaporation requires latent heat, taken from the body).
\end{itemize}
These responses lower blood temperature, which reduces the stimulus to the hypothalamus: this is \cle{negative feedback}.

\textbf{Task 2(c).} Sweat only cools the body when it \emph{evaporates}. In humid air the air is already nearly saturated with water vapour, so sweat evaporates slowly and drips off instead, cooling the body much less effectively.

\textbf{Task 3(a).} Contracting muscles respire aerobically at a very high rate and need a continuous supply of \cle{glucose} as respiratory substrate to produce ATP.

\textbf{Task 3(b).} As blood glucose begins to fall during exercise, the \cle{islets of Langerhans} (alpha cells) of the pancreas secrete \cle{glucagon} (and adrenaline is also released from the adrenal medulla). These hormones act on \cle{liver cells (hepatocytes)}, which: (i) accelerate \cle{glycogenolysis} — the breakdown of stored glycogen to glucose; and (ii) carry out \cle{gluconeogenesis} — the production of glucose from amino acids (from muscle protein breakdown) and glycerol (from lipolysis). Glucose is released into the blood, raising blood glucose concentration.

\textbf{Task 3(c).} If the race continued for hours without food, liver glycogen stores would become exhausted, blood glucose would \emph{fall} (hypoglycaemia) and glucagon secretion would increase further while insulin secretion is suppressed; the athlete would fatigue ("hit the wall") because muscles and the brain would lack glucose.

\textbf{Task 4(a).} During the race, large volumes of water are lost in sweat (\SI{1.4}{dm^{3}.h^{-1}}), so the \cle{water potential of the blood falls} (blood becomes more concentrated). This is detected by \cle{osmoreceptors} in the hypothalamus, which stimulate the release of more \cle{ADH} (confirmed by the rise from 1.0 to 4.5 units). ADH makes the kidney reabsorb more water, so a \emph{small volume} of \emph{concentrated} urine is produced (15 instead of \SI{60}{cm^{3}.h^{-1}}; concentration 3.2 instead of 1.0 units).

\textbf{Task 4(b).} ADH, secreted by the \cle{posterior pituitary gland}, travels in the blood to the kidney and increases the \cle{permeability of the walls of the collecting ducts} (and distal convoluted tubules) to water. More water is therefore reabsorbed by osmosis into the hypertonic medulla and back into the blood, leaving a small volume of concentrated urine.

\textbf{Task 4(c).} The feedback loop: fall in blood water potential $\rightarrow$ osmoreceptors in hypothalamus $\rightarrow$ posterior pituitary releases more ADH $\rightarrow$ collecting ducts more permeable $\rightarrow$ more water reabsorbed $\rightarrow$ blood water potential rises back to normal $\rightarrow$ ADH secretion falls (negative feedback).

\textbf{Task 5.} Drinking \SI{4}{dm^{3}} of plain water rapidly: (a) the blood becomes too dilute — its \cle{water potential rises} above normal; (b) the osmoreceptors detect this, ADH secretion is \emph{inhibited}, the collecting ducts become \emph{less permeable}, and large volumes of dilute urine should be produced; (c) however, the water was absorbed faster than the kidneys could excrete it, so blood solutes (especially sodium ions) became dangerously diluted (\cle{hyponatraemia}); water then moves by osmosis into body cells, including brain cells, which swell — causing dizziness, nausea and confusion.

\textbf{Task 6.} The diagram should show three interlocking loops: (i) temperature $\rightarrow$ hypothalamus $\rightarrow$ skin (vasodilation, sweat); (ii) blood glucose $\rightarrow$ pancreas $\rightarrow$ liver (glycogenolysis); (iii) blood water potential $\rightarrow$ hypothalamus/pituitary $\rightarrow$ ADH $\rightarrow$ kidney collecting ducts; with arrows showing that each response cancels the initial change.

\begin{popfigure}
\centering
\begin{tikzpicture}[
    node distance=1.8cm and 1.2cm,
    var/.style={ellipse, draw=popdark, thick, fill=popblueL, text width=2.2cm, align=center, minimum height=1.2cm},
    rec/.style={rectangle, draw=popgreen, thick, fill=popgreenL, text width=2.2cm, align=center, minimum height=1.2cm},
    cc/.style={rectangle, draw=poppurple, thick, fill=poppurple!20, text width=2.2cm, align=center, minimum height=1.2cm},
    eff/.style={rectangle, draw=poporange, thick, fill=poporange!20, text width=2.2cm, align=center, minimum height=1.2cm},
    resp/.style={rectangle, draw=poppink, thick, fill=poppink!20, text width=2.2cm, align=center, minimum height=1.2cm},
    arrow/.style={-latex, thick},
    labelarrow/.style={-latex, thick, font=\small},
    looparrow/.style={-latex, thick, bend left=40, font=\small, color=popdark}
]
% Loop 1: Thermoregulation (top)
\node[var, align=center] (temp){Blood temperature\\ \SI{39.2}{\ensuremath{{}^{\circ}\mathrm{C}}} (high)};
\node[rec, below left=of temp, align=center] (thermoR){Receptors:\\ Thermoreceptors in skin \& hypothalamus};
\node[cc, below=of thermoR, align=center] (hypoT){Coordination centre:\\ Hypothalamus (heat-loss centre)};
\node[eff, below right=of hypoT, align=center] (skinE){Effectors:\\ Skin arterioles \& sweat glands};
\node[resp, right=of skinE, align=center] (skinR){Responses:\\ Vasodilation + sweating};

\draw[arrow] (temp) -- node[left,font=\small] {stimulus} (thermoR);
\draw[arrow] (thermoR) -- (hypoT);
\draw[arrow] (hypoT) -- (skinE);
\draw[arrow] (skinE) -- (skinR);
\draw[looparrow] (skinR) to node[above,font=\small] {negative feedback} (temp);

% Loop 2: Blood glucose (middle)
\node[var, right=3.5cm of temp, align=center] (gluc){Blood glucose\\ \SI{110}{mg/100cm^3} (rising)};
\node[rec, below left=of gluc, align=center] (pancreasR){Receptors:\\ $\alpha$-cells of pancreatic islets};
\node[cc, below=of pancreasR, align=center] (pancreasCC){Coordination centre:\\ Pancreas (endocrine)};
\node[eff, below right=of pancreasCC, align=center] (liverE){Effectors:\\ Liver hepatocytes};
\node[resp, right=of liverE, align=center] (liverR){Responses:\\ Glycogenolysis \& gluconeogenesis};

\draw[arrow] (gluc) -- node[left,font=\small] {stimulus} (pancreasR);
\draw[arrow] (pancreasR) -- (pancreasCC);
\draw[arrow] (pancreasCC) -- node[left,font=\small] {glucagon} (liverE);
\draw[arrow] (liverE) -- (liverR);
\draw[looparrow] (liverR) to node[above,font=\small] {negative feedback} (gluc);

% Loop 3: Osmoregulation (bottom)
\node[var, right=3.5cm of gluc, align=center] (water){Blood water potential\\ low (concentrated)};
\node[rec, below left=of water, align=center] (osmoR){Receptors:\\ Osmoreceptors in hypothalamus};
\node[cc, below=of osmoR, align=center] (pituitaryCC){Coordination centre:\\ Hypothalamus / Posterior pituitary};
\node[eff, below right=of pituitaryCC, align=center] (kidneyE){Effectors:\\ Collecting duct cells};
\node[resp, right=of kidneyE, align=center] (kidneyR){Responses:\\ Increased water permeability (ADH)};

\draw[arrow] (water) -- node[left,font=\small] {stimulus} (osmoR);
\draw[arrow] (osmoR) -- (pituitaryCC);
\draw[arrow] (pituitaryCC) -- node[left,font=\small] {ADH} (kidneyE);
\draw[arrow] (kidneyE) -- (kidneyR);
\draw[looparrow] (kidneyR) to node[above,font=\small] {negative feedback} (water);

% Cross-links: show integration (nervous/endocrine)
\draw[dashed, thick, popmuted] (hypoT) -- node[above,font=\tiny,sloped] {nervous} (pancreasCC);
\draw[dashed, thick, popmuted] (hypoT) -- (pituitaryCC);
\draw[dashed, thick, popmuted] (pancreasCC) -- (liverE);
\draw[dashed, thick, popmuted] (pituitaryCC) -- (kidneyE);
\end{tikzpicture}
\legende{Integrated homeostatic control during the Mount Cameroon race: three negative feedback loops for temperature, blood glucose, and blood water potential. Dashed lines indicate neural/endocrine coordination.}
\end{popfigure}

\textbf{Task 7.} Limitations: only one athlete was studied (no replication); measurements "during the race" were taken only at the summit, not continuously; sweat was collected only from a small forearm patch and may not represent whole-body loss. Additional measurement: blood sodium ion concentration, heart rate, or continuous core temperature.
\end{corrige}

\begin{attention}[Marking grid — GCE style]
A class correction grid modelled on GCE marking schemes: \textbf{Data handling} 3 marks; \textbf{Thermoregulation} 6 marks (receptor, centre, both effectors explained, humidity point); \textbf{Blood glucose} 5 marks; \textbf{Osmoregulation} 6 marks (data interpretation, ADH mechanism, loop); \textbf{Prediction} 4 marks; \textbf{Synthesis diagram} 4 marks (three complete, correctly labelled loops); \textbf{Evaluation} 2 marks. \textbf{Total: 30 marks.} Award marks for correct \emph{terminology} (vasodilation, glycogenolysis, osmoreceptor, collecting duct) and for explicitly stating that each loop is \emph{negative feedback}.
\end{attention}

\begin{attention}[Common errors to avoid]
\begin{itemize}
\item Saying that sweat itself "cools the blood" without mentioning \cle{evaporation} — it is evaporation that removes heat.
\item Writing that ADH "makes the kidney produce less urine" without explaining the \cle{permeability of the collecting duct} — the mechanism is required at A Level.
\item Confusing \cle{glucagon} with \cle{glycogen}, or claiming insulin lowers temperature — keep the hormones and their targets distinct.
\item Describing feedback without naming all four components: receptor, coordination centre, effector, response.
\end{itemize}
\end{attention}

\newpage

\coursec{Methods and worked examples}

\courssub{Skill 1 — Analysing a Negative Feedback Loop}

\begin{methode}[Identifying and describing a feedback mechanism]
To analyse any homeostatic control system in a GCE question, proceed as follows.
\begin{enumerate}
\item \textbf{Identify the controlled variable}: the factor kept constant (e.g. blood glucose, body temperature, blood water potential).
\item \textbf{Identify the set point (norm)}: the reference value around which the variable fluctuates (e.g. about $90\ \mathrm{mg}$ glucose per $100\ \mathrm{cm}^3$ blood; $37\ ^\circ\mathrm{C}$ core temperature).
\item \textbf{Identify the receptor (detector)}: the cells or organ that sense the deviation (e.g. hypothalamus, islets of Langerhans).
\item \textbf{Identify the coordinator (control centre)}: the structure that compares the input with the set point and sends instructions (e.g. hypothalamus, pancreas).
\item \textbf{Identify the effector(s)}: muscles or glands that carry out the corrective response.
\item \textbf{State the correction}: how the response returns the variable towards the set point.
\item \textbf{Classify the loop}: \cle{negative feedback} if the response cancels the deviation (the common case); \cle{positive feedback} if it amplifies it (rare, e.g. oxytocin in labour, blood clotting).
\end{enumerate}
\textbf{Reflexes:} always mention the \emph{stimulus $\rightarrow$ receptor $\rightarrow$ coordinator $\rightarrow$ effector $\rightarrow$ response $\rightarrow$ restoration} chain in full sentences; examiners award marks for each named component.
\end{methode}

\begin{exoresolu}[Analysing temperature control]
A student sits in a cold examination hall in Bamenda at $15\ ^\circ\mathrm{C}$. His core body temperature begins to fall. (a) Describe, as a feedback loop, how his body restores normal temperature. (b) State why this is described as negative feedback.
\tcblower
\textbf{(a) The feedback loop.}
\begin{itemize}
\item \emph{Stimulus:} core temperature falls below the set point of about $37\ ^\circ\mathrm{C}$.
\item \emph{Receptor:} thermoreceptors in the skin detect the cold, and the \cle{thermoregulatory centre in the hypothalamus} detects the fall in temperature of the blood flowing through it.
\item \emph{Coordinator:} the hypothalamus compares the input with the set point and sends nerve impulses to effectors.
\item \emph{Effectors and responses:}
\begin{itemize}
\item \textbf{Erector pili muscles} contract, raising hairs to trap an insulating layer of air (limited effect in humans);
\item \textbf{Arteriolar muscles} constrict (\cle{vasoconstriction}), reducing blood flow to the skin surface so less heat is lost by radiation;
\item \textbf{Skeletal muscles} contract rhythmically (\cle{shivering}), releasing heat from respiration;
\item \textbf{Sweat glands} stop secreting sweat, so no heat is lost by evaporation;
\item the \textbf{thyroid/adrenal medulla} may increase metabolic rate (thyroxine, adrenaline), increasing heat production.
\end{itemize}
\item \emph{Result:} heat production rises and heat loss falls, so core temperature returns to $37\ ^\circ\mathrm{C}$. The hypothalamus then reduces the corrective output.
\end{itemize}
\textbf{(b) Why negative feedback:} the response (warming) \emph{opposes and cancels} the initial change (cooling). The output of the system switches off its own stimulus, keeping the variable oscillating narrowly around the set point. This self-correcting reversal is the defining feature of negative feedback.
\end{exoresolu}

\courssub{Skill 2 — Interpreting Blood Glucose Regulation Data}

\begin{methode}[Explaining glucose homeostasis by the liver and pancreas]
\begin{enumerate}
\item \textbf{Know the two hormones and their sources}: \cle{insulin} from the $\beta$-cells and \cle{glucagon} from the $\alpha$-cells of the islets of Langerhans in the pancreas.
\item \textbf{High blood glucose (after a meal)} $\rightarrow$ insulin released $\rightarrow$ liver and muscles convert glucose to \cle{glycogen} (glycogenesis); glucose uptake and respiration by cells increase; fat synthesis increases $\rightarrow$ blood glucose falls.
\item \textbf{Low blood glucose (fasting, exercise)} $\rightarrow$ glucagon released $\rightarrow$ liver converts glycogen to glucose (\cle{glycogenolysis}) and makes glucose from amino acids/glycerol (\cle{gluconeogenesis}) $\rightarrow$ blood glucose rises.
\item \textbf{When interpreting a graph}: identify the meal time (glucose peak), then match the hormone curves: insulin rises \emph{after} glucose rises; glucagon rises when glucose falls. Hormone changes always \emph{lag slightly behind} the glucose change that triggers them.
\item \textbf{Conclude with the feedback statement}: each hormone's effect removes the stimulus for its own release — negative feedback.
\end{enumerate}
\textbf{Key idea to remember:} the liver is the \emph{effector}; the pancreas is both \emph{receptor} and \emph{coordinator}.
\end{methode}

\begin{exoresolu}[Glucose and hormone curves after a meal]
A healthy student eats a meal of achu and vegetables at 12:00. Her blood glucose rises from $90$ to $140\ \mathrm{mg}/100\ \mathrm{cm}^3$ by 12:45, then falls back to $90\ \mathrm{mg}/100\ \mathrm{cm}^3$ by 14:00. (a) Explain the changes in insulin and glucagon secretion over this period. (b) Predict what would happen to her blood glucose if her $\beta$-cells were destroyed.
\tcblower
\textbf{(a) Hormone changes.}
\begin{itemize}
\item \emph{12:00–12:45:} absorption of glucose from the ileum raises blood glucose above the set point. The $\beta$-cells of the islets of Langerhans detect this rise and secrete more \textbf{insulin}. Insulin stimulates the liver and muscles to convert glucose to glycogen, increases glucose uptake and respiration by cells, and promotes conversion of glucose to fat. Meanwhile the fall in glucagon secretion removes the stimulus for glycogen breakdown. Blood glucose therefore falls.
\item \emph{12:45–14:00:} as glucose falls back towards $90\ \mathrm{mg}/100\ \mathrm{cm}^3$, insulin secretion decreases (the stimulus has been removed — negative feedback). If glucose dipped below the set point, $\alpha$-cells would secrete \textbf{glucagon}, stimulating glycogenolysis and gluconeogenesis in the liver to restore the level.
\end{itemize}
\textbf{(b) Destroyed $\beta$-cells:} no insulin would be secreted. After the meal, blood glucose would rise and \emph{remain abnormally high} (hyperglycaemia); the liver would not store glycogen efficiently and cells would take up little glucose. Glucose would eventually appear in the urine once the kidney's reabsorption capacity was exceeded. This is the basis of \cle{diabetes mellitus} (Type 1). The feedback loop is broken because the effector signal (insulin) is missing.
\end{exoresolu}

\courssub{Skill 3 — Linking Skin Structure to Thermoregulatory Function}

\begin{methode}[Answering "structure–function" questions on the skin]
\begin{enumerate}
\item \textbf{Draw/recall the labelled section}: epidermis (cornified layer, granular layer, Malpighian layer, melanin), dermis (sweat glands and ducts, hair follicle, erector pili muscle, sebaceous gland, blood capillaries/arterioles, sensory receptors, nerve endings), subcutaneous \cle{adipose tissue}.
\item \textbf{Match each structure to a function}: protection (cornified layer, melanin against UV), temperature control (sweat glands, hairs, blood vessels, fat), sensitivity (receptors for touch, pressure, pain, temperature), excretion (small amounts of water, salts, urea in sweat), vitamin D synthesis, fat storage.
\item \textbf{For "hot conditions" answers}: vasodilation (more blood to surface $\rightarrow$ more radiation loss), sweating (evaporation absorbs latent heat), hairs lie flat, reduced metabolic heat.
\item \textbf{For "cold conditions" answers}: vasoconstriction, shivering, hairs raised, no sweating.
\item \textbf{Always state the mechanism of heat transfer}: radiation, conduction, convection, evaporation — and which structure affects which.
\end{enumerate}
\end{methode}

\begin{attention}[Common trap]
Sweating only cools the body if the sweat \emph{evaporates}; in humid weather evaporation is slow, which is why humid heat feels more uncomfortable. Also, vasoconstriction does not "move blood away from the skin" entirely — it reduces flow in superficial vessels; blood is diverted through deeper vessels.
\end{attention}

\begin{exoresolu}[Explaining temperature data from a practical]
In an experiment, a student ran for 10 minutes at midday in Douala (air temperature $33\ ^\circ\mathrm{C}$, high humidity). His skin temperature and sweat rate both increased, but his core temperature still rose by $1.2\ ^\circ\mathrm{C}$. Explain (a) the skin responses observed and (b) why core temperature nevertheless rose.
\tcblower
\textbf{(a) Skin responses.}
\begin{itemize}
\item The hypothalamus detected rising blood temperature and sent impulses to skin effectors.
\item \textbf{Vasodilation} of arterioles in the dermis increased blood flow near the surface, raising skin temperature and increasing heat loss by radiation and convection.
\item \textbf{Sweat glands} increased secretion; evaporation of sweat from the skin surface absorbs latent heat of vaporisation from the body, cooling it.
\item Hairs lay flat (erector pili relaxed) so no insulating air layer was trapped.
\end{itemize}
\textbf{(b) Why core temperature still rose.}
\begin{itemize}
\item Vigorous muscular contraction greatly increased \textbf{respiration in skeletal muscles}, releasing large amounts of heat — heat production temporarily exceeded heat loss.
\item The \textbf{high humidity} of Douala reduced the rate of evaporation of sweat (air already nearly saturated with water vapour), so the main cooling mechanism was inefficient.
\item The small air–skin temperature gradient ($33$ vs $37\ ^\circ\mathrm{C}$) reduced heat loss by radiation and convection.
\end{itemize}
Hence the negative feedback system was \emph{overloaded}: corrective responses were maximal but could not fully compensate, so core temperature rose until exercise stopped. This illustrates that homeostatic mechanisms have \emph{limits} beyond which the internal environment cannot be stabilised (risk of heat stroke).
\end{exoresolu}

\courssub{Skill 4 — Matching Excretory Products to Excretory Structures}

\begin{methode}[Building an excretion summary answer]
\begin{enumerate}
\item \textbf{Define excretion precisely}: the removal from the body of \emph{waste products of metabolism} (and substances in excess). Distinguish it from \emph{egestion} (removal of undigested food from the gut — never a product of metabolism).
\item \textbf{List the products and their origins}: carbon dioxide (respiration $\rightarrow$ lungs); urea (\cle{deamination} of excess amino acids in the liver $\rightarrow$ kidneys); water and mineral salts (kidneys, skin, lungs); bile pigments (breakdown of haemoglobin $\rightarrow$ liver $\rightarrow$ gut); heat (skin).
\item \textbf{Match each product to its organ}: lungs — \ce{CO2} and water vapour; kidneys — urea, water, salts (urine); skin — water, salts, traces of urea (sweat); liver — produces urea, excretes bile pigments.
\item \textbf{Know the nitrogenous products across animals}: \cle{ammonia} (very toxic, needs much water — aquatic animals), \cle{urea} (moderately toxic — mammals, amphibians), \cle{uric acid} (non-toxic, insoluble, conserves water — birds, insects, reptiles).
\end{enumerate}
\end{methode}

\begin{exoresolu}[Excretion across animal groups]
(a) Construct a table giving the main nitrogenous excretory product of a bony fish, a mammal and a bird, with one advantage of each product for that animal's way of life. (b) Explain why mammals convert ammonia to urea in the liver even though this costs energy.
\tcblower
\textbf{(a) Table.}
\begin{center}
\begin{tabularx}{\linewidth}{|l|l|X|}
\hline
\rowcolor{popblueL} \textbf{Animal} & \textbf{Product} & \textbf{Advantage} \\
\hline
Bony fish & Ammonia & Very soluble and quickly diluted/carried away by the surrounding water; no energy spent converting it; water is abundant so its toxicity is not a problem. \\
\hline
Mammal & Urea & Much less toxic than ammonia, so it can be carried in the blood and concentrated in urine; requires only moderate amounts of water for excretion. \\
\hline
Bird & Uric acid & Almost insoluble and non-toxic; excreted as a semi-solid paste, conserving water — vital for flight (low body mass) and for the embryo inside a shelled egg, where toxic waste would accumulate. \\
\hline
\end{tabularx}
\end{center}
\textbf{(b) Why mammals detoxify ammonia.} Ammonia is extremely toxic to cells even at low concentrations (it disrupts pH and enzyme function). A terrestrial mammal cannot dilute it in a surrounding water medium, and storing it would be lethal. In the liver, ammonia is combined with carbon dioxide in the \cle{ornithine cycle} (urea cycle) to form urea, which is about $10^5$ times less toxic, can be transported safely in the blood to the kidneys, and is excreted in a reasonably small volume of water. The energy cost is justified because the animal gains \emph{safe internal transport} and \emph{water economy} — both essential for life on land.
\end{exoresolu}

\courssub{Skill 5 — Calculating Filtration, Reabsorption and Excretion in the Nephron}

\begin{methode}[Quantitative treatment of urine formation]
\begin{enumerate}
\item \textbf{Recall the three processes and their sites}: \cle{ultrafiltration} (glomerulus/Bowman's capsule), \cle{selective reabsorption} (mainly proximal convoluted tubule and loop of Henlé), \cle{tubular secretion} (distal tubule), with final concentration in the collecting duct under \cle{ADH}.
\item \textbf{Use the mass balance}:
\[ \text{amount excreted} = \text{amount filtered} - \text{amount reabsorbed} + \text{amount secreted}. \]
\item \textbf{Percentage reabsorbed}:
\[ \%\ \text{reabsorbed} = \frac{\text{filtered} - \text{excreted}}{\text{filtered}} \times 100. \]
\item \textbf{Know the typical figures}: about $125\ \mathrm{cm}^3$ filtrate per minute ($\approx 180\ \mathrm{dm}^3$ per day) but only about $1.5\ \mathrm{dm}^3$ urine per day — over $99\%$ of the water is reabsorbed. Glucose is \emph{completely} reabsorbed in a healthy person.
\item \textbf{Interpret composition tables}: a substance \emph{more concentrated} in urine than in filtrate (e.g. urea) has had water reabsorbed around it; a substance \emph{absent} from urine (glucose, proteins) has been fully reabsorbed or never filtered.
\end{enumerate}
\end{methode}

\begin{exoresolu}[Glomerular filtration calculation]
A patient's kidneys produce $170\ \mathrm{dm}^3$ of glomerular filtrate per day. The filtrate contains $1.0\ \mathrm{g\,dm^{-3}}$ of glucose and $0.3\ \mathrm{g\,dm^{-3}}$ of urea. The daily urine output is $1.4\ \mathrm{dm}^3$, containing no glucose and $20\ \mathrm{g\,dm^{-3}}$ of urea. Calculate (a) the mass of glucose filtered per day and the percentage reabsorbed; (b) the mass of urea filtered and excreted per day, and the percentage of filtered urea reabsorbed; (c) the percentage of filtered water reabsorbed.
\tcblower
\textbf{(a) Glucose.}
\begin{align*}
\text{Mass filtered} &= 170\ \mathrm{dm}^3 \times 1.0\ \mathrm{g\,dm^{-3}} = 170\ \mathrm{g/day}.\\
\text{Mass excreted} &= 0 \quad (\text{no glucose in urine}).\\
\%\ \text{reabsorbed} &= \frac{170 - 0}{170} \times 100 = 100\ \%.
\end{align*}
All filtered glucose is reabsorbed in the proximal convoluted tubule by active transport — normal for a healthy kidney.

\textbf{(b) Urea.}
\begin{align*}
\text{Mass filtered} &= 170 \times 0.3 = 51\ \mathrm{g/day}.\\
\text{Mass excreted} &= 1.4\ \mathrm{dm}^3 \times 20\ \mathrm{g\,dm^{-3}} = 28\ \mathrm{g/day}.\\
\text{Mass reabsorbed} &= 51 - 28 = 23\ \mathrm{g/day}.\\
\%\ \text{reabsorbed} &= \frac{23}{51} \times 100 \approx 45\ \%.
\end{align*}
Only about $45\%$ of urea is reabsorbed (passively, with water); the rest is excreted. Note the urine urea concentration ($20\ \mathrm{g\,dm^{-3}}$) is about $67$ times the filtrate concentration, because water was reabsorbed while much urea stayed in the tubule.

\textbf{(c) Water.}
\[ \%\ \text{water reabsorbed} = \frac{170 - 1.4}{170} \times 100 = \frac{168.6}{170} \times 100 \approx 99.2\ \%. \]
\textbf{Interpretation:} the kidney is fundamentally a \emph{reabsorption} organ: a huge volume is filtered, then almost all water and all useful solutes (glucose, most salts) are reclaimed, while wastes are left behind in a small, concentrated volume of urine.
\end{exoresolu}

\courssub{Skill 6 — Explaining Osmoregulation and ADH Action}

\begin{methode}[Answering ADH/osmoregulation questions]
\begin{enumerate}
\item \textbf{State the detector}: \cle{osmoreceptors} in the hypothalamus detect changes in the water potential (solute concentration) of the blood.
\item \textbf{State the effector pathway}: the posterior \cle{pituitary gland} releases \cle{ADH} (antidiuretic hormone); ADH acts on the \emph{collecting ducts} (and distal tubules) of the kidneys, increasing their \textbf{permeability to water}.
\item \textbf{Dehydration scenario}: blood too concentrated $\rightarrow$ more ADH $\rightarrow$ more water reabsorbed $\rightarrow$ small volume of \emph{concentrated} urine $\rightarrow$ blood water potential restored.
\item \textbf{Overhydration scenario}: blood too dilute $\rightarrow$ less ADH $\rightarrow$ collecting ducts impermeable $\rightarrow$ large volume of \emph{dilute} urine.
\item \textbf{Close the loop}: the response removes the stimulus for ADH release — negative feedback.
\item \textbf{For adaptation questions}, link structure to environment: desert mammals (e.g. kangaroo rat) have very \emph{long loops of Henlé} $\rightarrow$ steep medullary salt gradient $\rightarrow$ very concentrated urine, conserving water; freshwater fish excrete copious dilute urine and actively take in salts at the gills; marine bony fish drink sea water and excrete salts via the gills.
\end{enumerate}
\end{methode}

\begin{exoresolu}[ADH and a desert mammal]
(a) A student drinks $2\ \mathrm{dm}^3$ of water after a football match. Describe how her kidneys respond over the next two hours. (b) The desert rodent of the Sahel produces urine about four times more concentrated than that of humans. Suggest two structural or physiological adaptations that make this possible.
\tcblower
\textbf{(a) Response to excess water.}
\begin{itemize}
\item The large water intake is absorbed from the gut, so the blood becomes \emph{more dilute} (its water potential rises / solute concentration falls).
\item \textbf{Osmoreceptors} in the hypothalamus detect the dilution and reduce stimulation of the posterior pituitary.
\item \textbf{Less ADH} is released into the blood.
\item The walls of the \textbf{collecting ducts} become \emph{less permeable} to water, so most of the water in the filtrate passing through them is \emph{not} reabsorbed.
\item A \textbf{large volume of dilute urine} is produced over the next hours, removing the excess water.
\item As blood concentration returns to the set point, ADH secretion returns to normal — negative feedback.
\end{itemize}
\textbf{(b) Desert rodent adaptations.}
\begin{itemize}
\item \textbf{Very long loops of Henlé} (and a relatively thick kidney medulla): the counter-current multiplier establishes a very steep sodium chloride gradient deep in the medulla, so water can be withdrawn from the collecting duct against a very low external water potential, yielding extremely concentrated urine.
\item \textbf{High sensitivity of collecting ducts to ADH / sustained high ADH levels}, maximising water reabsorption; in addition, such rodents produce very dry faeces and obtain most water from \emph{metabolic water} (oxidation of food), and many are nocturnal, avoiding evaporative loss.
\end{itemize}
Each feature reduces water loss, allowing the animal to maintain blood water potential in an environment where free water is scarce.
\end{exoresolu}

\courssub{Skill 7 — Planning an Integrated Essay (Liver, Kidney, Skin, Plants)}

\begin{methode}[Writing a high-mark synoptic homeostasis essay]
\begin{enumerate}
\item \textbf{Define the key term} in the first sentence (homeostasis: maintenance of a constant internal environment despite external changes).
\item \textbf{Plan one paragraph per organ}, each following the pattern: \emph{what is regulated $\rightarrow$ how it is detected $\rightarrow$ what the organ does $\rightarrow$ the feedback outcome}.
\item \textbf{Liver}: carbohydrate metabolism (glycogenesis, glycogenolysis, gluconeogenesis), deamination and urea formation, detoxification, bile production, heat production, storage (iron, vitamins, glycogen), breakdown of old red blood cells.
\item \textbf{Kidneys}: osmoregulation, excretion of urea, salt balance, pH regulation.
\item \textbf{Skin}: temperature regulation, minor excretion, protection.
\item \textbf{Plants}: no special excretory organs; wastes handled by \emph{storage in vacuoles, bark, leaves and fruits} (then shed), \emph{gaseous exchange} of \ce{O2}/\ce{CO2} through stomata and lenticels, exudation of resins/gums/latex, and conversion of excess nitrogen into \emph{protein stored in seeds}; plants excrete little because they reuse \ce{CO2} and \ce{H2O} in photosynthesis and make proteins only as needed.
\item \textbf{Conclude} by naming the nervous and endocrine systems as the \emph{integrators} coordinating all organs.
\end{enumerate}
\end{methode}

\begin{attention}[Examiner's warning]
Do not describe organs in isolation without linking them to the \emph{constancy of the internal environment} — that link is what earns the top band. And never say plants "excrete through roots into the soil" as their main strategy; the syllabus emphasises storage and shedding.
\end{attention}

\begin{exoresolu}[Synoptic GCE essay]
"Describe how the liver, kidneys and skin contribute to homeostasis in humans, and explain how excretion in flowering plants differs from excretion in animals." (GCE-style, 20 marks) — give a model plan and the key marking points.
\tcblower
\textbf{Introduction (2 marks).} Homeostasis is the maintenance of a constant internal environment (temperature, blood glucose, water potential, pH) within narrow limits, achieved by negative feedback involving receptors, coordinators and effectors.

\textbf{Liver (5 marks).} Regulates blood glucose: stores glucose as glycogen under insulin, releases glucose by glycogenolysis and gluconeogenesis under glucagon. Carries out \textbf{deamination} of excess amino acids, producing ammonia which is converted to \textbf{urea} (ornithine cycle) for safe transport to the kidneys. Detoxifies drugs, alcohol and poisons. Produces heat from its intense metabolism, distributed by the blood. Stores glycogen, iron and vitamins; breaks down worn-out red blood cells, excreting bile pigments.

\textbf{Kidneys (5 marks).} Ultrafiltration at the glomerulus; selective reabsorption of all glucose, most water and salts; urea, excess salts and water excreted as urine. Osmoregulation via ADH: concentrated blood $\rightarrow$ more ADH $\rightarrow$ permeable collecting ducts $\rightarrow$ concentrated urine, and the reverse when blood is dilute. Also regulate blood pH and salt balance.

\textbf{Skin (4 marks).} Thermoregulation: sweating and evaporation, vasodilation when hot; vasoconstriction, shivering (with muscles) and raised hairs when cold; subcutaneous fat insulates. Minor excretion of water, salts and urea in sweat. Cornified layer and melanin protect against entry of pathogens and UV damage.

\textbf{Plants vs animals (3 marks).} Plants have \emph{no specialised excretory organs}: they produce little waste because \ce{CO2} and \ce{H2O} from respiration are reused in photosynthesis, and protein synthesis is matched to need, so little nitrogenous waste forms. Gaseous wastes leave through stomata and lenticels; other wastes (tannins, alkaloids, calcium oxalate crystals, resins, gums, latex) are stored in vacuoles, old leaves, bark, flowers and fruits and removed when these are shed; some are exuded (e.g. gum arabic). Animals, by contrast, cannot reuse their nitrogenous waste and need dedicated organs (kidneys, lungs, skin).

\textbf{Conclusion (1 mark).} The nervous and endocrine systems integrate these organs into a single coordinated homeostatic system, keeping the internal environment stable.
\end{exoresolu}

\newpage

\coursec{Exercises}

\courssub{I check my knowledge}

\begin{exercice}[Multiple choice questions]
For each question, choose the single correct answer.
\begin{enumerate}
\item Homeostasis is best defined as:
\begin{enumerate}
\item the production of heat by the body;
\item the maintenance of a constant internal environment despite external changes;
\item the elimination of nitrogenous waste;
\item the control of breathing rate.
\end{enumerate}
\item The hormone that lowers blood glucose concentration is:
\begin{enumerate}
\item glucagon; \item adrenaline; \item insulin; \item ADH.
\end{enumerate}
\item In a negative feedback loop, the effector:
\begin{enumerate}
\item detects the change in the controlled factor;
\item amplifies the initial change;
\item acts to reverse the change detected by the receptor;
\item stores the controlled factor.
\end{enumerate}
\item The part of the nephron where selective reabsorption of glucose is completed is:
\begin{enumerate}
\item the loop of Henl\'e; \item the proximal convoluted tubule; \item the collecting duct; \item the Bowman's capsule.
\end{enumerate}
\item The main nitrogenous excretory product of a bony fish such as tilapia is:
\begin{enumerate}
\item urea; \item uric acid; \item ammonia; \item creatinine.
\end{enumerate}
\end{enumerate}
\end{exercice}

\begin{exercice}[True or false — justify your answer]
State whether each statement is true or false, and justify your answer in one or two sentences.
\begin{enumerate}
\item Urea is produced in the liver from excess amino acids.
\item Sweating cools the body because sweat is colder than the skin.
\item ADH increases the permeability of the collecting duct to water.
\item Glucose is normally found in the urine of a healthy person.
\item Plants excrete large quantities of nitrogenous waste like mammals do.
\end{enumerate}
\end{exercice}

\begin{exercice}[Key definitions]
Define each of the following terms precisely, as would be expected at the GCE Advanced Level:
\begin{enumerate}
\item homeostasis;
\item osmoregulation;
\item deamination;
\item ultrafiltration;
\item negative feedback.
\end{enumerate}
\end{exercice}

\begin{exercice}[Direct calculations]
\begin{enumerate}
\item A patient produces $180\ \mathrm{dm^3}$ of glomerular filtrate per day and excretes $1.5\ \mathrm{dm^3}$ of urine per day. Calculate the volume of filtrate reabsorbed per day and express it as a percentage of the filtrate.
\item The glomerular filtrate contains $1.0\ \mathrm{g\,dm^{-3}}$ of glucose. Assuming complete reabsorption, calculate the mass of glucose reabsorbed per day from $180\ \mathrm{dm^3}$ of filtrate.
\item A man loses $500\ \mathrm{cm^3}$ of water in urine, $400\ \mathrm{cm^3}$ in sweat and $350\ \mathrm{cm^3}$ in expired air per day, and gains $250\ \mathrm{cm^3}$ from metabolic water. Calculate the volume of water he must obtain from food and drink to remain in water balance.
\end{enumerate}
\end{exercice}

\courssub{I practise}

\begin{exercice}[Classifying animals by excretory product]
Copy and complete the table below for the following animals: \emph{Amoeba}, tilapia (freshwater bony fish), toad, lizard, man, locust. For each animal, state its main nitrogenous excretory product and explain, in terms of water availability and toxicity, how the product is adapted to the animal's habitat.
\begin{center}
\begin{tabularx}{\linewidth}{|l|X|X|}
\hline
\rowcolor{popblueL} \textbf{Animal} & \textbf{Main excretory product} & \textbf{Link to habitat} \\
\hline
Amoeba & & \\
\hline
Tilapia & & \\
\hline
Toad & & \\
\hline
Lizard & & \\
\hline
Man & & \\
\hline
Locust & & \\
\hline
\end{tabularx}
\end{center}
\end{exercice}

\begin{exercice}[Analysing plasma, filtrate and urine]
The table below shows the concentrations of some substances in blood plasma, glomerular filtrate and urine of a healthy mammal (in $\mathrm{g\,dm^{-3}}$).
\begin{center}
\begin{tabularx}{\linewidth}{|l|X|X|X|}
\hline
\rowcolor{popblueL} \textbf{Substance} & \textbf{Plasma} & \textbf{Filtrate} & \textbf{Urine} \\
\hline
Proteins & 80 & 0 & 0 \\
\hline
Glucose & 1.0 & 1.0 & 0 \\
\hline
Urea & 0.3 & 0.3 & 20 \\
\hline
Sodium ions & 3.2 & 3.2 & 6.0 \\
\hline
\end{tabularx}
\end{center}
\begin{enumerate}
\item Explain why proteins are absent from the glomerular filtrate.
\item Identify the substance(s) that are completely reabsorbed, and state precisely where and how this reabsorption occurs in the nephron.
\item Explain why the concentration of urea is much higher in urine than in the filtrate.
\item Suggest one reason why the sodium ion concentration differs between filtrate and urine.
\end{enumerate}
\end{exercice}

\begin{exercice}[Blood glucose and insulin after a meal]
A healthy student eats a meal rich in carbohydrates at 12:00. Blood glucose and blood insulin concentrations are measured every 30 minutes for three hours.
\begin{enumerate}
\item Sketch, on the same axes, the curves you expect for blood glucose and blood insulin concentration against time. Label the axes and the curves.
\item Describe the changes in each curve after the meal.
\item Explain, step by step, the negative feedback mechanism that returns blood glucose to its normal level, naming the receptor organ, the endocrine gland, the hormone and the target organ involved, and stating precisely the role of the liver.
\item Predict what would happen to the two curves if the meal were skipped and the student then did vigorous exercise.
\end{enumerate}
\end{exercice}

\begin{exercice}[Investigating catalase in the liver]
A Form Upper Sixth student in Buea wants to show that the liver contains an enzyme that breaks down hydrogen peroxide.
\begin{enumerate}
\item Describe a simple experiment to demonstrate this, using pieces of fresh liver and hydrogen peroxide solution. State how the gas evolved can be identified.
\item State two control experiments you would set up and explain why each is necessary.
\item Predict the results you would obtain if the liver were boiled before the experiment, and explain your prediction.
\item State two precautions needed to obtain reliable results.
\end{enumerate}
\end{exercice}

\courssub{I prepare for the GCE}

\begin{exercice}[Structured question: urine output and its disorders — 20 marks]
A patient at a district hospital in Bamenda complains of excessive thirst and produces about $5\ \mathrm{dm^3}$ of dilute urine per day. Tests show that his urine contains glucose.
\begin{enumerate}
\item Describe how ultrafiltration in the Bowman's capsule produces glomerular filtrate, and explain why the filtrate of a healthy person contains glucose but no protein. \hfill (6 marks)
\item Explain how glucose is normally reabsorbed in the nephron, naming the region and the transport mechanism involved. \hfill (4 marks)
\item The patient's condition is diagnosed as diabetes mellitus. Explain how a failure of insulin production leads to the presence of glucose in the urine and to the production of large volumes of dilute urine. \hfill (6 marks)
\item A second patient produces large volumes of dilute urine but has no glucose in it and normal blood insulin. Suggest a hormonal cause of this condition and justify your answer by describing the normal role of the hormone concerned. \hfill (4 marks)
\end{enumerate}
\end{exercice}

\begin{exercice}[Essay: response to a fall in core temperature — 20 marks]
Write an essay entitled: ``Describe how the mammalian body responds to a fall in core body temperature.''

Your essay should include:
\begin{enumerate}
\item the role of the hypothalamus as receptor and coordination centre;
\item the structure of the skin relevant to temperature control (with a clearly labelled diagram);
\item the responses of at least four effectors (blood vessels of the skin, sweat glands, erector pili muscles, skeletal muscles, and the role of hormones such as thyroxine/adrenaline);
\item an explanation of how these responses form a negative feedback loop that restores the normal temperature.
\end{enumerate}
(Marks: structure and receptors 4; effector responses 10; feedback loop 4; quality of diagram and organisation 2.)
\end{exercice}

\begin{exercice}[Data analysis and comparison: osmoregulation in fishes and excretion in plants — 20 marks]
\textbf{Part A.} A freshwater bony fish (tilapia) produces large volumes of very dilute urine, while a marine bony fish produces small volumes of urine and actively excretes salts.
\begin{enumerate}
\item Using the principle of osmosis, explain why the freshwater fish tends to gain water and lose salts, and how its kidneys and gills correct this. \hfill (5 marks)
\item Explain the opposite osmotic problem faced by the marine fish and how it is solved. \hfill (5 marks)
\end{enumerate}
\textbf{Part B.} ``Plants excrete differently from animals.''
\begin{enumerate}
\item Discuss this statement with reference to the nature of plant excretory products (oxygen, carbon dioxide, water, tannins, resins, latex, alkaloids, crystals) and the reasons why plants have no specialised excretory organs. \hfill (6 marks)
\item Give two examples of plant excretory or stored waste products of economic importance in Cameroon and state their uses. \hfill (4 marks)
\end{enumerate}
\end{exercice}

\newpage

\coursec{Solutions to the exercises}

\begin{corrige}[Exercise 1 — Multiple choice questions]
\begin{enumerate}
\item \textbf{Answer: (b).} Homeostasis is the maintenance of a relatively constant internal environment (temperature, blood glucose, water potential, pH, etc.) despite changes in the external environment. Options (a), (c) and (d) describe only single processes, not the general principle.
\item \textbf{Answer: (c).} Insulin, secreted by the $\beta$-cells of the islets of Langerhans in the pancreas, lowers blood glucose by stimulating its uptake by cells and its conversion to glycogen in the liver and muscles. Glucagon and adrenaline \emph{raise} blood glucose; ADH acts on water balance, not glucose.
\item \textbf{Answer: (c).} In negative feedback the effector produces a response that \emph{reverses} (counteracts) the change detected by the receptor, bringing the controlled factor back towards its set point. Detection is the role of the receptor (a), and amplification (b) is characteristic of \emph{positive} feedback.
\item \textbf{Answer: (b).} All the glucose in the filtrate is reabsorbed in the \cle{proximal convoluted tubule}, by active transport (via sodium-glucose co-transporters) across the microvilli-bearing epithelial cells.
\item \textbf{Answer: (c).} Freshwater bony fishes such as tilapia excrete \cle{ammonia}, which is very toxic but highly soluble; it is diluted and washed away rapidly by the abundant water of their habitat, mainly across the gills.
\end{enumerate}
\end{corrige}

\begin{corrige}[Exercise 2 — True or false]
\begin{enumerate}
\item \textbf{True.} Excess amino acids cannot be stored; in the liver they undergo \cle{deamination} (removal of the amino group), and the ammonia formed is converted into the less toxic urea in the ornithine cycle.
\item \textbf{False.} Sweat is not colder than the skin; cooling occurs because the \emph{evaporation} of sweat from the skin surface removes latent heat of vaporisation from the body.
\item \textbf{True.} ADH (antidiuretic hormone) binds to the cells of the distal convoluted tubule and especially the collecting duct, increasing their permeability to water so that more water is reabsorbed into the blood, producing a smaller volume of concentrated urine.
\item \textbf{False.} In a healthy person all the glucose filtered at the glomerulus is reabsorbed in the proximal convoluted tubule, so urine is normally glucose-free. Glucose appears in urine only when blood glucose exceeds the renal threshold (e.g. in diabetes mellitus).
\item \textbf{False.} Plants produce very little nitrogenous waste because their protein metabolism is limited and they reuse nitrogen compounds; their main excretory products are gases ($\ce{O2}$, $\ce{CO2}$), water vapour and organic substances such as tannins, resins and alkaloids, often stored in dead tissues or lost in fallen leaves.
\end{enumerate}
\end{corrige}

\begin{corrige}[Exercise 3 — Key definitions]
\begin{enumerate}
\item \textbf{Homeostasis:} the maintenance of a relatively constant internal environment (physical and chemical factors such as temperature, blood glucose concentration, water potential and pH of the blood and tissue fluid) despite fluctuations in the external environment, by means of self-regulating feedback mechanisms.
\item \textbf{Osmoregulation:} the control of the water potential (and solute concentration) of the body fluids, i.e. the regulation of the balance between water and solutes gained and lost by the organism, so that cells are neither shrunken nor swollen by osmosis.
\item \textbf{Deamination:} the removal of the amino group ($-\ce{NH2}$) from an amino acid in the liver, producing ammonia (which is converted to urea) and a keto acid that can be respired or converted to carbohydrate or fat.
\item \textbf{Ultrafiltration:} the filtration of blood under high pressure in the glomerulus, forcing water and small solutes (glucose, urea, salts, amino acids) from the capillaries into the Bowman's capsule, while retaining blood cells and plasma proteins in the blood.
\item \textbf{Negative feedback:} a control mechanism in which a change in a controlled factor is detected by a receptor, and the response of the effector counteracts (reverses) the initial change, returning the factor towards its set point and thereby switching off the response.
\end{enumerate}
\end{corrige}

\begin{corrige}[Exercise 4 — Direct calculations]
\begin{enumerate}
\item Volume reabsorbed $=$ filtrate $-$ urine $= 180 - 1.5 = 178.5\ \mathrm{dm^3}$ per day.
\[ \text{Percentage reabsorbed} = \frac{178.5}{180} \times 100 = 99.17\ \% \approx 99.2\ \%. \]
This illustrates the enormous efficiency of tubular reabsorption: only about $0.8\ \%$ of the filtrate leaves the body as urine.
\item Mass of glucose reabsorbed $=$ concentration $\times$ volume:
\[ m = 1.0\ \mathrm{g\,dm^{-3}} \times 180\ \mathrm{dm^3} = 180\ \mathrm{g} \text{ of glucose per day.} \]
All of this is reabsorbed in the proximal convoluted tubule, which is why none appears in the urine.
\item Typical daily water losses: urine $500\ \mathrm{cm^3}$, sweat $400\ \mathrm{cm^3}$, faeces and expired air $350\ \mathrm{cm^3}$.
Total water loss $= 500 + 400 + 350 = 1250\ \mathrm{cm^3}$ per day.

For water balance, total input $=$ total output:
\[ \text{water from food and drink} + \text{metabolic water} = 1250 \]
\[ \text{water from food and drink} = 1250 - 250 = 1000\ \mathrm{cm^3} = 1.0\ \mathrm{dm^3}. \]
\end{enumerate}
\end{corrige}

\begin{corrige}[Exercise 5 — Classifying animals by excretory product]
\begin{center}
\begin{tabularx}{\linewidth}{|l|X|X|}
\hline
\rowcolor{popblueL} \textbf{Animal} & \textbf{Main excretory product} & \textbf{Link to habitat} \\
\hline
Amoeba & Ammonia & Lives in fresh water; ammonia is very toxic but highly soluble and is immediately diluted and washed away by the surrounding water, diffusing out across the cell surface. \\
\hline
Tilapia & Ammonia & Freshwater fish with abundant water available; ammonia diffuses out across the gills and is diluted at once, so there is no need to spend energy converting it. \\
\hline
Toad & Urea & Amphibian living partly on land where water is less abundant; ammonia is converted in the liver to the less toxic urea, which can be retained in the blood at low concentration and excreted with moderate water loss. \\
\hline
Lizard & Uric acid & Terrestrial reptile, often in dry habitats; uric acid is almost insoluble and non-toxic, so it is excreted as a semi-solid paste with almost no water loss — essential for water conservation. \\
\hline
Man & Urea & Terrestrial mammal; urea is far less toxic than ammonia and can be transported in the blood to the kidneys and eliminated in a concentrated urine, conserving water while avoiding toxicity. \\
\hline
Locust & Uric acid & Terrestrial flying insect; excreting semi-solid uric acid via the Malpighian tubules minimises water loss and body mass — adaptations to dry land life and flight. \\
\hline
\end{tabularx}
\end{center}
\textbf{General principle:} the more toxic the product, the more water is needed to eliminate it safely. Ammonia $\rightarrow$ aquatic animals with unlimited water; urea $\rightarrow$ animals with moderate water availability; uric acid $\rightarrow$ terrestrial animals that must conserve water.
\end{corrige}

\begin{corrige}[Exercise 6 — Analysing plasma, filtrate and urine]
\begin{enumerate}
\item \textbf{Proteins absent from filtrate:} plasma proteins are too large (relative molecular mass above about 68\,000) to pass through the basement membrane of the glomerular capillaries, which acts as a selective filter during ultrafiltration. Only water and small solutes are forced through; proteins and blood cells remain in the blood.
\item \textbf{Completely reabsorbed substance: glucose.} Its concentration falls from $1.0\ \mathrm{g\,dm^{-3}}$ in the filtrate to zero in urine. It is reabsorbed entirely in the \cle{proximal convoluted tubule} by \textbf{active transport} (sodium-glucose co-transport) across the epithelial cells, which have microvilli (large surface area) and many mitochondria (ATP supply).
\item \textbf{Urea more concentrated in urine:} urea is filtered but only partially reabsorbed, while a very large proportion of the \emph{water} of the filtrate (about $99\ \%$) is reabsorbed along the tubule and collecting duct. The same mass of urea is therefore dissolved in a much smaller volume, so its concentration rises from $0.3$ to $20\ \mathrm{g\,dm^{-3}}$.
\item \textbf{Sodium ions:} sodium is actively reabsorbed (proximal tubule, loop of Henl\'e, distal tubule), but the \emph{amount} reabsorbed is regulated according to the body's salt balance (under hormonal control, e.g. aldosterone). Since water reabsorption and salt reabsorption are regulated independently, the final urine can be more concentrated in $\ce{Na+}$ than the filtrate while still excreting excess salt.
\end{enumerate}
\end{corrige}

\begin{corrige}[Exercise 7 — Blood glucose and insulin after a meal]
\begin{enumerate}
\item Sketch of the expected curves:
\begin{popfigure}
\begin{tikzpicture}[scale=0.95]
\draw[->, thick] (0,0) -- (7.2,0) node[below left, align=center, text width=2.2cm] {time after meal (h)};
\draw[->, thick] (0,0) -- (0,4.6) node[above left, align=center, text width=3.2cm] {blood concentration (a.u.)};
\foreach \x/\t in {0/0, 1/0.5, 2/1.0, 3/1.5, 4/2.0, 5/2.5, 6/3.0} \draw (\x,0.05) -- (\x,-0.05) node[below, font=\small] {\t};
\draw[popblue, very thick, smooth] plot coordinates {(0,1.6) (0.5,2.1) (1,3.0) (1.5,2.6) (2,2.0) (2.5,1.7) (3,1.6) (6,1.6)};
\draw[poporange, very thick, smooth] plot coordinates {(0,1.0) (0.5,1.4) (1,2.3) (1.5,2.0) (2,1.4) (2.5,1.1) (3,1.0) (6,1.0)};
\node[popblue, font=\small] at (4.6,2.9) {blood glucose};
\node[poporange, font=\small] at (4.6,1.9) {blood insulin};
\draw[dashed, gray] (0,1.6) -- (6.5,1.6) node[pos=0, left, font=\small, black] {normal};
\end{tikzpicture}
\legende{Blood glucose and blood insulin after a carbohydrate-rich meal.}
\end{popfigure}
\item \textbf{Description:} after the meal, digestion and absorption of carbohydrates cause blood glucose to rise, peaking about 30--60 minutes after eating. The rise in glucose stimulates insulin secretion, so the insulin curve rises shortly after (and roughly parallel to) the glucose curve. As insulin acts, blood glucose falls back to the normal (set-point) level within about 2--3 hours, and insulin secretion then falls back to its basal level too.
\item \textbf{Step-by-step negative feedback:}
\begin{itemize}
\item \textbf{Stimulus:} blood glucose rises above the set point after absorption of the meal.
\item \textbf{Receptor:} the $\beta$-cells of the \cle{islets of Langerhans} in the \textbf{pancreas} detect the rise (the pancreas acts as both receptor and endocrine gland).
\item \textbf{Hormone:} the $\beta$-cells secrete \cle{insulin} into the blood.
\item \textbf{Target organs and role of the liver:} insulin acts on the \textbf{liver} and muscles, stimulating them to take up glucose and convert it to \textbf{glycogen} (glycogenesis); it also increases glucose uptake and respiration by cells and inhibits the breakdown of glycogen and the formation of glucose from other sources (gluconeogenesis).
\item \textbf{Result and feedback:} blood glucose falls back to the set point; the fall is detected by the pancreas, which reduces insulin secretion — the response switches itself off, which is the essence of negative feedback.
\end{itemize}
\item \textbf{Skipped meal + vigorous exercise:} respiring muscles use large amounts of glucose, so \textbf{blood glucose would fall below normal}. The $\alpha$-cells of the pancreas then secrete \textbf{glucagon} (and adrenaline is released), stimulating the liver to convert glycogen back to glucose (glycogenolysis) and to produce glucose from amino acids. Blood glucose rises back to normal. The \textbf{insulin curve would fall below its basal level}, since low blood glucose suppresses insulin secretion. Thus both curves move in the opposite direction to part (a).
\end{enumerate}
\end{corrige}

\begin{corrige}[Exercise 8 — Investigating catalase in the liver]
\begin{enumerate}
\item \textbf{Experiment:} cut two equal small cubes of fresh raw liver. Place one cube in a test tube and add about $5\ \mathrm{cm^3}$ of hydrogen peroxide solution. Effervescence occurs as the enzyme \cle{catalase} in the liver tissue catalyses the decomposition:
\[ 2\ce{H2O2} \xrightarrow{\text{catalase}} 2\ce{H2O} + \ce{O2}. \]
\textbf{Identification of the gas:} insert a \textbf{glowing splint} into the mouth of the tube; it \textbf{rekindles (relights)}, showing the gas is oxygen.
\item \textbf{Controls:}
\begin{itemize}
\item \emph{Control 1 — hydrogen peroxide alone, no liver:} shows that the peroxide does not decompose significantly by itself at room temperature, so the gas produced in the test comes from the action of the liver.
\item \emph{Control 2 — boiled (and cooled) liver + hydrogen peroxide:} shows that the activity is due to a heat-sensitive component (an enzyme), not to a non-living chemical property of the tissue.
\end{itemize}
\item \textbf{Boiled liver:} little or no effervescence occurs; the glowing splint is not relit (or only very weakly). Boiling \textbf{denatures} catalase: the high temperature destroys the tertiary structure of the enzyme and the shape of its active site, so it can no longer bind hydrogen peroxide. Enzyme denaturation by heat is irreversible.
\item \textbf{Precautions (any two):}
\begin{itemize}
\item use equal masses/sizes of liver pieces and equal volumes and concentrations of hydrogen peroxide in all tubes (fair test);
\item use fresh liver and grind or cut it to expose a comparable surface area;
\item keep all tubes at the same (room) temperature;
\item repeat the experiment and compare results for reliability;
\item allow boiled liver to cool before adding peroxide (otherwise heat alone could decompose the peroxide).
\end{itemize}
\end{enumerate}
\end{corrige}

\begin{corrige}[Exercise 9 — Structured question: urine output and its disorders]
\textbf{(a) Ultrafiltration (6 marks) — indicative mark scheme:}
\begin{itemize}
\item Blood enters the glomerulus via the afferent arteriole, which is wider than the efferent arteriole, so a \textbf{high hydrostatic pressure} builds up in the glomerular capillaries (1).
\item This pressure forces water and small solutes — glucose, urea, salts, amino acids — out of the blood through the capillary walls and the \textbf{basement membrane} into the Bowman's capsule (1); the podocytes and basement membrane act as a selective filter (1).
\item \textbf{Glucose present:} glucose molecules are small enough to pass through the filter, so filtrate contains glucose at the same concentration as plasma (1).
\item \textbf{Protein absent:} plasma proteins are too large to pass through the basement membrane (1) and are retained in the blood together with the blood cells (1).
\end{itemize}
\textbf{(b) Glucose reabsorption (4 marks):}
\begin{itemize}
\item All glucose is reabsorbed in the \textbf{proximal convoluted tubule} (1).
\item The mechanism is \textbf{active transport} against the concentration gradient (1), via sodium-glucose co-transport (1).
\item The tubule cells are adapted by microvilli (large surface area) and numerous mitochondria supplying ATP (1).
\end{itemize}
\textbf{(c) Diabetes mellitus (6 marks):}
\begin{itemize}
\item Insulin normally lowers blood glucose by promoting its uptake by cells and conversion to glycogen in the liver (1).
\item Failure of insulin production means blood glucose rises above the \textbf{renal threshold} after meals and stays high (1).
\item The proximal tubule cannot reabsorb all the filtered glucose (its carriers are saturated), so \textbf{glucose remains in the tubule fluid and appears in the urine} (1).
\item The glucose in the tubule \textbf{lowers the water potential} of the filtrate (1), so less water is reabsorbed by osmosis in the tubules and collecting duct (1); the result is a \textbf{large volume of dilute urine} and, through water loss, intense thirst (1).
\end{itemize}
\textbf{(d) Second patient (4 marks):}
\begin{itemize}
\item The cause is a deficiency of \textbf{ADH} (antidiuretic hormone) — or failure of the kidney to respond to it — the condition called diabetes insipidus (1).
\item ADH is normally secreted by the posterior pituitary when the blood becomes too concentrated (1); it increases the \textbf{permeability of the collecting duct (and distal tubule) to water} (1), so more water is reabsorbed and a small volume of concentrated urine is produced. Without ADH the collecting duct stays impermeable, water is not reabsorbed, and large volumes of dilute, glucose-free urine result (1).
\end{itemize}
\textbf{Examiner expectations:} precise terminology (ultrafiltration, basement membrane, active transport, renal threshold, water potential, collecting duct permeability); clear cause-effect chains; no confusion between diabetes mellitus and diabetes insipidus.
\end{corrige}

\begin{corrige}[Exercise 10 — Essay: response to a fall in core temperature]
\textbf{Indicative content and mark scheme (20 marks):}

\textbf{Receptor and coordination (4 marks).} The fall in core temperature is detected by \textbf{thermoreceptors in the hypothalamus} (which monitor the temperature of the blood) and by cold receptors in the skin (1). The \textbf{thermoregulatory centre of the hypothalamus} compares the temperature with the set point (about $37\ ^\circ\mathrm{C}$) (1) and sends nerve impulses to effectors (1); it also triggers hormonal responses (1).

\textbf{Skin structure and diagram (part of the 2 organisation marks + supports effector marks).} A clearly labelled vertical section of the skin should show: epidermis and dermis; \textbf{hair with erector pili muscle}; \textbf{sweat gland} and sweat duct; \textbf{blood capillaries / arterioles} with shunt vessels; sensory nerve endings; subcutaneous fat (insulation).

\begin{popfigure}
\begin{tikzpicture}[scale=1.0]
\fill[popgold!25] (0,2.4) rectangle (8,3.2);
\fill[popblueL] (0,0.8) rectangle (8,2.4);
\fill[gray!15] (0,0) rectangle (8,0.8);
\node[font=\small] at (7.4,2.85) {epidermis};
\node[font=\small] at (7.4,1.7) {dermis};
\node[font=\small] at (7.2,0.4) {fat layer};
\draw[popdark, thick] (2.0,3.2) -- (2.35,4.3);
\draw[popdark, thick] (2.0,3.2) -- (1.7,1.6);
\draw[poporange, thick] (1.75,1.7) -- (2.5,2.9);
\node[font=\small, align=center, text width=2.4cm] at (0.9,3.9) {hair + erector pili muscle};
\draw[->, thin] (1.6,3.7) -- (2.1,3.1);
\draw[poppurple, thick] (4.2,1.1) circle (0.28);
\draw[poppurple, thick] (4.2,1.38) .. controls (4.3,2.2) and (4.1,2.6) .. (4.2,3.2);
\node[font=\small] at (5.3,1.1) {sweat gland};
\draw[popink, very thick] (0.6,1.9) .. controls (2,1.5) and (3,2.2) .. (4.8,1.8) .. controls (6,1.55) and (6.8,2.0) .. (7.6,1.85);
\node[font=\small, popink] at (3.1,2.35) {blood capillaries};
\end{tikzpicture}
\legende{Vertical section of mammalian skin showing structures involved in thermoregulation.}
\end{popfigure}

\textbf{Effector responses (10 marks — about 2 marks each):}
\begin{itemize}
\item \textbf{Blood vessels of the skin — vasoconstriction:} arterioles near the surface constrict and shunt vessels open, diverting blood away from the skin surface, so less heat is lost by radiation, conduction and convection.
\item \textbf{Sweat glands:} sweating is reduced or stops, so less heat is lost by evaporation of sweat.
\item \textbf{Erector pili muscles:} they contract, raising the hairs; the trapped layer of still air is a poor conductor and insulates the body (more effective in furry mammals than in man, where it appears as ``goose pimples'').
\item \textbf{Skeletal muscles:} \textbf{shivering} — rapid involuntary contractions — increases respiration in muscle and therefore heat production.
\item \textbf{Hormones:} increased secretion of \textbf{adrenaline} and \textbf{thyroxine} raises the metabolic rate (rate of respiration) in the liver and muscles, increasing heat production over longer periods.
\end{itemize}

\textbf{Negative feedback loop (4 marks).} Fall in temperature $\rightarrow$ detected by hypothalamus $\rightarrow$ effector responses reduce heat loss and increase heat production $\rightarrow$ core temperature rises back towards $37\ ^\circ\mathrm{C}$ (2). As the set point is reached, the hypothalamus reduces the outgoing impulses and hormonal stimulation, so the responses are switched off (1); the output cancels the input that triggered it — the defining feature of \textbf{negative feedback}, keeping temperature fluctuating narrowly around the set point (1).

\textbf{Examiner expectations:} a logical essay structure (receptor $\rightarrow$ coordination $\rightarrow$ effectors $\rightarrow$ feedback); a labelled diagram; correct pairing of each effector with the mechanism by which it conserves or produces heat; explicit use of the terms stimulus, receptor, effector, response and set point.
\end{corrige}

\begin{corrige}[Exercise 11 — Data analysis: osmoregulation in fishes and excretion in plants]
\textbf{Part A (a) — Freshwater fish (5 marks):}
\begin{itemize}
\item The body fluids of the tilapia have a \textbf{lower water potential} (are more concentrated in solutes) than the surrounding fresh water (1).
\item Water therefore enters by \textbf{osmosis} across the gills and body surface, while salts tend to be lost by diffusion down their concentration gradient (1).
\item Correction by the \textbf{kidneys}: the glomeruli filter large volumes and the tubules reabsorb salts, producing \textbf{large volumes of very dilute urine} to expel the excess water (1).
\item Correction by the \textbf{gills}: special chloride (ion-transporting) cells \textbf{actively take up salts} from the surrounding water against the gradient (1), replacing the salts lost in urine and by diffusion (1).
\end{itemize}
\textbf{Part A (b) — Marine fish (5 marks):}
\begin{itemize}
\item Sea water has a \textbf{lower water potential} than the fish's body fluids, so the fish \textbf{loses water by osmosis} across the gills and gains excess salts (1).
\item To replace the lost water the fish \textbf{drinks sea water} (1).
\item Excess salts from the swallowed sea water are \textbf{actively excreted by the gills} (chloride cells) (1).
\item The \textbf{kidneys produce only small volumes of urine}, conserving water (1); nitrogen is lost mainly as ammonia across the gills, so little water is needed for excretion (1).
\end{itemize}
\textbf{Part B (a) — Discussion (6 marks):}
\begin{itemize}
\item Plants produce gaseous wastes — \textbf{oxygen} from photosynthesis and \textbf{carbon dioxide} from respiration — which diffuse out through stomata and lenticels; \textbf{water vapour} is lost in transpiration (1).
\item They produce \textbf{very little nitrogenous waste} because protein metabolism is limited and nitrogen compounds are recycled within the plant (1).
\item Organic wastes such as \textbf{tannins, resins, latex, alkaloids and calcium oxalate crystals} are not eliminated but \textbf{stored in dead or senescent tissues} — bark, old leaves, fruits — and removed when leaves fall or bark peels (1).
\item Reasons for the absence of specialised excretory organs: wastes are produced slowly; gaseous wastes diffuse away directly; many wastes are stored rather than excreted; plants reuse many substances (e.g. $\ce{CO2}$ and $\ce{O2}$ are mutual raw materials of respiration and photosynthesis) (any 3 valid reasons, 3 marks).
\end{itemize}
\textbf{Part B (b) — Economic examples (4 marks, 2 each):} any two of:
\begin{itemize}
\item \textbf{Rubber latex} from \emph{Hevea brasiliensis} (plantations in the South-West and Littoral regions) — raw material for rubber and tyre manufacture.
\item \textbf{Resins} (e.g. gum arabic/resins from savanna trees) — used in varnishes, adhesives, incense and pharmaceuticals.
\item \textbf{Tannins} from barks (e.g. mangrove, acacia) — used in tanning leather and in dyes.
\item \textbf{Alkaloids} (e.g. quinine-type compounds from \emph{Cinchona}, caffeine from coffee) — used in medicine and beverages.
\end{itemize}
\textbf{Examiner expectations:} correct use of water potential gradients (never ``the fish absorbs water because it is thirsty''); explicit mention of osmosis and active transport; for plants, the key idea that storage in disposable tissues replaces excretory organs; Cameroonian examples must be genuine products with correct uses.
\end{corrige}

\newpage

\coursec{Summary sheet}

\begin{popfigure}
\centering
\begin{tikzpicture}[
  mindmap,
  concept color=popblue,
  level 1/.style={sibling angle=360/7, level distance=5cm},
  level 2/.style={sibling angle=360/7, level distance=3cm},
  every node/.style={concept, execute at begin node=\hskip0pt},
  root concept/.append style={concept color=popdark, font=\large\bfseries},
  concept 1/.append style={concept color=poporange},
  concept 2/.append style={concept color=popgreen},
  concept 3/.append style={concept color=poppurple},
  concept 4/.append style={concept color=poppink},
  concept 5/.append style={concept color=popgold},
  concept 6/.append style={concept color=popblueL},
  concept 7/.append style={concept color=popgreenL},
]
\node[root concept] {Homeostasis}
  child[concept 1] {node[concept] {Feedback Mechanisms}}
  child[concept 2] {node[concept] {Liver Homeostasis}}
  child[concept 3] {node[concept] {Skin Functions}}
  child[concept 4] {node[concept] {Excretory Structures}}
  child[concept 5] {node[concept] {Renal System}}
  child[concept 6] {node[concept] {Osmoregulation}}
  child[concept 7] {node[concept] {Plant Excretion}};
\end{tikzpicture}
\legende{Mind map of Homeostasis chapter (Upper Sixth Biology)}
\end{popfigure}

\begin{bilan}
\textbf{Key Definitions}
\begin{itemize}
\item \cle{Homeostasis}: Maintenance of a stable internal environment despite external fluctuations.
\item \cle{Set point}: The optimal value of a physiological variable around which oscillations occur.
\item \cle{Negative feedback}: A control loop that counteracts a deviation from the set point (e.g., blood glucose regulation).
\item \cle{Positive feedback}: A loop that amplifies a change; rare in homeostasis (e.g., oxytocin in childbirth).
\item \cle{Osmoregulation}: Regulation of water and solute concentrations in body fluids.
\item \cle{Excretion}: Removal of metabolic waste products (urea, uric acid, ammonia, creatinine).
\item \cle{Egestion}: Elimination of undigested food residues; not a metabolic waste process.
\end{itemize}

\textbf{Laws / Principles / Formulas}
\begin{itemize}
\item \cle{Feedback loop sequence}: Stimulus $\to$ Receptor $\to$ Control centre $\to$ Effector $\to$ Response.
\item \cle{Renal clearance}: $C = \dfrac{U \times V}{P}$ where $U$ = urine concentration, $V$ = urine flow rate (mL/min), $P$ = plasma concentration.
\item \cle{Glomerular filtration rate (GFR)}: $GFR = \dfrac{U_{inulin} \times V}{P_{inulin}}$ (inulin clearance).
\item \cle{Counter-current multiplier}: Loop of Henle creates a medullary osmotic gradient essential for water reabsorption.
\item \cle{Hormonal control}: ADH increases collecting duct water permeability; aldosterone stimulates Na$^+$ reabsorption in distal tubule; insulin/glucagon regulate blood glucose.
\end{itemize}

\textbf{Methods / Step-by-step Procedures}
\begin{itemize}
\item \textbf{Analysing a feedback loop}: Identify stimulus, receptor, control centre, effector, response; classify as negative or positive.
\item \textbf{Describing urine formation}: 1) Glomerular filtration (pressure-driven), 2) Proximal tubule reabsorption (glucose, amino acids, ions, water), 3) Loop of Henle (counter-current multiplication), 4) Distal tubule and collecting duct (fine-tuning under ADH/aldosterone), 5) Excretion.
\item \textbf{Comparing excretory products}: Ammonia (high toxicity, high water loss — aquatic animals), Urea (moderate toxicity, moderate water loss — mammals), Uric acid (low toxicity, low water loss — birds, reptiles, insects).
\item \textbf{Drawing a labelled nephron}: Include glomerulus, Bowman's capsule, proximal convoluted tubule, loop of Henle (descending/ascending limbs), distal convoluted tubule, collecting duct; indicate blood vessels (afferent/efferent arterioles, peritubular capillaries, vasa recta).
\end{itemize}

\textbf{Common Mistakes (Traps)}
\begin{itemize}
\item Confusing \cle{excretion} (metabolic wastes) with \cle{egestion} (undigested food).
\item Overlooking the liver's role in glucose homeostasis (glycogen storage, gluconeogenesis) and urea synthesis (ornithine cycle).
\item Misidentifying nephron segments: proximal tubule = bulk reabsorption; loop of Henle = gradient formation; distal tubule/collecting duct = hormonal regulation.
\item Forgetting that ADH acts on collecting duct aquaporins, while aldosterone acts on distal tubule Na$^+$/K$^+$ pumps.
\item Stating that plants excrete urea; they mainly store wastes in vacuoles, shed leaves, or release into soil (e.g., tannins, alkaloids).
\item Mixing up diabetes insipidus (ADH deficiency) with diabetes mellitus (insulin deficiency).
\end{itemize}

\textbf{GCE Advanced Level Tips}
\begin{itemize}
\item Always sketch a nephron diagram in essay questions; label all segments and associated blood vessels.
\item Structure long answers: Definition $\to$ Mechanism $\to$ Example $\to$ Physiological significance.
\item In calculations, show units clearly; use clearance formula with $V$ in mL/min and concentrations in same units.
\item Know the hormonal cascade: hypothalamus $\to$ posterior pituitary (ADH) $\to$ kidney; adrenal cortex (aldosterone) $\to$ kidney; pancreas (insulin/glucagon) $\to$ liver/muscle/adipose.
\item Relate to Cameroon context: Sickle cell trait confers malaria resistance (homeostasis of haemoglobin); kidney disease prevalence linked to hypertension and infections.
\item Practice past GCE questions on feedback loops, kidney function, and osmoregulation adaptations (e.g., desert rodents, freshwater fish).
\end{itemize}
\end{bilan}

\findecours

\end{document}
